\documentclass[11pt]{article}

\usepackage[a4paper,margin=1in]{geometry}
\usepackage{amsmath,amssymb,amsthm,mathtools}
\usepackage[hidelinks]{hyperref}
\hypersetup{
  pdfauthor={Bangyen Pham},
  pdftitle={Every Row Is a Top Row: Coefficient Order Statistics of Multiples of a Power},
  pdfsubject={Exact coefficient order statistics of polynomial multiples of (x-r)^L},
  pdfkeywords={polynomial multiples, coefficient mass, reduced height, Chebyshev systems, exponential sums},
  pdflang={en-US}
}

\newtheorem{theorem}{Theorem}[section]
\newtheorem{lemma}[theorem]{Lemma}
\newtheorem{corollary}[theorem]{Corollary}
\newtheorem{proposition}[theorem]{Proposition}
\newtheorem*{mainthm}{Main Theorem}
\newcommand{\R}{\mathbb R}
\newcommand{\N}{\mathbb N}
\DeclareMathOperator{\Tail}{Tail}
\DeclareMathOperator{\sgn}{sgn}

\title{Every Row Is a Top Row: Coefficient Order Statistics of Multiples of a Power}
\author{\shortstack{Bangyen Pham\\
\small Independent Researcher\\
\small\href{mailto:bangyenp@gmail.com}{\texttt{bangyenp@gmail.com}}\\
\small\href{https://orcid.org/0009-0006-9928-2088}{ORCID: 0009-0006-9928-2088}}}
\date{September 2026}

\begin{document}
\maketitle

\begin{abstract}
For a real polynomial $F=\sum_jf_jx^j$ of degree $D$ let $b_k(F)$ be the
$k$-th largest of the nonleading coefficient magnitudes $|f_j|$, $j<D$.  Fix
a real $r>1$ and let $V_r(L,k)$ be the infimum of $b_k(F)$ over monic
multiples $F$ of $(x-r)^L$.  We show that for every $r\ge2$ every row is a
top row: $V_r(L,k)=V_r(L-k+1,1)$ for $1\le k\le L$.  Exempting the $k-1$
largest coefficients thus costs exactly $k-1$ of the $L$ roots.  At $r=2$
this sharpens the bound $b_k\ge L-k+1$.  The proof reduces each row to an
$\ell^1$ problem for polynomials with prescribed integer zeros, covers all
exempted positions beyond the first gap of a minimizer by one crossing step,
and bounds the cost of that step by first-order optimality.  For $1<r<2$ the
rows are not top rows, and we conjecture that they are given by
prefix-constrained problems.  We prove this whenever the top row is at most
one, and on each diagonal $L-k=\mathrm{const}$ for all but finitely many
rows.  Along the way we treat several roots, finite certificates for the
rows, and the hole products of top-row optima.
\end{abstract}

\noindent\textbf{Keywords.} Polynomial multiples; coefficient order
statistics; reduced height; linear programming duality.

\noindent\textbf{2020 Mathematics Subject Classification.} Primary 11C08; Secondary 26C05, 90C05.

\section{Introduction}

Let $F=\sum_jf_jx^j$ be a nonzero real polynomial of degree $D$, and write
$b_k(F)$ for the $k$-th largest of its nonleading coefficient magnitudes
$|f_j|$, $j<D$.  The companion paper~\cite{coefficient-mass} proves, for real
roots $2\le r_1\le\cdots\le r_L$, that every nonzero multiple $F$ of
$\prod_i(x-r_i)$ has
\[
  b_k(F)\ \ge\ |f_D|\prod_{i=k}^L(r_i-1)\qquad(1\le k\le L)
\]
\cite[Theorem~3.2]{coefficient-mass}, and that at the root $2$ every monic
multiple of $(x-2)^L$ has $b_k(F)\ge L-k+1$
\cite[Theorem~4.16]{coefficient-mass}.  We call the bound at a given $k$
the row $k$.  This paper determines the rows when all roots are equal.

\paragraph{The rows at a single root.}
Fix a real $r>1$ and put
\[
  V_r(L,k)=\inf\{b_k(F):\ F\text{ a monic real multiple of }(x-r)^L\},
  \qquad \beta_r(n)=V_r(n,1),
\]
so $\beta_r(n)$ is the top row of $(x-r)^n$.  Exempting the positions of
the $k-1$ largest coefficients can never help by more than it costs in
roots, and $V_r(L,k)\le\beta_r(L-k+1)$
(Proposition~\ref{prop:rowvalue}).  Our main result is that equality holds
at every root $r\ge2$.

\begin{mainthm}
For every real $r\ge2$ and all $1\le k\le L$,
$V_r(L,k)=\beta_r(L-k+1)$ (Theorem~\ref{thm:allrowstwo}).
\end{mainthm}

At $r=2$ this sharpens $b_k\ge L-k+1$ to $b_k\ge\beta_2(L-k+1)$, and
$\beta_2(3)\ge\frac{40}{11}$.  Earlier partial results proved the identity
at $r\ge\frac92$, $r\ge4$ and $r\ge\frac{25}7$ by bounding the hole product
of a top-row optimum.  That route cannot reach $r=2$, and the appendix
keeps it for the information it gives about top-row optima.

For $1<r<2$ the rows are not top rows.  The last row is
$V_r(L,L)=(r-1)^L<\beta_r(1)$ (Corollary~\ref{cor:lastrowr}).  The natural
comparison values are the prefix values $\nu_i(n)$: the least weighted
$\ell^1$ norm of a polynomial that vanishes on $[1,i-1]$.  They satisfy
\[
  \frac1{\nu_k(n)+E_k(r)}\ \le\ V_r(L,k)\ \le\ \min_{1\le i\le k}\frac1{\nu_i(n)},
  \qquad n=L-k+1,
\]
with $E_k(r)/\nu_k(n)\to0$ as $k\to\infty$ (Theorem~\ref{thm:prefixrows}).
We conjecture that the upper bound is always attained.  We prove this for
every $k$ whenever $\beta_r(n)\le1$ (Theorem~\ref{thm:onepolyrows}).  We
also prove it whenever $k\ge\max(K,\alpha(n-1))$ for constants depending on
$r$ (Theorem~\ref{thm:longdiag}, Proposition~\ref{prop:longdiagexplicit}),
and for complete diagonals at $r=\frac54$, $\frac32$ and $\frac74$
(Corollary~\ref{cor:completediag}).

\paragraph{Method.}
Divisibility by $(x-r)^L$ is a set of linear conditions on the
coefficients, and duality for linear programs turns each row into a problem
about polynomials
$q$ of degree less than $L$ with $q(0)=1$: minimize
\[
  \Phi_r(q)=\sum_{s\ge1}|q(s)|\,r^{-s}
\]
subject to $q$ vanishing on the set $S$ of backward distances of the
exempted positions (Proposition~\ref{prop:rowvalue}).  Minimizers can be taken with integer
zeros.  Two facts then settle $r\ge2$.
\begin{itemize}
\item A convexity argument, crossing, shows that one insertion at the
first gap of a minimizer handles every later exempted position at once
(Theorem~\ref{thm:crossing}).
\item First-order optimality of the minimizer in a single direction bounds
the cost of that insertion by the factor $\max\{1,1/(r-1)\}$
(Lemma~\ref{lem:optins}).
\end{itemize}
Below $2$ this factor exceeds $1$, and the proofs use instead one
polynomial for all exempted sets (Theorem~\ref{thm:onepoly}) and the
supermultiplicativity of the prefix values (Lemma~\ref{lem:prefixshift}).

\paragraph{Organization.}
Section~\ref{sec:rowvalues} sets up the row values, extends the tail bound
of~\cite{coefficient-mass} to arbitrary nodes, and proves the last row and
first bounds at every root.  Section~\ref{sec:prefix} brackets the rows
below $2$ by prefix values and treats several roots.
Sections~\ref{sec:holeshifts}--\ref{sec:crossing} develop the tools: hole
shifts and insertions, finite certificates, monotone prefix values and
crossing.  Section~\ref{sec:allrowstwo} proves that every row is a top row
at $r\ge2$, and Section~\ref{sec:belowtwo} treats $1<r<2$.
Appendix~\ref{sec:holeproducts} collects bounds on the hole products of
top-row optima.

\paragraph{Conventions.}
All logarithms are natural and $\N_{>0}=\{1,2,\ldots\}$.  For a finite set
$Z\subset\N_{>0}$ we write $r_Z(s)=\prod_{z\in Z}(1-s/z)$, so $r_Z(0)=1$
and $\deg r_Z=|Z|$.  Results of the companion paper are cited by their
numbers there.

\section{Row values and certificates}\label{sec:rowvalues}

Let $r>1$ be real.  Throughout the paper
\[
  a=\frac1{r-1},\qquad
  \Phi_r(q)=\sum_{s\ge1}|q(s)|\,r^{-s},\qquad
  \Phi_r(Z)=\Phi_r(r_Z),\qquad
  \psi_r(y)=\frac{1+a^y}{ry}\quad(y\ge1),
\]
for real polynomials $q$ and finite $Z\subset\N_{>0}$, so $\Phi_2$ is the sum $\Phi$ of \cite[Section~4.4]{coefficient-mass} and
$\Phi_r(\varnothing)=a$.  Write $[u,v]$ for the integers in that range.  For
$1\le k\le L$ put $n=L-k+1$ and
\[
  V_r(L,k)=\inf\{b_k(F): F\text{ a monic real multiple of }(x-r)^L\},
\]
the exact value of the row.  For finite $S\subset\N_{>0}$ with $|S|<L$ let
$\mu_r(S,L)$ be the infimum of $\Phi_r(q)$ over real $q$ with $\deg q\le L-1$,
$q(0)=1$ and $q|_S=0$, and put $\beta_r(n)=1/\mu_r(\varnothing,n)$.  At
$r=2$, \cite[Theorem~4.16]{coefficient-mass} says $V_2(L,k)\ge n$.

\begin{proposition}[The row values]\label{prop:rowvalue}
Let $r>1$ and $1\le k\le L$.
\begin{enumerate}
\item[(a)] If $F$ is a monic multiple of $(x-r)^L$, $S$ is the set of
backward distances of $k-1$ positions carrying the $k-1$ largest nonleading
magnitudes, and $q$ is real with $\deg q\le L-1$, $q(0)=1$ and $q|_S=0$, then
$b_k(F)\,\Phi_r(q)\ge1$.  In particular $b_k(F)\ge1/\Phi_r(Z)$ for every
finite $Z\supseteq S$ with $|Z|\le L-1$.
\item[(b)] $V_r(L,k)=\inf_{|S|=k-1}1/\mu_r(S,L)$, the infimum over
$S\subset\N_{>0}$.
\item[(c)] $V_r(L,k)\le\beta_r(n)$, and $V_r(n,1)=\beta_r(n)$.
\end{enumerate}
\end{proposition}

\begin{proof}
(a) This is the proof of \cite[Lemma~4.8]{coefficient-mass} with $2$ replaced by $r$:
$(q(D-\theta)F)(r)=0$ with $\theta=x\,d/dx$ gives
$\sum_{j\le D}f_jq(D-j)r^{j-D}=0$, so
$1\le\sum_{s\ge1,\,s\notin S}|f_{D-s}||q(s)|r^{-s}\le b_k(F)\Phi_r(q)$.

(b) By (a), $b_k(F)\ge1/\mu_r(S_F,L)$ for the set $S_F$ of (a), which
gives ``$\ge$''.  For ``$\le$'' fix $S$ with $|S|=k-1$ and $D\ge\max(L,\max S)$.
Writing $F=\sum_{s=0}^Dg_sx^{D-s}$ with $g_0=1$, divisibility by
$(x-r)^L$ is the set of $L$ linear conditions
$\sum_{s=0}^Dg_sq(s)r^{-s}=0$ for $\deg q\le L-1$ (as in (a), these are the
conditions $(\theta^iF)(r)=0$, $i<L$).  Minimizing
$\max_{s\in[1,D]\setminus S}|g_s|$ over this affine set is a finite linear
program, feasible because $D\ge L$, and by linear-programming duality its value is
$1/\mu_{r,D}(S,L)$, where $\mu_{r,D}$ is defined like $\mu_r$ with the sum
$\Phi_r(q)$ truncated to $s\le D$ (the dual variables are the values
$q(s)r^{-s}$ of a polynomial $q$ of degree at most $L-1$ vanishing on $S$;
for weak duality, $|q(0)|=|\sum_{s\ge1,\,s\notin S}g_sq(s)r^{-s}|$).  Since
the $k-1$ positions at distances $S$ can carry at most the $k-1$ largest
magnitudes, $b_k(F)\le\max_{s\notin S}|g_s|$, so
$V_r(L,k)\le1/\mu_{r,D}(S,L)$ for every $D$.  Finally
$\mu_{r,D}(S,L)\to\mu_r(S,L)$: the truncated sums increase with $D$ and
are at most $\Phi_r$; on the affine space $\{\deg q\le L-1,\ q(0)=1\}$ the
values $q(1),\ldots,q(L-1)$ are coordinates, so minimizers $q_D$ of the
truncated problems ($D\ge L$) have bounded coefficients, and a limit point
$q_*$ satisfies $q_*(0)=1$, $q_*|_S=0$ and, for each fixed $M$,
$\sum_{s\le M}|q_*(s)|r^{-s}\le\lim_D\mu_{r,D}(S,L)$; letting $M\to\infty$
gives $\mu_r(S,L)\le\Phi_r(q_*)\le\lim_D\mu_{r,D}(S,L)$.

(c) Let $S_N=[N+1,N+k-1]$.  A polynomial $q$ of degree at most $L-1$ with
$q(0)=1$ and $q|_{S_N}=0$ factors as $q=q_0\pi_N$ with
$\pi_N(s)=\prod_{\sigma\in S_N}(1-s/\sigma)$, $\deg q_0\le n-1$ and
$q_0(0)=1$.  For $1\le s\le M<N$ we have $\pi_N(s)\ge(1-M/N)^{k-1}$, so
$\Phi_r(q)\ge(1-M/N)^{k-1}\mu_{r,M}(\varnothing,n)$.  Hence
$\liminf_N\mu_r(S_N,L)\ge\mu_{r,M}(\varnothing,n)$ for every $M$, and by
the convergence in (b) it is at least $\mu_r(\varnothing,n)$.  By (b),
$V_r(L,k)\le\inf_N1/\mu_r(S_N,L)\le\beta_r(n)$.  For $k=1$, (b) with
$S=\varnothing$ is $V_r(n,1)=\beta_r(n)$.
\end{proof}

So every row is at most the top row of a multiple of $(x-r)^n$, the value
approached when the $k-1$ exempted positions lie far below the leading one.
The lower bounds come from certificates, and the first tool is the
displaced-zero tail bound without the cutoff $\rho_i\ge2$.

\begin{theorem}[Tail bound at arbitrary nodes]\label{thm:tailgen}
Let $0<y_1<\cdots<y_c<1$, let $Z\subset\N_{>0}$ with $|Z|=c-1$ have leading
run $\ell$, let $w$ be the exponential sum on these nodes with $w_0=1$ and
$w|_Z=0$, and put $\varphi_i=y_i/(1-y_i)$.  Then
\[
  \Tail(w)\le\prod_{i=1}^{\ell+1}\varphi_i\prod_{i=\ell+2}^c\max(1,\varphi_i).
\]
Moreover, if $Z\ne[1,c-1]$ and $F$ is the sum on $y_1,\ldots,y_{c-1}$ with
$F_0=1$ and $F|_{Z\setminus\{\max Z\}}=0$, then
$\Tail(w)\le\max(1,\varphi_c)\Tail(F)$.
\end{theorem}

\begin{proof}
The sum $w$ exists and is unique for any distinct positive nodes, by the
Vandermonde argument at the start of \cite[Section~2]{coefficient-mass}.  Induct on $c-1-\ell$ as in
the proof of \cite[Theorem~2.2]{coefficient-mass}.  The base case $Z=[1,c-1]$ is
\cite[Lemma~2.1]{coefficient-mass}, which gives $\prod_{i=1}^c\varphi_i$ and needs
no cutoff.  In the displaced step, the sign comparison, the splitting
identity and the first correction-ratio inequality in the proof of
\cite[Theorem~2.2]{coefficient-mass},
$\Gamma_G(z)\le\Gamma_F(z)/(1-y_c)$, use only the zero bound,
\cite[Lemma~2.3]{coefficient-mass} (which needs no upper bound on the nodes) and
$y_c<1$.  If $y_c\le\frac12$, the proof of \cite[Theorem~2.2]{coefficient-mass} gives
$\Tail(w)<\Tail(F)$.  If $y_c>\frac12$, then $\Sigma_G\ge0$ and
the splitting identity give
\[
  \Tail(w)\le\Tail(F)+\mu\Gamma_F(z)\Bigl(\frac1{1-y_c}-2\Bigr),
\]
where the bracket is positive and
$\mu\Gamma_F(z)=\sum_{d\ge z}|F_d|\le\Tail(F)$ (if $F_z=0$ then $w=F$), so
$\Tail(w)\le(\frac1{1-y_c}-1)\Tail(F)=\varphi_c\Tail(F)$.  In both cases
$\Tail(w)\le\max(1,\varphi_c)\Tail(F)$.  The set $Z\setminus\{\max Z\}$ has
the same leading run $\ell$ and $\ell+1\le c-1$, so the induction hypothesis
for $F$ on $y_1,\ldots,y_{c-1}$ gives the bound.
\end{proof}

For $y_c\le\frac12$ every $\max(1,\varphi_i)$ is $1$ and this is
\cite[Theorem~2.2]{coefficient-mass}.  The order statistics extend to all roots above $1$.

\begin{corollary}[Order statistics for roots above one]\label{cor:ordergen}
Let $1<r_1\le\cdots\le r_L$ and let $F$ be a nonzero real multiple of
$\prod_{i=1}^L(x-r_i)$.  Then for every $0\le u\le L-1$,
\[
  b_{u+1}(F)\ge|f_D|\prod_{i=u+1}^L(r_i-1)\prod_{i=1}^u\min(1,r_i-1).
\]
\end{corollary}

\begin{proof}
Repeat the proofs of \cite[Lemma~3.1]{coefficient-mass} and \cite[Theorem~3.2]{coefficient-mass}.  In
\cite[Lemma~3.1]{coefficient-mass} the reciprocal nodes $y_1<\cdots<y_L$ are the
$1/r_i$ in decreasing order of $r_i$ and the leading run of $Z$ is
$\ell\ge L-u-1$; Theorem~\ref{thm:tailgen}, together with
$\varphi_i\le\max(1,\varphi_i)$ for $L-u<i\le\ell+1$, bounds the tail by
$\prod_{i\le L-u}\varphi_i\prod_{i>L-u}\max(1,\varphi_i)$, which is the
reciprocal of the right side (with $|f_D|=1$) because $\varphi_i=1/(\rho_i-1)$.
The perturbation of \cite[Theorem~3.2]{coefficient-mass} applies unchanged, the bound
being continuous in the roots.
\end{proof}

For $r_1\ge2$ this is \cite[Theorem~3.2]{coefficient-mass}.  At a single root
$r\in(1,2]$ it gives $b_k\ge(r-1)^L$ in every row, which is sharp in the
last row (Corollary~\ref{cor:lastrowr}).

\begin{lemma}[Deleting the largest zero at the root $r$]\label{lem:confdelr}
Let $r>1$.  Then $\Phi_r([1,\ell])=a^{\ell+1}$ for $\ell\ge0$, and for every
nonempty finite $Z$,
\[
  \Phi_r(Z)\le\max(1,a)\,\Phi_r\bigl(Z\setminus\{\max Z\}\bigr).
\]
Consequently $\Phi_r(Z)\le\Phi_r(Z\cap[1,t])$ for all $t$ when $r\ge2$, and
$\Phi_r(Z)\le a^{|Z|+1}$ when $1<r\le2$.
\end{lemma}

\begin{proof}
$|r_{[1,\ell]}(s)|=\binom{s-1}\ell$ for $s>\ell$ and
$\sum_{s>\ell}\binom{s-1}\ell r^{-s}=(r^{-1}/(1-r^{-1}))^{\ell+1}=a^{\ell+1}$;
this also settles $Z=[1,\ell]$, the ratio being $a$.  For
$Z\ne[1,|Z|]$ run the limit in the proof of \cite[Lemma~4.9]{coefficient-mass} with
the nodes $y_i=\frac1r-(c-i)\varepsilon$ and the comparison
$\Tail(w^\varepsilon)\le\max(1,\varphi_c)\Tail(F^\varepsilon)$ of
Theorem~\ref{thm:tailgen}; the domination
$0\le h_{d-i+1}(y_1,\ldots,y_i)\le\binom d{i-1}r^{-(d-i+1)}$ and the
invertibility of the limit matrix are as there, and $\varphi_c\to a$.  The
consequences follow by deleting the largest element repeatedly: for $r\ge2$
each step does not increase $\Phi_r$, and for $r\le2$ each of the $|Z|$
steps from $Z$ to $\varnothing$ costs at most the factor $a$, with
$\Phi_r(\varnothing)=a$.
\end{proof}

\begin{corollary}[The last row]\label{cor:lastrowr}
$V_r(L,L)=r-1$ for $r\ge2$, and $V_r(L,L)=(r-1)^L$ for $1<r\le2$.
\end{corollary}

\begin{proof}
Lower bounds: Proposition~\ref{prop:rowvalue}(a) with $Z=S$, $|S|=L-1$, and
Lemma~\ref{lem:confdelr}, which gives $\Phi_r(S)\le a$ for $r\ge2$ and
$\Phi_r(S)\le a^L$ for $r\le2$.  Upper bounds: for $r\ge2$,
Proposition~\ref{prop:rowvalue}(c) with $n=1$ and
$\beta_r(1)=1/\Phi_r(1)=r-1$; for $r\le2$, Proposition~\ref{prop:rowvalue}(b)
with $S=[1,L-1]$, where $q=r_S$ is forced and $\mu_r(S,L)=a^L$.
\end{proof}

At $r=2$ both values are $1$.  For the other rows at $r\ge2$ we need the
hole integral at the root $r$.  For finite nonempty $C\subset\N_{>0}$ with
$m=\max C$ and $w_c$ as in \cite[Lemma~4.10]{coefficient-mass}, put
\[
  T_C(u)=\sum_{c\in C}w_cr^{-c}u^{m-c}(1+u)^{c-1},\qquad
  H_C(t)=\sum_{c\in C}|w_c|\,r^{-c}t^{m-c}(1-t)^{c-1}.
\]

\begin{lemma}[Hole integral at the root $r$]\label{lem:holeintr}
Let $r>1$, $k\ge0$ and $P=[1,m+k]\setminus C$.  Then
\[
  \Phi_r(P)=\int_0^1t^k\bigl(a^{k+1}T_C(at)+H_C(t)\bigr)\,dt,
\]
where the second summand integrates to $\sum_{c\in C}|r_P(c)|r^{-c}$ and
the first to $\sum_{s>m+k}|r_P(s)|r^{-s}$.  In particular, for $c\ge1$ and
$t_0\ge0$,
\[
  S_r(c,t_0):=\Phi_r\bigl([1,c+t_0]\setminus\{c\}\bigr)
  =\frac{r^{-c}}{\binom{c+t_0}c}
   +c\,r^{-c}\sum_{i=0}^{c-1}\binom{c-1}i\frac{a^{t_0+i+1}}{t_0+i+1},
\]
with $S_r(1,t_0)=\psi_r(t_0+1)$ and $S_r(c,0)=a^c$.
\end{lemma}

\begin{proof}
Put $x=m+k$.  As in the proof of \cite[Lemma~4.10]{coefficient-mass}, with
$\sum_{s>x}\binom{s-1}xy^s=(y/(1-y))^{x+1}$ at $y=u/r$ and then
$v=u/(r-u)$,
\[
  \sum_{s>x}\binom{s-1}x\frac{r^{-s}}{s-c}
  =\int_0^1\frac{u^{x-c}\,du}{(r-u)^{x+1}}
  =r^{-c}\int_0^av^{x-c}(1+v)^{c-1}\,dv
  =r^{-c}a^{x-c+1}\int_0^1t^{x-c}(1+at)^{c-1}\,dt,
\]
and $a^{x-c+1}t^{x-c}=a^{k+1}t^k(at)^{m-c}$.  The values at the holes are
$|r_P(c)|=|w_c|\int_0^1t^{x-c}(1-t)^{c-1}dt$ exactly as there.  For
$C=\{c\}$, $w_c=c$, the hole value is $r^{-c}/\binom xc$, and expanding
$(1+v)^{c-1}$ gives the sum.  For $c=1$ it is $(1+a^{t_0+1})/(r(t_0+1))$;
for $t_0=0$ the set is $[1,c-1]$.
\end{proof}

At $r=2$, $a=1$ and $T_C+H_C=M_C$, which is \cite[Lemma~4.10]{coefficient-mass}.  For
$r>2$ the tail carries the extra factor $a^{k+1}$; the target $\psi_r$
splits accordingly as $\frac1{r(k+1)}+\frac{a^{k+1}}{r(k+1)}$.

\begin{lemma}[A sufficient condition]\label{lem:psiprin}
Let $r\ge2$, $k\ge0$ and $G_C(t)=r\bigl(H_C(t)+T_C(at)\bigr)$.  If
\[
  \int_0^1t^kG_C(t)\,dt\le\frac2{k+1}
  \qquad\text{and}\qquad
  \int_0^1t^k\,rH_C(t)\,dt\le\frac1{k+1},
\]
then $\Phi_r([1,m+k]\setminus C)\le\psi_r(k+1)$.
\end{lemma}

\begin{proof}
By Lemma~\ref{lem:holeintr},
$r\Phi_r=\int_0^1t^k\bigl(a^{k+1}G_C+(1-a^{k+1})rH_C\bigr)dt$, and
$0<a\le1$, so $r\Phi_r\le\frac{2a^{k+1}+1-a^{k+1}}{k+1}=r\psi_r(k+1)$.
\end{proof}

\begin{proposition}[Rows with the target $\psi_r$]\label{prop:psirows}
Let $r\ge2$, $n\ge1$ and $A\subset\N_{>0}$ finite.  Some finite $Z\supseteq A$
has $|Z|\le|A|+n-1$ and $\Phi_r(Z)\le\psi_r(n)$.
\end{proposition}

\begin{proof}
For $n=1$ take $Z=A$: $\Phi_r(A)\le\Phi_r(\varnothing)=a=\psi_r(1)$ by
Lemma~\ref{lem:confdelr}.  Let $n\ge2$, let $j$ be the leading run of $A$,
and let $c_1=j+1<c_2<\cdots$ be the positive integers outside $A$.  Keep the
holes $K=\{1\}$, $\{2\}$, $\{3,c_2\}$, $\{c_1,\ldots,c_{j+1}\}$ for $j=0$,
$1$, $2$, $j\ge3$, fill the next $n-1$ holes, and let $x$ be the last filled
hole, $P=[1,x]\setminus K$ and $Z=P\cup A$, exactly as in
\cite[Proposition~4.12]{coefficient-mass} and \cite[Proposition~4.15]{coefficient-mass}.  Then $Z\supseteq A$,
$|Z|\le|A|+n-1$, and $\Phi_r(Z)\le\Phi_r(P)$ by Lemma~\ref{lem:confdelr},
the elements of $Z$ above $x$ being deleted one at a time.  Put
$m=\max K$ and $k=x-m\ge n-1$.  Since $\psi_r$ is decreasing ($a\le1$), it
suffices to show $\Phi_r(P)\le\psi_r(k+1)$, or, for $j\ge3$,
$\Phi_r(P)\le(1+a^n)/(r(k+1))$.

\emph{$j=0$.}  $P=[2,x]$ and $\Phi_r(P)=\psi_r(x)=\psi_r(k+1)$ by
Lemma~\ref{lem:holeintr}.

\emph{$j=1$.}  $C=\{2\}$: $rH_C=\frac2r(1-t)\le1$ and
$G_C=\frac2r\bigl(2-(1-a)t\bigr)\le\frac4r\le2$, so
Lemma~\ref{lem:psiprin} applies.

\emph{$j=2$, $b=c_2\ge5$.}  $C=\{3,b\}$, $w_b=-w_3=\kappa:=3b/(b-3)$.  Writing
$B_i=\binom{b-1}it^i(1-t)^{b-1-i}$ and $1+at=(1-t)+rat$, a direct expansion gives
$G_C=\kappa\sum_{i=0}^{b-1}\gamma_iB_i$ with
\[
  \gamma_0=2r^{1-b},\quad
  \gamma_i=r^{1-b}(ra)^i\ (1\le i\le b-4),\quad
  \gamma_{b-3}=\frac{a^{b-3}(1-\epsilon)+\epsilon}{r^2},\quad
  \gamma_{b-2}=\frac{a^{b-2}}r\cdot\frac{b-3}{b-1},\quad
  \gamma_{b-1}=0,
\]
where $\epsilon=1/\binom{b-1}2$.  With $r\ge2$, $a\le1$ and $ra\le2$:
$\gamma_0\le2^{2-b}\le\frac18$ and $\gamma_i\le2^{1-b+i}\le\frac18$ for $1\le i\le b-4$, so $\kappa\gamma_i\le\frac{3b}{8(b-3)}\le1$ for $i\le b-4$;
$\kappa\gamma_{b-3}\le\kappa/4\le\frac{15}8$ and
$\kappa\gamma_{b-2}\le\frac{3b}{2(b-1)}\le2$.  So $G_C\le2$.  Next
$rH_C=\kappa\bigl(r^{-2}t^{b-3}(1-t)^2+r^{1-b}(1-t)^{b-1}\bigr)$, and
$\int_0^1t^kt^{b-3}(1-t)^2\,dt\le\frac2{(k+1)b(b-1)}$,
$\int_0^1t^k(1-t)^{b-1}dt\le\frac1{k+1}$, so
$(k+1)\int t^k rH_C\le\frac3{2(b-3)(b-1)}+\frac{3b\,2^{1-b}}{b-3}\le\frac3{16}+\frac{15}{32}<1$.

\emph{$j=2$, $b=4$.}  Here $T_C(at)=12r^{-4}(1+at)^2(1-t)$ because
$1+at-rat=1-t$, so
\[
  G_C=12(1-t)\Bigl(\frac{t(1-t)}{r^2}+\frac{(1-t)^2}{r^3}
       +\frac1r\Bigl(\frac{1+at}r\Bigr)^2\Bigr)\le3(1-t^2),
\]
each term being largest at $r=2$ ($(1+at)/r$ decreases in $r$), and
$\int_0^1t^k3(1-t^2)\,dt=\frac6{(k+1)(k+3)}\le\frac2{k+1}$.  Also
$rH_C\le3t(1-t)^2+\frac32(1-t)^3$, whose moment is
$\frac{6k+15}{(k+1)(k+2)(k+3)(k+4)}\le\frac1{k+1}$.

\emph{$j\ge3$.}  Here $k\ge1$.  Split
$\Phi_r(P)=H_r+T_r$ into the values at the kept holes and the tail, as in
the paragraph before \cite[Lemma~4.13]{coefficient-mass}.  Every kept hole is at least
$j+1$ and $2/r\le1$, so with $H(C,k)$ of that paragraph and
\cite[Lemma~4.13]{coefficient-mass},
\[
  H_r\le(2/r)^{j+1}H(C,k)\le\frac{(2/r)^{j+1}}{2(k+1)}\le\frac1{r(k+1)} .
\]
For the tail, the first two steps of the proof of \cite[Lemma~4.14]{coefficient-mass}
do not involve the root, and summing
$\binom{s-1}xr^{-s}(1-t)^{s-m-1}$ over $s>x$ gives
$(1-t)^k(r-1+t)^{-x-1}=a^{x+1}(1-t)^k(1+at)^{-x-1}$, so
\[
  T_r\le(j+1)\binom mj a^{m+k+1}\int_0^1t^j
     \Bigl(\frac{1-t}{1+at}\Bigr)^k(1+at)^{-m-1}\,dt .
\]
The function $\log(1+at)-\log(1-t)$ vanishes at $0$, has derivative $1+a$
there and is convex (as $a/(1+at)\le1\le1/(1-t)$), so
$(1-t)/(1+at)\le e^{-(1+a)t}$; and $(k+1)e^{-(1+a)kt}\le2/(e(1+a)t)$.
Extending the integral to $(0,\infty)$ and using
$\int_0^\infty t^{j-1}(1+at)^{-m-1}dt=a^{-j}\mathrm B(j,m-j+1)$,
$\binom mj\mathrm B(j,m-j+1)=1/j$ and $1+a=ra$,
\[
  (k+1)T_r\le\frac{2(j+1)}{e\,j\,r}\,a^{m+k-j}
  \le\frac{2(j+1)}{e\,j}\,a^{j}\cdot\frac{a^n}r\le\frac{a^n}r,
\]
because $m\ge2j+1$ and $k+1\ge n$ give $m+k-j\ge n+j$, and
$2(j+1)/(ej)\le8/(3e)<1$ for $j\ge3$.  Hence
$\Phi_r(P)\le(1+a^n)/(r(k+1))\le\psi_r(n)$.
\end{proof}

At $r=2$ the proposition is $T(n)$.  For larger $r$ its target $\psi_r(n)$
is weaker than the row $a^n$ of \cite[Theorem~3.2]{coefficient-mass} once $(r-1)^n+1>nr$; the
remedy is to keep one hole in the middle.

\begin{lemma}[Moving the hole]\label{lem:shift}
For $r\ge2$ and $t_0\ge0$, $S_r(2,t_0)\le S_r(1,t_0)$.  For $r\ge\frac52$,
$h\ge1$ and $t_0\ge0$, $S_r(h+1,t_0)\le S_r(h,t_0)$.
\end{lemma}

\begin{proof}
By Lemma~\ref{lem:holeintr},
$S_r(2,t_0)=\frac2{r^2}\bigl(\frac{a^{t_0+1}}{t_0+1}+\frac{a^{t_0+2}}{t_0+2}
+\frac1{(t_0+1)(t_0+2)}\bigr)$, so
$r(t_0+1)S_r(2,t_0)=\frac2r\bigl(a^{t_0+1}+\frac{(t_0+1)a^{t_0+2}+1}{t_0+2}\bigr)
\le\frac2r(1+a^{t_0+1})\le1+a^{t_0+1}=r(t_0+1)S_r(1,t_0)$.
For $h\ge2$ write $S_r(h,t_0)=\eta(h)+\tau(h)$ with
$\eta(h)=r^{-h}/\binom{h+t_0}h$ and
$\tau(h)=hr^{-h}\int_0^av^{t_0}(1+v)^{h-1}dv$.  Then
$\eta(h+1)=\frac{h+1}{r(h+1+t_0)}\eta(h)\le\eta(h)$, and $1+v\le1+a=ra$
gives $\tau(h+1)\le\frac{h+1}h\,a\,\tau(h)\le\frac32a\,\tau(h)\le\tau(h)$
for $r\ge\frac52$.
\end{proof}

\begin{proposition}[Rows with one kept hole]\label{prop:oneholerows}
Let $r\ge2$, $1\le c\le n$ and $A\subset\N_{>0}$ finite.  Some finite
$Z\supseteq A$ has $|Z|\le|A|+n-1$ and
$\Phi_r(Z)\le\sup_{h\ge c}S_r(h,n-c)$.  The supremum is $S_r(c,n-c)$ when
$r\ge\frac52$, and also when $c=n$, where it is $a^n$.
\end{proposition}

\begin{proof}
Let $c_1<c_2<\cdots$ be the positive integers outside $A$, keep the hole
$h=c_c\ge c$, fill the holes $c_i$, $i\le n$, $i\ne c$, and put $x=c_n$ and
$Z=([1,x]\setminus\{h\})\cup A$.  Then $|Z\setminus A|=n-1$, and deleting the
elements of $Z$ above $h+n-c\le x$ one at a time (Lemma~\ref{lem:confdelr})
gives $\Phi_r(Z)\le S_r(h,n-c)$.  For $r\ge\frac52$,
Lemma~\ref{lem:shift} applied $h-c$ times gives $S_r(h,n-c)\le
S_r(c,n-c)$; for $c=n$, $S_r(h,0)=a^h\le a^n$.
\end{proof}

\begin{theorem}[Every row at the root $r$]\label{thm:everyrowr}
Let $r\ge2$, $L\ge1$, let $F$ be a monic multiple of $(x-r)^L$, $1\le k\le L$
and $n=L-k+1$.  Then
\[
  b_k(F)\ \ge\ \max\Bigl\{(r-1)^n,\ \frac{nr(r-1)^n}{(r-1)^n+1}\Bigr\},
\]
and if $r\ge\frac52$,
\[
  b_k(F)\ \ge\ B_r(n):=\max_{1\le c\le n}\frac1{S_r(c,n-c)},\qquad
  S_r(c,n-c)=\frac{r^{-c}}{\binom nc}
   +c\,r^{-c}\sum_{i=0}^{c-1}\binom{c-1}i\frac{(r-1)^{-(n-c+i+1)}}{n-c+i+1}.
\]
\end{theorem}

\begin{proof}
Let $S$ be as in Proposition~\ref{prop:rowvalue}(a), $|S|=k-1$.
Proposition~\ref{prop:psirows} with $A=S$ gives $Z\supseteq S$,
$|Z|\le L-1$, $\Phi_r(Z)\le\psi_r(n)$, and $1/\psi_r(n)$ is the second
term; the first is \cite[Theorem~3.2]{coefficient-mass} (or Proposition~\ref{prop:oneholerows}
with $c=n$).  For $r\ge\frac52$ use Proposition~\ref{prop:oneholerows} for
each $c$.  In each case Proposition~\ref{prop:rowvalue}(a) concludes.
\end{proof}

At $r=2$ the first bound is $\max\{1,n\}=L-k+1$, which is
\cite[Theorem~4.16]{coefficient-mass}.  The terms $c=1$ and $c=n$ of $B_r(n)$ are the
two terms of the first bound, so $B_r(n)$ dominates it.  Both rows at the
bottom are exact.

\begin{corollary}[The last two rows]\label{cor:tworowsr}
Let $r\ge2$ and $L\ge2$.  Then $V_r(L,L)=r-1$ and
\[
  V_r(L,L-1)=(r-1)^2\max\Bigl\{1,\frac{2r}{r^2-2r+2}\Bigr\},
\]
the maximum being the second term exactly when $r\le2+\sqrt2$.
\end{corollary}

\begin{proof}
The first is Corollary~\ref{cor:lastrowr}.  For $n=2$ the lower bound is
Theorem~\ref{thm:everyrowr}, as $(r-1)^2+1=r^2-2r+2$.  For the upper bound,
$V_r(L,L-1)\le\beta_r(2)$ by Proposition~\ref{prop:rowvalue}(c), and
$\mu_r(\varnothing,2)$ is the minimum over $\omega\in\R$ of
$\sum_s|1-\omega s|r^{-s}$, a convex function of $\omega$ that is linear
between consecutive points $1/s$, exceeds its value $a$ at $\omega=0$ for
$\omega<0$, and tends to infinity; so the minimum is $a$ or some
$\Phi_r(\{z\})$.  Summing directly,
\[
  \Phi_r(\{z\})=\frac{z(r-1)-r+2r^{1-z}}{z(r-1)^2}=a-\frac{a^2}z\,g(z),\qquad
  g(z)=r-2r^{1-z}.
\]
Now $g(1)=r-2$, $g(2)/2-g(3)/3=(r-2)(r^2+2r-2)/(6r^2)\ge0$, and for $z\ge4$,
$g(z)/z<r/4\le r/2-1/r=g(2)/2$.  So the minimum is $\min\{\Phi_r(\{1\}),
\Phi_r(\{2\})\}=\min\{a^2,(r^2-2r+2)a^2/(2r)\}$, which is at most $a$, and
$\beta_r(2)$ is the stated value.  The comparison $2r\ge r^2-2r+2$ is
$r^2-4r+2\le0$.
\end{proof}

At $r=2$ the corollary says $V_2(L,L-1)=2$: \cite[Theorem~4.16]{coefficient-mass} is
sharp in its two lowest rows.  The threshold $2+\sqrt2$ is where the
criterion quoted after \cite[Proposition~4.7]{coefficient-mass} starts to hold for
$P_{L-2}=(x-r)^2$: its first partial sum $1-2r$ has absolute value at most
$(r-1)^2$ exactly when $r^2-4r+2\ge0$.

\paragraph{How far the bound is from the rows.}
By Proposition~\ref{prop:rowvalue}, $B_r(n)\le V_r(L,k)\le\beta_r(n)$ for
$r\ge\frac52$, and $\beta_r(n)$ is the top row of $(x-r)^n$.  We computed
$1/\Phi_r(Z_0)$ exactly for the best zero set $Z_0$ of size $n-1$ inside
$[1,\min(3n+3,18)]$; this is a lower bound for $\beta_r(n)$, and presumably
its value.  It equals $B_r(n)$ for $r=6$, $n\le7$; $r=4$, $n\le5$; $r=3$,
$n\le3$; $r=\frac52$, $n\le2$.  For example $B_3(3)=\frac{27}2$ at $c=2$ and
$B_4(4)=128$ at $c=3$.  Beyond that the two separate, for instance
$B_3(5)=80$ against $\frac{31104}{307}=101.3\ldots$ and $B_3(7)=407.2\ldots$
against $815.8\ldots$; in the computed range the ratio decreases with $n$.
At $r=2$ the computed values are $1,2,\frac{40}{11},\frac{672}{115},
\frac{10368}{851}$ for $n\le5$, against $n$.  The one-hole family is not
optimal in general: the best zero sets of the top row, such as
$\{1,3,5\}$ at $(r,n)=(3,4)$, spread their gaps.  Numerically
Lemma~\ref{lem:shift} holds for all $h,t_0<50$ at $r=2.3$ and fails at
$r=2.2$ ($h=3$, $t_0=13$), so the constant $\frac52$ is not the limit of
the method.

\paragraph{Rows equal the top row.}
For $r\ge2$ we conjecture $V_r(L,k)=\beta_r(L-k+1)$, and
Theorem~\ref{thm:allrowstwo} proves it: exempting $k-1$
positions exactly offsets $k-1$ of the $L$ zeros at $r$.  In the form of
Proposition~\ref{prop:rowvalue}(b), this says that every finite $S$ admits
$Z\supseteq S$ with $|Z|\le|S|+n-1$ and $\Phi_r(Z)\le\mu_r(\varnothing,n)$.
It holds for $n=1$ by Lemma~\ref{lem:confdelr}, and for $n=2$ by
Corollary~\ref{cor:tworowsr}.  For $r\in\{2,\frac94,\frac52,3,4\}$,
$n\le4$ and every $S\subseteq[1,9]$ with $|S|\le3$ we found such a $Z$,
with $\mu_r(\varnothing,n)$ replaced by $\Phi_r(Z_0)$ for the $Z_0$ of the
previous paragraph and the inequality checked in exact arithmetic.  The
worst $S$ found always lies at the top of the search range, as the limit in
Proposition~\ref{prop:rowvalue}(c) suggests.  At $r=2$ the conjecture would
sharpen \cite[Theorem~4.16]{coefficient-mass} to $b_k\ge\beta_2(L-k+1)$, with
$\beta_2(3)\ge\frac{40}{11}>3$.  It fails for $r<2$: there
$V_r(L,L)=(r-1)^L<r-1=\beta_r(1)$ by Corollary~\ref{cor:lastrowr}.

\paragraph{Roots between one and two.}
For $1<r<2$ the deletion of Lemma~\ref{lem:confdelr} can increase $\Phi_r$
by the factor $a>1$, and the exempted set $S=[1,k-1]$ is costly: every
$Z\supseteq[1,k-1]$ has
$\Phi_r(Z)=a^k\sum_{s\ge k}\binom{s-1}{k-1}(1-r^{-1})^kr^{-(s-k)}|q(s)|$
with $q=r_{Z\setminus[1,k-1]}$, a negative-binomial mean inflated by $a^k$.
What is proved here is the top row, $V_r(L,1)=\beta_r(L)\ge\max_c1/S_r(c,L-c)$
(any certificate is admissible when $S=\varnothing$), and the last row,
$V_r(L,L)=(r-1)^L$; Corollary~\ref{cor:ordergen} gives only $(r-1)^L$ for
the rows in between.  In a numerical search at $r=\frac32$ and $\frac74$
with $L\le5$, $S\subseteq[1,8]$ and zero sets in a window, the worst
exempted set is $S=[1,k-1]$ in all but one case, and the approximate row
values at $r=\frac32$, $L=5$ are $2.98,\,1.99,\,0.80,\,0.20,\,\frac1{32}$
for $k=1,\ldots,5$ (the last is exact).  Theorem~\ref{thm:prefixrows} below
determines these rows up to a factor tending to $1$, and
Theorems~\ref{thm:onepolyrows} and~\ref{thm:longdiag} determine many of
them exactly.

\section{Prefix sets, several roots, and the top row}\label{sec:prefix}

This section keeps the notation of the previous one.  It adds three
things.  For $1<r<2$ it identifies the rows up to a factor tending to $1$
(Theorem~\ref{thm:prefixrows}, Corollary~\ref{cor:belowtwo}).  For several
roots it evaluates the one-hole certificate in closed form and derives a
row bound (Proposition~\ref{prop:oneholegen}, Theorem~\ref{thm:rowsgen}).
Finally it proves that rows equal the top row at $r\ge\frac52$ whenever an
explicit criterion holds (Theorem~\ref{thm:rowstop}), and it reduces the
general conjecture to a single insertion inequality
(Proposition~\ref{prop:insertion}).

\paragraph{Rows through prefix sets.}
For $i\ge1$ and $n\ge1$ put
\[
  \nu_i(n)=\inf\Bigl\{\sum_{s\ge1}\binom{s-1}{i-1}r^{-s}|p(s)|:\
  p\in\R[s],\ \deg p\le n-1,\ p(0)=1\Bigr\},
\]
so $\nu_1(n)=\mu_r(\varnothing,n)=1/\beta_r(n)$.  Since
$|r_{[1,i-1]}(s)|=\binom{s-1}{i-1}$ for $s\ge1$, we have
$\nu_i(n)=\mu_r([1,i-1],i+n-1)$.  In the negative-binomial form of the
paragraph on roots between one and two, $\nu_i(n)=a^im_i(n)$, where
$m_i(n)=\inf_p\mathbb E|p(X_i)|$ and $X_i$ counts the trials up to the
$i$-th success at success probability $1-1/r$.  So $1/\nu_i(n)$ is a top-row
value for a transformed measure.  Put also
\[
  W_k(s)=\binom{s-1}{\min(k-1,\lfloor(s-1)/2\rfloor)},\qquad
  \Omega_k(n)=\inf_p\sum_{s\ge1}W_k(s)r^{-s}|p(s)|,\qquad
  E_k(r)=\sum_{s=1}^{2k-2}\binom{s-1}{\lfloor(s-1)/2\rfloor}r^{-s},
\]
with $p$ as in $\nu_i$.  Since the binomial coefficients are unimodal,
$W_k(s)=\max\{\binom{s-1}j:0\le j\le\min(k-1,s-1)\}$.

\begin{lemma}[Integer zero sets suffice]\label{lem:intzeros}
Let $i,n\ge1$ and let $\omega(s)\ge0$ ($s\ge1$) vanish for $s<i$, be
positive for $s\ge i$, and satisfy $\sum_s\omega(s)s^{n-1}<\infty$.  Then
$\inf_p\sum_s\omega(s)|p(s)|$, over real $p$ with $\deg p\le n-1$ and
$p(0)=1$, equals the infimum over the polynomials
$p=\prod_{y\in Y}(1-s/y)$ with $Y$ a set of $n-1$ integers at least $i$.
\end{lemma}

\begin{proof}
Let $D\ge i+n$ and $f_D(p)=\sum_{s=i}^D\omega(s)|p(s)|$ on the affine space
$\mathcal A=\{\deg p\le n-1,\ p(0)=1\}$ of dimension $n-1$.  On each cell
$\{p\in\mathcal A:\ \epsilon_sp(s)\ge0,\ i\le s\le D\}$
($\epsilon_s=\pm1$), $f_D$ is linear and nonnegative.  Each cell is
pointed, because a polynomial of degree at most $n-1$ that vanishes at $0$
and at the $D-i+1\ge n$ points $[i,D]$ is zero.  So the minimum of $f_D$ over
a cell is attained at a vertex.  A vertex satisfies $n-1$ linearly
independent equations $p(y)=0$ with $y\in[i,D]$, so
$p=\prod_{y\in Y}(1-s/y)$ with $|Y|=n-1$.

Hence some such $p_D$ has $f_D(p_D)=\min_{\mathcal A}f_D\le
\sum_s\omega(s)|p(s)|$ for all $p\in\mathcal A$.  Since
$|p_D(s)|\le(1+s)^{n-1}$, the full sum at $p_D$ exceeds $f_D(p_D)$ by at
most $\tau_D=\sum_{s>D}\omega(s)(1+s)^{n-1}\to0$.  So the infimum lies
between $\min f_D$ and $\min f_D+\tau_D$, and the truncated minima increase
to it.
\end{proof}

\begin{lemma}[A uniform bound on the exempted factor]\label{lem:rowpointwise}
Let $S\subset\N_{>0}$ be finite, $s\ge1$, and $j_s=|S\cap[1,s-1]|$.  Then
$|r_S(s)|\le\binom{s-1}{j_s}$.  If $|S|\le k-1$, then
$|r_S(s)|\le W_k(s)$, and $W_k(s)=\binom{s-1}{k-1}$ for $s\ge2k-1$.
\end{lemma}

\begin{proof}
Let $\sigma_1<\cdots<\sigma_{j}$ be the elements of $S$ below $s$, so
$\sigma_l\ge l$ and $j=j_s\le s-1$.  The factor $(s-x)/x$ decreases in $x$, so
$|1-s/\sigma_l|\le(s-l)/l$.  Every $\sigma\in S$ with $\sigma\ge s$ has
$0\le1-s/\sigma<1$.  Hence
$|r_S(s)|\le\prod_{l\le j}(s-l)/l=\binom{s-1}j$.  The rest follows from
$j_s\le\min(k-1,s-1)$, unimodality, and the fact that
$\lfloor(s-1)/2\rfloor\ge k-1$ exactly when $s\ge2k-1$.
\end{proof}

\begin{theorem}[Rows through prefix sets]\label{thm:prefixrows}
Let $r>1$, $1\le k\le L$ and $n=L-k+1$.  Then
\begin{equation}\label{eq:prefixsandwich}
  \frac1{\nu_k(n)+E_k(r)}\ \le\ \frac1{\Omega_k(n)}\ \le\ V_r(L,k)\ \le\
  \min_{1\le i\le k}\frac1{\nu_i(n)} .
\end{equation}
\end{theorem}

\begin{proof}
\emph{Upper bound.}  Fix $1\le i\le k$ and, for large $N$, put
$S_N=[1,i-1]\cup[N+1,N+k-i]$, so $|S_N|=k-1$.  A real $q$ with
$\deg q\le L-1$, $q(0)=1$ and $q|_{S_N}=0$ factors as
$q=r_{[1,i-1]}\pi_Nq_0$ with $\pi_N(s)=\prod_{\sigma=N+1}^{N+k-i}(1-s/\sigma)$,
$\deg q_0\le n-1$ and $q_0(0)=1$.  For $s\le M<N$ we have
$\pi_N(s)\ge(1-M/N)^{k-i}$, so
\[
  \Phi_r(q)\ge\Bigl(1-\frac MN\Bigr)^{k-i}\nu_{i,M}(n),
\]
where $\nu_{i,M}$ is $\nu_i$ with the sum cut at $s\le M$.
Proposition~\ref{prop:rowvalue}(b) gives
$V_r(L,k)\le1/\mu_r(S_N,L)$.  Let $N\to\infty$, then $M\to\infty$.  The
proof of Lemma~\ref{lem:intzeros} gives $\nu_{i,M}(n)\to\nu_i(n)$, so
$V_r(L,k)\le1/\nu_i(n)$.

\emph{Lower bound.}  Let $F$ and $S$ be as in
Proposition~\ref{prop:rowvalue}(a), so $|S|=k-1$.  For any $p$ as in
$\Omega_k$, the polynomial $q=r_Sp$ is admissible there, and by
Lemma~\ref{lem:rowpointwise}
\[
  \Phi_r(q)\le\sum_sW_k(s)r^{-s}|p(s)| .
\]
Taking the infimum over $p$ gives $b_k(F)\ge1/\Omega_k(n)$.

Now take $p=\prod_{y\in Y}(1-s/y)$ with $|Y|=n-1$ and $Y\subset[k,\infty)$.
For $1\le s\le2k-2\le2y$ every factor satisfies $|1-s/y|\le1$, so
$|p(s)|\le1$.  For $s\ge2k-1$, $W_k(s)=\binom{s-1}{k-1}$.  Hence
\[
  \sum_sW_k(s)r^{-s}|p(s)|\le E_k(r)+\sum_{s\ge1}\binom{s-1}{k-1}r^{-s}|p(s)| .
\]
By Lemma~\ref{lem:intzeros} with $\omega(s)=\binom{s-1}{k-1}r^{-s}$ and
$i=k$, the infimum of the last sum over such $p$ is $\nu_k(n)$.  This gives
$\Omega_k(n)\le\nu_k(n)+E_k(r)$.
\end{proof}

The upper bound says that the prefix sets $S=[1,i-1]$, completed by a far
block, are admissible exempted sets.  The lower bound compares every $S$
with the prefix $[1,k-1]$.  The loss $E_k(r)$ comes from the first $2k-2$
positions only.  For $r<2$ it is exponentially small against $\nu_k(n)$.

\begin{lemma}[The negative-binomial mass]\label{lem:nbmass}
Let $r>1$, $k,n\ge1$, $\varpi=1-1/r$, and let $M_k$ be a mode of the
negative-binomial law $P(t)=\binom{t-1}{k-1}\varpi^k r^{-(t-k)}$, $t\ge k$.
Then $M_k\le(k-1)r/(r-1)+1$, and
\[
  \frac{\nu_k(n)}{a^k}=m_k(n)\ \ge\
  \frac{(n-1)!\,P(M_k)}{r^{n-1}\bigl(2(M_k+n-1)\bigr)^{n-1}},
  \qquad
  P(M_k)\ge\frac3{4\bigl(4\sqrt{kr}/(r-1)+1\bigr)} .
\]
\end{lemma}

\begin{proof}
First, $\binom{s-1}{k-1}r^{-s}=a^kP(s)$, which gives $\nu_k=a^km_k$.  The
ratio $P(t+1)/P(t)=t/(r(t-k+1))$ is at least $1/r$.  It is at least $1$
exactly when $t\le(k-1)/\varpi$, which locates the mode.

Let $t_l=M_k+l$ for $0\le l\le n-1$, and let $\ell_l$ be the Lagrange basis
on these nodes.  Then
$|\ell_l(0)|=\prod_{l'\ne l}t_{l'}/|l-l'|\le(M_k+n-1)^{n-1}/(l!\,(n-1-l)!)$.
Evaluating $1=p(0)=\sum_l\ell_l(0)p(t_l)$ gives
$\max_l|p(t_l)|\ge(n-1)!/(2(M_k+n-1))^{n-1}$.  Moreover
$P(t_l)\ge r^{-(n-1)}P(M_k)$.

The law has variance $\sigma^2=kr/(r-1)^2$.  By Chebyshev's inequality at
least $\frac34$ of its mass lies within $2\sigma$ of the mean, on at most
$4\sigma+1$ integers.  So the mode carries at least $3/(4(4\sigma+1))$.
\end{proof}

\begin{corollary}[Rows for roots between one and two]\label{cor:belowtwo}
Let $1<r<2$ and $\theta=4(r-1)/r^2<1$.  For all $1\le k\le L$ and $n=L-k+1$,
\[
  1\ \ge\ \frac{V_r(L,k)\,m_k(n)}{(r-1)^k}\ \ge\
  \frac1{1+\varepsilon},\qquad
  \varepsilon\le\frac{r^2\,\theta^k}{4(2-r)\,m_k(n)} .
\]
With Lemma~\ref{lem:nbmass}, $\varepsilon\le C_{r,n}k^{n-1/2}\theta^k$.  So
for fixed $n$ (indeed for $n=o(k/\log k)$),
$V_r(L,k)=(1+o(1))(r-1)^k/m_k(n)$ as $k\to\infty$, and the exempted set
$S=[1,k-1]$ is asymptotically the worst.
\end{corollary}

\begin{proof}
Theorem~\ref{thm:prefixrows} with $i=k$ gives the upper bound, and
$1/(\nu_k+E_k)=(1/\nu_k)/(1+E_k/\nu_k)$ gives the lower one.  Since
$\binom{s-1}{\lfloor(s-1)/2\rfloor}\le2^{s-1}$ and $2/r>1$,
\[
  E_k(r)\le\frac12\sum_{s=1}^{2k-2}\Bigl(\frac2r\Bigr)^s\le
  \frac{(2/r)^{2k-1}}{2(2/r-1)}=\frac{r^2}{4(2-r)}\Bigl(\frac4{r^2}\Bigr)^k .
\]
Finally $E_k/\nu_k=(r-1)^kE_k/m_k$, and $\theta<1$ is $(r-2)^2>0$.
\end{proof}

For $r\ge2$ the loss $E_k(r)$ is not small and
Theorem~\ref{thm:prefixrows} adds nothing to Theorem~\ref{thm:everyrowr}.
At the same time $\theta\to1$ as $r\to2$, and for small $L$ the bound of
Corollary~\ref{cor:ordergen} is better.  We conjecture
$V_r(L,k)=\min_{1\le i\le k}1/\nu_i(n)$ for every $r>1$;
Theorem~\ref{thm:allrowstwo} proves this for $r\ge2$.  For $1<r<2$ the evidence is a
floating-point linear-programming search over exempted sets.  For
$r\in\{1.1,1.25,1.5,1.75,1.9\}$, $n\le4$ and all $S\subseteq[1,10]$ with
$|S|\le3$, no $S$ gave a value below the upper bound
of~\eqref{eq:prefixsandwich}.

At $r\ge2$ the conjecture says the minimum sits at $i=1$, that is, rows
equal the top row.  At $r=\frac32$ and $L\le8$ the minimum is at $i=k$
except for $(L,k)=(6,2),(8,2),(8,3)$, where it is at $i=1$.  For $n\ge2$
at $r=1.9$ and $n\ge3$ at $r=1.75$ it is at $i=1$.  This explains the
exception noted in the previous paragraph: there the worst exempted set lies
far out.  For $1<r<2$, Theorem~\ref{thm:onepolyrows} proves the conjecture
whenever $\beta_r(n)\le1$, for every $k$ and with the minimum at $i=k$.
Theorem~\ref{thm:longdiag} proves it for every $n$ once $k$ is large
compared with $n$, again with the minimum at $i=k$
(Corollary~\ref{cor:longdiagfinite}).  So for each fixed $n$ and $1<r<2$
the minimum moves to $i=k$ for all large $k$.

\paragraph{One kept hole at several roots.}
For a multiset $R=(\rho_1,\ldots,\rho_m)$ of reals greater than $1$ write
$e_j(R)$ for its elementary symmetric functions and $R-1=(\rho_i-1)_i$.  Put
\[
  Q(R)=\prod_{i=1}^m(\rho_i-1),\qquad
  P_h(R)=\sum_{i=0}^{h-1}(-1)^ie_i(R),
\]
and, for $1\le h\le m$,
\begin{equation}\label{eq:Tgen}
  T_h(R)=\frac{Q(R)+\sum_{i=0}^{h-1}\binom{m-1-i}{h-1-i}e_i(R-1)}
             {Q(R)\,e_h(R)} .
\end{equation}

\begin{lemma}[Partial sums]\label{lem:partialsums}
For $1\le h\le m$,
$(-1)^{h-1}P_h(R)=\sum_{i=0}^{h-1}\binom{m-1-i}{h-1-i}e_i(R-1)>0$.  Also
$P_{m+1}(R)=\prod_i(1-\rho_i)$, so $|P_{m+1}(R)|=Q(R)$.
\end{lemma}

\begin{proof}
$P_h(R)$ is the coefficient of $u^{h-1}$ in $\prod_i(1-\rho_iu)/(1-u)$.
Writing $1-\rho_iu=(1-u)-(\rho_i-1)u$ gives
\[
  \frac{\prod_i(1-\rho_iu)}{1-u}=
  \sum_{i=0}^m(-1)^ie_i(R-1)\,u^i(1-u)^{m-1-i}.
\]
For $h\le m$ only $i\le h-1\le m-1$ contributes to $u^{h-1}$, with
$(-1)^i\cdot(-1)^{h-1-i}\binom{m-1-i}{h-1-i}$.  The last claim is the value
of $\prod_i(1-\rho_iu)$ at $u=1$.
\end{proof}

\begin{proposition}[One hole at arbitrary roots]\label{prop:oneholegen}
Let $\rho_1,\ldots,\rho_m>1$ be distinct, $R=(\rho_i)$, $1\le h\le m$, and
let $w$ be the exponential sum on the nodes $y_i=1/\rho_i$ with $w_0=1$ and
$w_z=0$ for $z\in[1,m]\setminus\{h\}$.  Then $|w_h|=1/e_h(R)$, the terms
$w_d$ with $d>m$ have one sign, and
\[
  \Tail(w)=\frac1{e_h(R)}\Bigl(1+\frac{|P_h(R)|}{Q(R)}\Bigr)=T_h(R).
\]
For $R=(r,\ldots,r)$ the right side equals
$S_r(h,m-h)=\Phi_r([1,m]\setminus\{h\})$ of Lemma~\ref{lem:holeintr}.
\end{proposition}

\begin{proof}
By the zero bound, $w$ exists, is unique, and has exactly the $m-1$
prescribed zeros, all of them simple.  So $w$ has one sign on $(m,\infty)$.
Let $E(X)=\prod_i(1-y_iX)=\sum_j(-1)^je_j(y)X^j$ and
$U(X)=\sum_{d\ge0}w_dX^d=N(X)/E(X)$ with $\deg N\le m-1$, and write
$U=1+w_hX^h+T$ with $T=\sum_{d>m}w_dX^d$.

The coefficient of $X^m$ in $N=(1+w_hX^h)E+TE$ vanishes, so
$w_h=(-1)^{h+1}e_m(y)/e_{m-h}(y)$.  The coefficients beyond $X^m$ vanish
as well, so
$T=-w_hX^h\sum_{j>m-h}(-1)^je_j(y)X^j\big/E(X)$.  This is analytic for
$|X|<1/\max y_i$, so it may be evaluated at $X=1$.  By one-signedness,
\[
  \Tail(w)=|w_h|+|T(1)|
  =|w_h|\Bigl(1+\Bigl|\sum_{j>m-h}(-1)^je_j(y)\Bigr|\Big/E(1)\Bigr).
\]
With $e_j(y)=e_{m-j}(R)/\prod_i\rho_i$ this reads
\[
  |w_h|=\frac1{e_h(R)},\qquad
  \Bigl|\sum_{j>m-h}(-1)^je_j(y)\Bigr|=\frac{|P_h(R)|}{\prod_i\rho_i},\qquad
  E(1)=\frac{Q(R)}{\prod_i\rho_i},
\]
and Lemma~\ref{lem:partialsums} converts $|P_h(R)|$ into the numerator
of~\eqref{eq:Tgen}.

For a single root, take the nodes $y_i=\frac1r-(m-i)\epsilon$ of the proof
of \cite[Lemma~4.9]{coefficient-mass}.  Its limit gives
$\Tail(w^\epsilon)\to\Phi_r([1,m]\setminus\{h\})$ (as in
Lemma~\ref{lem:confdelr}), while $T_h$ is continuous.
\end{proof}

The case $h=1$ is $T_1(R)=(1+Q(R))/(Q(R)\sum_i\rho_i)$, the several-root
form of $\psi_r(m)$.  The case $h=m$ is $T_m(R)=1/Q(R)$, which is
\cite[Theorem~3.2]{coefficient-mass} with $u=0$.

\begin{theorem}[Rows at several roots]\label{thm:rowsgen}
Let $2\le r_1\le\cdots\le r_L$, let $F$ be a nonzero real multiple of
$\prod_{i=1}^L(x-r_i)$, $1\le k\le L$ and $n=L-k+1$, and let
$R_m=(r_{L-m+1},\ldots,r_L)$ be the $m$ largest roots.  Then
\[
  b_k(F)\ \ge\ |f_D|\max_{1\le c\le n}\ \min_{c\le h\le c+k-1}
  \frac1{T_h(R_{h+n-c})} .
\]
The term $c=n$ is \cite[Theorem~3.2]{coefficient-mass}.  For $r_1=\cdots=r_L=r\ge\frac52$ the
right side is $B_r(n)$.
\end{theorem}

\begin{proof}
Let the roots be distinct and $f_D=1$, and let $W$ be the set of backward
distances of $k-1$ positions carrying the $k-1$ largest nonleading
magnitudes.  Let $c_1<c_2<\cdots$ be the positive integers outside $W$ and
fix $c$.  Put
\[
  Z=W\cup\{c_i:\ i\le n,\ i\ne c\},\qquad h=c_c,\qquad x=c_n,
\]
so $|Z|=L-1$ and $c\le h\le c+k-1$.  Let $w$ be the sum on all $L$
reciprocal roots with $w_0=1$ and $w|_Z=0$.  The dual identity
of \cite[Lemma~3.1]{coefficient-mass} gives $1\le b_k(F)\Tail(w)$,
because $w$ vanishes at the $k-1$ exempted distances.

Now delete the largest zero repeatedly, each time together with the largest
remaining node, that is, the smallest remaining root.  First delete the
elements of $W$ above $x$, then $x,x-1,\ldots,h+n-c+1$.  Every set met
before a deletion contains an element larger than $h$ and misses $h$, so it
is not an initial segment.  The deletion step in the proof of
\cite[Theorem~2.2]{coefficient-mass} (nodes at most $\frac12$) makes the tail smaller at
each step.  The last set is $[1,h+n-c]\setminus\{h\}$ on the reciprocals of
$R_{h+n-c}$, since $x\ge h+n-c$.  By Proposition~\ref{prop:oneholegen},
$\Tail(w)\le T_h(R_{h+n-c})$.

Repeated roots follow by the perturbation in the proof of
\cite[Theorem~3.2]{coefficient-mass}, because $T_h$ is continuous in the roots.

For $c=n$ the value is $1/T_h(R_h)=Q(R_h)\ge Q(R_n)$, since $h\ge n$ and
every factor is at least $1$.  For equal roots,
$T_h(R_{h+n-c})=S_r(h,n-c)$, and Lemma~\ref{lem:shift} gives
$\max_hS_r(h,n-c)=S_r(c,n-c)$.
\end{proof}

At distinct roots the minimum over $h$ need not sit at $h=c$.  The one step
that is proved in general is the move of a hole from $2$ to $1$.

\begin{lemma}[Moving the hole from $2$ to $1$]\label{lem:holetwo}
Let $R$ be a multiset of $m\ge2$ reals, all at least $2$, and let $R'$ be
$R$ with one element $\rho$ removed.  Then $T_2(R)\le T_1(R')$.
\end{lemma}

\begin{proof}
Write $s=\rho-1\ge1$, $Q'=Q(R')\ge1$, $E_1=e_1(R')$ and $E_2=e_2(R')$.  By
\eqref{eq:Tgen}, $T_2(R)=(sQ'+E_1+s)/(sQ'(E_2+(s+1)E_1))$ and
$T_1(R')=(Q'+1)/(Q'E_1)$.  After clearing denominators the claim is
$E_1^2\le s(Q'+1)E_2+s^2(Q'+1)E_1$, and with $s\ge1$ it suffices that
$E_1^2\le(Q'+1)(E_1+E_2)$.

Now $E_1^2=\sum_j\rho_j'^2+2E_2$ and $(Q'+1)E_2\ge2E_2$.  Also
$\sum_j\rho_j'^2\le(\max_j\rho_j')E_1\le(Q'+1)E_1$, because
$Q'\ge\max_j\rho'_j-1$ when all factors are at least $1$.
\end{proof}

\begin{corollary}[Several roots, away from the top two positions]\label{cor:rowsgenfree}
In the setting of Theorem~\ref{thm:rowsgen}, suppose that the $k-1$ largest
nonleading magnitudes can be chosen at positions not including both $D-1$
and $D-2$.  Then
\[
  b_k(F)\ \ge\ |f_D|\,\frac{\bigl(\sum_{i=k}^Lr_i\bigr)\,Q}{Q+1},
  \qquad Q=\prod_{i=k}^L(r_i-1).
\]
\end{corollary}

\begin{proof}
In the proof of Theorem~\ref{thm:rowsgen} take $c=1$.  Then $h=c_1\le2$.
If $h=1$ the bound is $1/T_1(R_n)$.  If $h=2$, Lemma~\ref{lem:holetwo}
applied to $R_{n+1}$ and its smallest root gives
$T_2(R_{n+1})\le T_1(R_n)$.
\end{proof}

At $r_i=2$ this is $b_k\ge n$ for such positions, as in
\cite[Proposition~4.12]{coefficient-mass}.  For equal roots it is the target
$\psi_r(n)$ of Proposition~\ref{prop:psirows} for $j\le1$.  The general
step $T_h(R)\le T_{h-1}(R')$ ($R'$ without the smallest root) would make
the bound of Theorem~\ref{thm:rowsgen} equal to
$\max_c1/T_c(R_n)$, the several-root form of $B_r(n)$.  In about $15{,}000$
random exact tests with smallest root at least $\frac52$ it held, with
largest ratio $0.75$, attained near equal roots.  It fails
at equal roots near $2.2$, as Lemma~\ref{lem:shift} does.

\paragraph{The top row, and rows equal to it.}
Put $\beta(R)=\inf b_1(F)$ over monic real multiples $F$ of
$\prod_i(x-\rho_i)$, so $\beta_r(n)=\beta(r,\ldots,r)$.  With
$A_z=P_{z+1}(R)$ for $0\le z\le m$ (so $A_0=1$ and $|A_m|=Q(R)$), consider
the $m+1$ affine functions
\[
  \ell_0(\lambda)=A_m(1+\lambda),\qquad
  \ell_z(\lambda)=A_z+\lambda A_{z-1}\quad(1\le z\le m),
\]
and write $\ell_z=\alpha_z+\lambda\gamma_z$.  Every $\gamma_z\neq0$ by
Lemma~\ref{lem:partialsums}.  Put
$X(z,z')=|\alpha_z\gamma_{z'}-\alpha_{z'}\gamma_z|/(|\gamma_z|+|\gamma_{z'}|)$.

\begin{proposition}[The top row at arbitrary roots]\label{prop:toprowgen}
For every multiset $R$ of $m$ reals greater than $1$,
\[
  \max_{1\le c\le m}\frac1{T_c(R)}\ \le\ \beta(R)\ \le\
  U(R):=\min_{\lambda\in\R}\max_{0\le z\le m}|\ell_z(\lambda)|
  =\max_{0\le z<z'\le m}X(z,z'),
\]
and $X(0,c)=1/T_c(R)$.  Hence $\beta(R)=\max_c1/T_c(R)$ whenever
$X(z,z')\le\max_c1/T_c(R)$ for all $1\le z<z'\le m$.
\end{proposition}

\begin{proof}
\emph{Lower bound.}  For distinct roots, the dual identity of \cite[Lemma~3.1]{coefficient-mass}
with the one-hole sum of Proposition~\ref{prop:oneholegen} gives
$1\le b_1(F)T_c(R)$.  Repeated roots follow by perturbation.

\emph{Upper bound.}  Fix $\lambda$ and let
\[
  N(y)=(1+\lambda y)\prod_i(1-\rho_iy)=\sum_{j=0}^{m+1}N_jy^j .
\]
Its partial sums are $\sum_{j\le z}N_j=\ell_z(\lambda)$ for $1\le z\le m$
and $\ell_0(\lambda)=N(1)$ for $z\ge m+1$.  For $K\ge m+1$ let
$h_j=1$ for $0\le j\le K$, $h_j=(2K-j)/K$ for $K\le j\le2K$ and $h_j=0$
beyond, and let $\widetilde G(y)=N(y)\sum_jh_jy^j=\sum_sg_sy^s$.

The polynomial $F(x)=x^{\deg\widetilde G}\widetilde G(1/x)$ is monic,
because $g_0=1$.  It is divisible by $\prod_i(x-\rho_i)$, repetitions
included, because $N$ is divisible by $\prod_i(1-\rho_iy)$.  Its nonleading
coefficients are $g_1,g_2,\ldots$.  For $s\le K$, $g_s$ is the partial sum
up to $s$, so it is $\ell_s(\lambda)$ or $\ell_0(\lambda)$.  For $s>K$,
$h_{s-j}=h_s+O(j/K)$ gives $|g_s|\le|\ell_0(\lambda)|+K^{-1}\sum_jj|N_j|$.
Letting $K\to\infty$ gives $\beta(R)\le\max_z|\ell_z(\lambda)|$ for every
$\lambda$.

\emph{The minimax.}  For two indices,
$|\alpha_z\gamma_{z'}-\alpha_{z'}\gamma_z|
=|\gamma_{z'}\ell_z(\lambda)-\gamma_z\ell_{z'}(\lambda)|
\le(|\gamma_z|+|\gamma_{z'}|)\max(|\ell_z|,|\ell_{z'}|)$, with equality at
the crossing point.  So $X(z,z')$ is the minimax of that pair.  The
sublevel sets $\{\lambda:|\ell_z(\lambda)|\le t\}$ are compact intervals,
and by Helly's theorem on the line they have a common point once they meet
pairwise; this gives $U(R)=\max X$.

Finally, $X(0,c)=Q|A_{c-1}-A_c|/(Q+|A_{c-1}|)=Qe_c(R)/(Q+|P_c(R)|)$, which
is $1/T_c(R)$ by Lemma~\ref{lem:partialsums}.
\end{proof}

At $\lambda=0$ the functions $\ell_z$ are the partial sums of the
coefficients of $\prod_i(x-\rho_i)$ and $\ell_0=\pm Q(R)$.  This recovers
the criterion quoted after \cite[Proposition~4.7]{coefficient-mass}: if every partial
sum has absolute value at most $Q(R)$, then $\beta(R)=Q(R)$.  For a single
root, write $X_r(z,z')$ for $X(z,z')$ at $R=(r,\ldots,r)$ of length $n$,
with $A_z=\sum_{j\le z}\binom nj(-r)^j$.

\begin{theorem}[Rows equal the top row]\label{thm:rowstop}
Let $r\ge\frac52$ and $n\ge1$, and suppose $X_r(z,z')\le B_r(n)$ for all
$1\le z<z'\le n$.  Then for every $L\ge n$ and $k=L-n+1$,
\[
  V_r(L,k)=\beta_r(n)=B_r(n)=\max_{1\le c\le n}\frac1{S_r(c,n-c)} .
\]
\end{theorem}

\begin{proof}
Theorem~\ref{thm:everyrowr} gives $V_r(L,k)\ge B_r(n)$, and
Proposition~\ref{prop:rowvalue}(c) gives $V_r(L,k)\le\beta_r(n)$.  By
Proposition~\ref{prop:toprowgen} with $T_c(r,\ldots,r)=S_r(c,n-c)$
(Proposition~\ref{prop:oneholegen}), $\beta_r(n)\le U=\max\{B_r(n),
\max X_r(z,z')\}=B_r(n)$.
\end{proof}

\begin{corollary}[Explicit ranges]\label{cor:rowstopranges}
The hypothesis of Theorem~\ref{thm:rowstop}, and hence
$V_r(L,L-n+1)=\beta_r(n)=B_r(n)$ for all $L$, holds for
\[
  n=3,\ r\ge\tfrac{13}5;\quad n=4,\ r\ge\tfrac{16}5;\quad n=5,\ r\ge\tfrac{39}{10};\quad
  n=6,\ r\ge\tfrac{47}{10};\quad n=7,\ r\ge\tfrac{53}{10};\quad n=8,\ r\ge6 .
\]
\end{corollary}

\begin{proof}
With the known sign $(-1)^z$ of $A_z$ (Lemma~\ref{lem:partialsums}), each
inequality $X_r(z,z')\le1/S_r(c,n-c)$ is a pair of polynomial inequalities
in $r$.  For each $n$ and each pair $z<z'$, the real roots above the stated
threshold of all these polynomials cut the half-line into intervals.  No
polynomial changes sign inside an interval, and on each interval one $c$
satisfies both inequalities at a rational sample point.  This is a finite
exact computation (real-root isolation over $\mathbb Q$).
\end{proof}

For $n\le2$ the rows are Corollary~\ref{cor:tworowsr}.  In floating point
the least $r$ with the hypothesis is about $2.60$, $3.19$, $3.85$, $4.6$,
$5.2$, $5.9$, $6.6$, $7.2$ for $n=3,\ldots,10$.  The best $c$ is $n-2$,
$n-1$ or $n$ there.  Above roughly $r\approx n$ it is $c=n$, where
$\beta_r(n)=(r-1)^n$ and the rows follow already from \cite[Theorem~3.2]{coefficient-mass}.
Below the threshold $U_r(n)>B_r(n)$.  There the optimal top-row zero sets
have two or more holes (for instance $\{1,2,4,6\}$ at $(r,n)=(3,5)$), and
$B_r(n)\le V_r\le\beta_r(n)\le U_r(n)$ is all we know.

The general conjecture reduces to one inequality.  For finite
$Z\subset\N_{>0}$ and $\tau\ge1$ let
\[
  \mathrm{ins}(Z,\tau)=\{z\in Z:z<\tau\}\cup\{\tau\}\cup\{z+1:z\in Z,\ z\ge\tau\}.
\]

\begin{proposition}[Insertion implies rows equal the top row]\label{prop:insertion}
Let $r>1$ and suppose that $\Phi_r(\mathrm{ins}(Z,\tau))\le\Phi_r(Z)$ for
every finite $Z\subset\N_{>0}$ and every $\tau\ge1$.  Then
$V_r(L,k)=\beta_r(L-k+1)$ for all $1\le k\le L$.
\end{proposition}

\begin{proof}
Let $S=\{s_1<\cdots<s_{k-1}\}$ be the exempted set of
Proposition~\ref{prop:rowvalue}(a), and let $Z_0\subset\N_{>0}$ have
$|Z_0|=n-1$.  Put $Z_j=\mathrm{ins}(Z_{j-1},s_j)$.  Each insertion adds
$s_j$ and moves only elements at least $s_j$, so it keeps
$s_1,\ldots,s_{j-1}$.  Thus $Z_{k-1}\supseteq S$,
$|Z_{k-1}|=L-1$, and $\Phi_r(Z_{k-1})\le\Phi_r(Z_0)$.
Proposition~\ref{prop:rowvalue}(a) gives $b_k(F)\ge1/\Phi_r(Z_0)$.

The infimum of $\Phi_r(Z_0)$ over such $Z_0$ is $\mu_r(\varnothing,n)$, by
Lemma~\ref{lem:intzeros} with $\omega(s)=r^{-s}$.  So
$b_k(F)\ge\beta_r(n)$, and Proposition~\ref{prop:rowvalue}(c) gives the
reverse inequality.
\end{proof}

The composite map $Z_0\mapsto Z_{k-1}$ keeps $S$ and sends the $z$-th
positive integer outside $\varnothing$ to the $z$-th one outside $S$.  On
one-hole sets the hypothesis is Lemma~\ref{lem:shift} together with
Lemma~\ref{lem:confdelr}, and Proposition~\ref{prop:oneholerows} is this
argument for $Z_0=[1,n]\setminus\{c\}$.

The hypothesis fails for every $r<2$, since
$\Phi_r(\{1\})=a\,\Phi_r(\varnothing)$.  It also fails at
$r=2,\frac94,\frac{23}{10}$.  In exact arithmetic
$\Phi_2(\mathrm{ins}(\{1,3,4,6\},1))=\frac{63}{50}\Phi_2(\{1,3,4,6\})$.
For $r=\frac94$ a set in $[1,19]$ with $\tau=1$ gives the ratio
$1.044$.  For $r=\frac{23}{10}$ the set $[1,23]\cup\{30,34,36,38,41,49,52,53,61\}$
with $\tau=24$ gives $1.018$.  In random tests with sets up to $140$ it held
at every $r\ge\frac52$ tried, with ratio below $0.98$; the maxima occur at
$\tau=\max Z$, where it is Lemma~\ref{lem:confdelr}.

The conjecture itself survived an exact search.  For
$r\in\{2,\frac{21}{10},\frac94\}$, $n\le6$ and all $S\subseteq[1,12]$ with
$|S|\le4$, a witness $Z\supseteq S$ with $|Z|=|S|+n-1$ had
$\Phi_r(Z)\le0.80\,\Phi_r(Z_0)$, where $Z_0$ is the best top-row set found.

\section{Hole shifts and insertions}\label{sec:holeshifts}

This section proves two of the steps left open above.  At several roots
all at least $\frac52$, the hole moves from $h$ to $h-1$
(Lemma~\ref{lem:holeshiftgen}), so Theorem~\ref{thm:rowsgen} becomes
$b_k\ge|f_D|\max_c1/T_c(R_n)$ (Theorem~\ref{thm:rowsgenfull}).  At a single
root, the insertion inequality of Proposition~\ref{prop:insertion} holds at
every hole $h$ with $\prod_{c\ge h}(1+1/c)\le r-1$, the product over the holes
(Lemma~\ref{lem:insertpi}).  Hence every row is at least $1/\Phi_r(Z_0)$ for
each $Z_0$ whose hole product is at most $r-1$ (Theorem~\ref{thm:rowspi}).
With an exact certificate for the top row (Proposition~\ref{prop:toprowcert}),
this proves that rows equal the top row for $r\ge3$ and $4\le n\le10$
(Corollary~\ref{cor:rowspiranges}).  Together with
Corollaries~\ref{cor:tworowsr} and~\ref{cor:rowstopranges}, this covers every
$n\le10$ at $r\ge3$.

\begin{lemma}[Moving the hole at several roots]\label{lem:holeshiftgen}
Let $R$ be a multiset of $m\ge2$ reals, let $\rho=\min R\ge2$, let $R'$ be $R$
with one copy of $\rho$ removed, and let $2\le h\le m$ and $t_0=m-h$.  If
$S_\rho(d+1,t_0)\le S_\rho(d,t_0)$ for $1\le d\le h-1$, then
\[
  T_h(R)\le T_{h-1}(R') .
\]
By Lemma~\ref{lem:shift} the hypothesis holds for every $h$ when
$\rho\ge\frac52$, and for $h=2$ when $\rho\ge2$.
\end{lemma}

\begin{proof}
Put $s=\rho-1\ge1$, $m'=m-1$, $j=h-1$, so $1\le j\le m'$, and write
$E_i=e_i(R')$ (zero for $i>m'$), $Q'=Q(R')$ and $p_i=|P_i(R')|$ for
$1\le i\le m'+1$.  By Lemma~\ref{lem:partialsums}, $P_i(R')$ has the sign
$(-1)^{i-1}$ for these $i$.

\emph{Reduction.}  We have $e_h(R)=E_h+\rho E_j$ and $Q(R)=sQ'$.  The
polynomial $P_h(R)$ is the coefficient of $u^{h-1}$ in
$(1-\rho u)\prod_{\rho'\in R'}(1-\rho'u)/(1-u)$, so
$P_h(R)=P_h(R')-\rho P_j(R')$ and $|P_h(R)|=p_h+\rho p_j$.  Also
$P_h(R')-P_j(R')=(-1)^jE_j$ gives $p_h+p_j=E_j$.  By
Proposition~\ref{prop:oneholegen},
\[
  T_h(R)=\frac{sQ'+p_h+\rho p_j}{sQ'(E_h+\rho E_j)},\qquad
  T_j(R')=\frac{Q'+p_j}{Q'E_j}.
\]
Clear the positive denominators and substitute $p_h=E_j-p_j$ and $\rho=s+1$.
The claim $T_h(R)\le T_j(R')$ becomes $\Delta\ge0$, where
\[
  \Delta=s\,(Q'+p_j)\,(E_{j+1}+sE_j)-E_j^2 .
\]
More precisely, $s(Q')^2E_je_h(R)\bigl(T_j(R')-T_h(R)\bigr)=Q'\Delta$.

\emph{Expansion about $\rho$.}  Write the elements of $R'$ as $\rho+y_i$ with
$y_i\ge0$, and put $g_l=e_l(y)\ge0$ and $M_l=m'-l$ for $0\le l\le m'$.  Write
$\rho^M$ for the multiset of $M$ copies of $\rho$.  The identities
$\prod_i(1+(\rho+y_i)u)=\sum_lg_lu^l(1+\rho u)^{M_l}$ and
$\prod_i(s+y_i)=\sum_lg_ls^{M_l}$ give
\[
  E_i=\sum_lg_l\,e_{i-l}(\rho^{M_l}),\qquad Q'=\sum_lg_l\,s^{M_l}.
\]
Next, $\prod_i(1-(\rho+y_i)u)/(1-u)=\sum_l(-1)^lg_lu^l(1-\rho u)^{M_l}/(1-u)$
gives
\[
  P_i(R')=\sum_{l<i}(-1)^lg_l\,P_{i-l}(\rho^{M_l}).
\]
Here $1\le i-l\le M_l+1$.  By Lemma~\ref{lem:partialsums} (for $M_l=0$,
$P_1(\varnothing)=1$), every term has the sign $(-1)^{i-1}$, so
$p_i=\sum_{l<i}g_l|P_{i-l}(\rho^{M_l})|$.  Hence
\[
  E_j=\sum_lg_l\epsilon_l,\qquad E_{j+1}+sE_j=\sum_lg_l\kappa_l,\qquad
  Q'+p_j=\sum_lg_l\alpha_l,
\]
with nonnegative coefficients
\[
  \epsilon_l=e_{j-l}(\rho^{M_l}),\qquad
  \kappa_l=e_{j+1-l}(\rho^{M_l})+s\,e_{j-l}(\rho^{M_l}),\qquad
  \alpha_l=s^{M_l}+|P_{j-l}(\rho^{M_l})| .
\]
In $\alpha_l$ the second term is absent when $l\ge j$.

\emph{The diagonal.}  We claim $\epsilon_l^2\le s\alpha_l\kappa_l$ for every
$l$.
\begin{itemize}
\item If $l>j$, then $\epsilon_l=0$.
\item If $l=j$, then $\epsilon_l=1$, $\alpha_l=s^{M_l}$ and $\kappa_l\ge s$,
so $s\alpha_l\kappa_l\ge s^{M_l+2}\ge1$.
\item If $l<j$, put $d=j-l\ge1$ and $M=M_l$, so $d\le M$.  Then
$s\alpha_l\kappa_l-\epsilon_l^2$ is the quantity $\Delta$ of the reduction for
the multiset $\rho^{M+1}$ (in place of $R$) and the index $d+1$ (in place of
$h$).  By the reduction it has the sign of
$T_d(\rho^M)-T_{d+1}(\rho^{M+1})=S_\rho(d,M-d)-S_\rho(d+1,M-d)$
(Proposition~\ref{prop:oneholegen}), and $M-d=m'-j=t_0$.  This is nonnegative
by hypothesis.
\end{itemize}

\emph{Conclusion.}  By the diagonal claim and the Cauchy--Schwarz inequality,
\[
  E_j=\sum_lg_l\epsilon_l\le\sum_lg_l\sqrt{s\alpha_l\kappa_l}
  \le\Bigl(s\sum_lg_l\alpha_l\Bigr)^{1/2}\Bigl(\sum_lg_l\kappa_l\Bigr)^{1/2},
\]
which is $E_j^2\le s(Q'+p_j)(E_{j+1}+sE_j)$, that is, $\Delta\ge0$.

For the last sentence of the lemma: Lemma~\ref{lem:shift} gives
$S_\rho(d+1,t_0)\le S_\rho(d,t_0)$ for all $d\ge1$ when $\rho\ge\frac52$, and
for $d=1$ when $\rho\ge2$.
\end{proof}

\begin{theorem}[Rows at several roots at least $\frac52$]\label{thm:rowsgenfull}
Let $\frac52\le r_1\le\cdots\le r_L$, let $F$ be a nonzero real multiple of
$\prod_{i=1}^L(x-r_i)$, $1\le k\le L$ and $n=L-k+1$, and let $R_n$ be the $n$
largest roots, as in Theorem~\ref{thm:rowsgen}.  Then
\[
  b_k(F)\ \ge\ |f_D|\max_{1\le c\le n}\frac1{T_c(R_n)} .
\]
If moreover $X(z,z')\le\max_c1/T_c(R_n)$ for all $1\le z<z'\le n$, with $X$
as in Proposition~\ref{prop:toprowgen} for $R=R_n$, then
$b_k(F)\ge|f_D|\,\beta(R_n)$.
\end{theorem}

\begin{proof}
Fix $c$ and $h$ with $c<h\le c+k-1$.  The smallest element of
$R_{h+n-c}$ is one of the $r_i$, so it is at least $\frac52$, and removing it
leaves $R_{h-1+n-c}$.  Lemma~\ref{lem:holeshiftgen} gives
$T_h(R_{h+n-c})\le T_{h-1}(R_{h-1+n-c})$.  Repeating this down to $h=c$ gives
$T_h(R_{h+n-c})\le T_c(R_n)$.  So the inner minimum in
Theorem~\ref{thm:rowsgen} is attained at $h=c$, which proves the first claim.
Under the extra hypothesis, Proposition~\ref{prop:toprowgen} gives
$\beta(R_n)=\max_c1/T_c(R_n)$.
\end{proof}

At equal roots the bound is the $B_r(n)$ of Theorem~\ref{thm:everyrowr}.  The
second part says that every row is at least the top row of
$\prod_{i\ge k}(x-r_i)$ whenever that top row is attained by one hole.

We now return to a single root $r$.  For a finite $Z\subset\N_{>0}$, its
\emph{holes} are the elements of $[1,\max Z]\setminus Z$.  Put
\[
  \Pi(Z)=\prod_{c\text{ a hole of }Z}\Bigl(1+\frac1c\Bigr),\qquad
  \Pi_h(Z)=\prod_{c\text{ a hole of }Z,\ c\ge h}\Bigl(1+\frac1c\Bigr).
\]
The empty product is $1$.

\begin{lemma}[Inserting at a hole]\label{lem:insertpi}
Let $r>1$, let $Z\subset\N_{>0}$ be finite, and let $h$ be a hole of $Z$.  Then
\[
  \Phi_r\bigl(\mathrm{ins}(Z,h)\bigr)\le\max\{1,\,a\,\Pi_h(Z)\}\,\Phi_r(Z).
\]
In particular $\Phi_r(\mathrm{ins}(Z,h))\le\Phi_r(Z)$ whenever
$\Pi_h(Z)\le r-1$.
\end{lemma}

\begin{proof}
Let $x=\max Z$, let $C$ be the set of holes, $m=\max C<x$ and $k=x-m\ge1$, so
$Z=[1,m+k]\setminus C$.  Write $C_<=\{c\in C:c<h\}$ and
$C_\ge=\{c\in C:c\ge h\}$.  Then $\mathrm{ins}(Z,h)=[1,m+1+k]\setminus C'$ with
\[
  C'=C_<\cup\{c+1:c\in C_\ge\}.
\]
Indeed the elements below $h$ are kept, $h$ is added, and every $z>h$ in
$Z$ moves to $z+1$.  So $\max C'=m+1$, and the number of filled positions
above the largest hole is again $k$.

Put $G=at$, $F=(1+at)/r$ and $u=(1-t)/r$ for $t\in[0,1]$.  Then $F-G=u\ge0$
and $G,F\le a$.  Since $r^{-c}(at)^{m-c}(1+at)^{c-1}=r^{-1}G^{m-c}F^{c-1}$ and
$r^{-c}t^{m-c}(1-t)^{c-1}=r^{-1}t^{m-c}u^{c-1}$, Lemma~\ref{lem:holeintr}
reads
\[
  \Phi_r(Z)=\frac1r\int_0^1t^k\Bigl(a^{k+1}\Omega_C(G,F)+\Theta_C(t)\Bigr)dt,
\]
where
\[
  \Omega_C(X,Y)=\sum_{c\in C}w_cX^{m-c}Y^{c-1},\qquad
  \Theta_C(t)=\sum_{c\in C}|w_c|\,t^{m-c}u^{c-1},
\]
and $w_c=\prod_{c'\in C}c'\big/\prod_{c'\in C\setminus\{c\}}(c-c')$.
The same formula holds for $\mathrm{ins}(Z,h)$, with $C'$, $m+1$ and the same
$k$.  We compare the two integrands pointwise.  Note that
$\prod_{c\in C'}c\big/\prod_{c\in C}c=\Pi_h(Z)$.

\emph{The tail term.}  For a finite set $A$ of reals let
$[A]f=\sum_{c\in A}f(c)/\prod_{c'\in A\setminus\{c\}}(c-c')$ be the divided
difference.  For $0<X\le Y$ put $\lambda=Y/X\ge1$.  Then
$\Omega_C(X,Y)=(\prod_{c\in C}c)\,X^mY^{-1}[C](\lambda^z)$.  By the
Hermite--Genocchi formula,
$[c_1,\ldots,c_\nu]f=\int_{\Delta}f^{(\nu-1)}(\theta_1c_1+\cdots+\theta_\nu c_\nu)\,d\theta$,
over the standard simplex $\Delta$.  For $f(z)=\lambda^z$ the integrand is
$(\log\lambda)^{\nu-1}\lambda^{\theta\cdot c}\ge0$.  List $C$ and $C'$
increasingly, so that $c'_i=c_i$ or $c_i+1$.  Then
$\theta\cdot c'\le\theta\cdot c+1$, and hence
$0\le[C'](\lambda^z)\le\lambda\,[C](\lambda^z)$.  Therefore
\[
  0\le\Omega_{C'}(X,Y)\le\Pi_h(Z)\,X\lambda\,\Omega_C(X,Y)=\Pi_h(Z)\,Y\,\Omega_C(X,Y).
\]
At $X=G$ and $Y=F$ (for $t>0$, and at $t=0$ by continuity) this gives
$\Omega_{C'}(G,F)\le a\,\Pi_h(Z)\,\Omega_C(G,F)$.

\emph{The hole term.}  Let $c\in C_<$.  Then $c$ is also a hole of $C'$, and
\[
  \frac{|w'_c|}{|w_c|}=\prod_{c'\in C_\ge}\frac{(c'+1)(c'-c)}{c'(c'+1-c)}\le1 .
\]
Its monomial $t^{m+1-c}u^{c-1}$ is $t$ times the old one.  Now let
$c\in C_\ge$.  Then
\[
  \frac{|w'_{c+1}|}{|w_c|}=\Pi_h(Z)\prod_{c'\in C_<}\frac{c-c'}{c+1-c'}\le\Pi_h(Z),
\]
and its monomial $t^{m-c}u^{c}$ is $u$ times the old one.  Since $t\le1$ and
$u\le1/r<a$, we get $\Theta_{C'}\le\max\{1,a\Pi_h(Z)\}\,\Theta_C$.

Both parts of the integrand are therefore multiplied by at most
$\max\{1,a\Pi_h(Z)\}$.
\end{proof}

For one hole, $\Pi_h=(h+1)/h$, and the lemma is Lemma~\ref{lem:shift}
together with the deletion of the top zero.  The lemma is useful when the
holes at and above $h$ are few or far out.  An insertion moves holes only
upward, and that is what makes it usable along the whole chain of
Proposition~\ref{prop:insertion}.

\begin{theorem}[Rows from zero sets with a small hole product]\label{thm:rowspi}
Let $r\ge2$ and $n\ge1$, and let $Z_0\subset\N_{>0}$ be finite with
$|Z_0|\le n-1$ and $\Pi(Z_0)\le r-1$.  Then every monic multiple $F$ of
$(x-r)^L$, $L\ge n$, satisfies
\[
  b_{L-n+1}(F)\ \ge\ \frac1{\Phi_r(Z_0)} .
\]
Consequently, if $\beta_r(n)=1/\Phi_r(Z_0)$ for such a $Z_0$, then
$V_r(L,L-n+1)=\beta_r(n)$ for every $L\ge n$.
\end{theorem}

\begin{proof}
Let $k=L-n+1$, and let $S=\{s_1<\cdots<s_{k-1}\}$ be the exempted set of
Proposition~\ref{prop:rowvalue}(a).  Build $Z_1,\ldots,Z_{k-1}$ as follows:
\begin{itemize}
\item if $s_j\in Z_{j-1}$, put $Z_j=Z_{j-1}$;
\item if $s_j\notin Z_{j-1}$ and $s_j<\max Z_{j-1}$, then $s_j$ is a hole;
put $Z_j=\mathrm{ins}(Z_{j-1},s_j)$;
\item otherwise put $Z_j=Z_{j-1}\cup\{s_j\}$.
\end{itemize}
Each step adds at most one element and moves only elements larger than
$s_j$.  So $Z_j\supseteq\{s_1,\ldots,s_j\}$ and $|Z_j|\le|Z_0|+j$.

Once a step of the third kind occurs, $s_j=\max Z_j$, and every later step is
of the third kind.  Each such step does not increase $\Phi_r$, by
Lemma~\ref{lem:confdelr} with $r\ge2$.

Before that, steps of the first two kinds keep the number of holes.  An
insertion at a hole $h$ keeps the holes below $h$ and moves those at or above
$h$ up by one.  So if $c_1<c_2<\cdots$ are the holes of $Z_0$ and
$c^{(j)}_1<c^{(j)}_2<\cdots$ those of $Z_j$, then $c^{(j)}_i\ge c_i$.
Hence at an insertion step
\[
  \Pi_{s_j}(Z_{j-1})\le\prod_i\Bigl(1+\frac1{c^{(j-1)}_i}\Bigr)
  \le\prod_i\Bigl(1+\frac1{c_i}\Bigr)=\Pi(Z_0)\le r-1,
\]
and Lemma~\ref{lem:insertpi} gives $\Phi_r(Z_j)\le\Phi_r(Z_{j-1})$.

Thus $Z_{k-1}\supseteq S$, $|Z_{k-1}|\le n-1+k-1=L-1$ and
$\Phi_r(Z_{k-1})\le\Phi_r(Z_0)$, and Proposition~\ref{prop:rowvalue}(a) gives
the bound.  The last claim follows with Proposition~\ref{prop:rowvalue}(c).
\end{proof}

To apply the theorem one needs the top row exactly.  The upper-bound
construction of Proposition~\ref{prop:toprowgen} works with any polynomial
multiplier.

\begin{proposition}[An exact certificate for the top row]\label{prop:toprowcert}
Let $r>1$, $n\ge2$, and let $Z_0\subset\N_{>0}$ have $|Z_0|=n-1$ and
$M=\max Z_0$.  Suppose that some real $p$ with $p(0)=1$ and $\deg p\le M+1-n$
has the following property.  Write $N(y)=p(y)(1-ry)^n=\sum_jN_jy^j$ and
$g_z=\sum_{j\le z}N_j$.  Suppose $|g_z|\le1/\Phi_r(Z_0)$ for
$1\le z\le M+1$.  Then $\beta_r(n)=1/\Phi_r(Z_0)$.
\end{proposition}

\begin{proof}
Proposition~\ref{prop:rowvalue}(a) with $k=1$ and $Z=Z_0$ gives
$\beta_r(n)\ge1/\Phi_r(Z_0)$.  For the other direction, the upper-bound part of
the proof of Proposition~\ref{prop:toprowgen} uses only $N(0)=1$ and the
divisibility of $N$ by $(1-ry)^n$.  It gives
$\beta_r(n)\le\sup_{z\ge1}|g_z|$.  Since $\deg N\le M+1$, $g_z=N(1)=g_{M+1}$ for
$z\ge M+1$, so the supremum is $\max_{1\le z\le M+1}|g_z|$.
\end{proof}

The polynomial $p$ is found from complementary slackness.  The $M+2-n$
equations $g_s=-\beta\,\sgn r_{Z_0}(s)$ for $s\in[1,M+1]\setminus Z_0$ are
linear in the coefficients of $p$ and in $\beta$.  Their solution has
$\beta=1/\Phi_r(Z_0)$, because
$\sum_{s\ge1}g_sr_{Z_0}(s)r^{-s}=-1$.  It then remains to check $|g_s|\le\beta$
on $Z_0$.  For example, at $r=2$ this proves the values
$\beta_2(3)=\frac{40}{11}$, $\beta_2(4)=\frac{672}{115}$ and
$\beta_2(5)=\frac{10368}{851}$ quoted after Corollary~\ref{cor:tworowsr}, with
$Z_0=\{2,5\}$, $\{1,3,7\}$ and $\{1,3,6,9\}$.

\begin{corollary}[Rows equal the top row, explicit ranges]\label{cor:rowspiranges}
For every $L\ge n$, $V_r(L,L-n+1)=\beta_r(n)$ in each of the following cases:
\begin{enumerate}
\item[(a)] $4\le n\le7$ and $r\ge\frac{23}8$;
\item[(b)] $8\le n\le10$ and $r\ge3$;
\item[(c)] $r=3$ and $n\le24$; $r=\frac72$ and $n\le23$; $r=4$ and $n\le20$; $r=5$ and $n\le14$.
\end{enumerate}
For instance, $\beta_3(4)=\frac{3240}{97}$, $\beta_3(5)=\frac{31104}{307}$,
$\beta_3(6)=\frac{1399680}{5387}$ and $\beta_3(7)=\frac{55112400}{67549}$,
attained at $Z_0=\{1,3,5\}$, $\{1,2,4,6\}$, $\{1,2,3,5,8\}$ and
$\{1,2,3,5,7,10\}$.
\end{corollary}

\begin{proof}
In each case we exhibit finitely many sets $Z_0$ with $|Z_0|=n-1$.  For
every $r$ in the range, one of them satisfies $\Pi(Z_0)\le r-1$ and the
hypothesis of Proposition~\ref{prop:toprowcert}.  Theorem~\ref{thm:rowspi}
then applies.

For (c) this is a finite exact computation at each point.  For (a) and (b)
the $p$ of Proposition~\ref{prop:toprowcert} is the solution of the linear
system above, so its coefficients are rational functions of $r$.  Every
condition becomes a polynomial inequality in $r$, together with the
nonvanishing of the determinant of the system.  The real roots of all these
polynomials above the threshold cut the half-line into intervals, and on each
interval one set is checked at a rational sample point.  At each root, one
uses a set that is valid on an adjacent interval and whose determinant does
not vanish there (checked by exact division by the minimal polynomial), and
continuity carries the inequalities over.  This is the same kind of exact
real-root computation as in Corollary~\ref{cor:rowstopranges}.
\end{proof}

Corollary~\ref{cor:rowspiranges} contains the cases $4\le n\le8$ of
Corollary~\ref{cor:rowstopranges} for $r\ge3$ (and for $r\ge\frac{23}8$ when
$n\le7$).  Below those thresholds of Corollary~\ref{cor:rowstopranges}, the
top-row optimum has two or more holes, and the corollary is new there.  The
limit $r\ge3$ is where the method stops, not where the rows stop being equal
to the top row.  Along the exactly certified optima, $\Pi(Z_0)$ behaves as
follows:
\begin{itemize}
\item at $r=3$ it rises from $1.6$ to about $1.99$ as $n$ goes from $5$ to
$24$;
\item at $r=\frac72$ it stays near $1.7$ for $n\le23$;
\item at $r=4$ it stays between $1.25$ and $1.6$ for $n\le20$;
\item at $r=\frac52$ it is between $1.86$ and $2.35$ for $n\le14$, far above
$\frac32$.
\end{itemize}
If, for some $r_0\ge3$, every $n$ has a top-row optimum with
$\Pi(Z_0)\le r-1$ for all $r\ge r_0$, then Theorem~\ref{thm:rowspi} proves
that rows equal the top row for every $r\ge r_0$.
Section~\ref{sec:rowsfinite} removes the condition on $\Pi(Z_0)$ at the
price of a finite computation (Theorem~\ref{thm:rowsfinite}).

\section{Finite certificates for the rows}\label{sec:rowsfinite}

Theorem~\ref{thm:rowspi} needs $\Pi(Z_0)\le r-1$ for a top-row optimum
$Z_0$.  Along the certified optima this fails for $r<3$ once $n$ is
moderately large, and at $r=3$ it fails for most $n\ge25$.  This section removes the hypothesis
at the price of a finite computation.  The insertion chain of
Theorem~\ref{thm:rowspi} only ever produces sets obtained from $Z_0$ by
shifting its holes upward (Lemma~\ref{lem:holestates}).  Once the holes at and
above the insertion point have product at most $r-1$, every further step is
free by Lemma~\ref{lem:insertpi}.  So for $r>2$ only finitely many shifted
sets have to be compared with $Z_0$ (Theorem~\ref{thm:rowsfinite}), and
the comparison is monotone in $r$ (Lemma~\ref{lem:rmono}), so it can be
carried out on whole intervals of roots.  For a fixed number of exempted
positions a second reduction works at every $r>1$, including $r=2$ and
$1<r<2$.  Exempted positions far beyond a witness can simply be appended
(Lemma~\ref{lem:farappend}), and the remaining positions form a finite tree
(Theorem~\ref{thm:rowstree}).

\paragraph{Hole shifts.}
Let $Z_0\subset\N_{>0}$ be finite with holes $c_1<\cdots<c_\nu$ ($\nu\ge0$)
and $M=\max Z_0$.  For a nondecreasing sequence $E=(e_1,\ldots,e_\nu)$ of
nonnegative integers put
\[
  Z_0^E=[1,M+e_\nu]\setminus\{c_1+e_1,\ldots,c_\nu+e_\nu\}
\]
(for $\nu=0$ the only sequence is empty and $Z_0^E=Z_0$).  Since
$c_i+e_i<c_{i+1}+e_{i+1}$ and $c_\nu+e_\nu<M+e_\nu$, the holes of $Z_0^E$
are exactly the $c_i+e_i$, its maximum is $M+e_\nu$, and
$|Z_0^E|=|Z_0|+e_\nu$.  Write $\mathbf 1_{\ge i}$ for the sequence with ones
in the positions $i,\ldots,\nu$ and zeros before.

\begin{lemma}[The chain moves holes]\label{lem:holestates}
For $1\le i\le\nu$ and $h=c_i+e_i$,
\[
  \mathrm{ins}(Z_0^E,h)=Z_0^{E+\mathbf 1_{\ge i}},\qquad
  \Pi_h(Z_0^E)=\prod_{l=i}^{\nu}\Bigl(1+\frac1{c_l+e_l}\Bigr).
\]
\end{lemma}

\begin{proof}
The insertion keeps the elements below $h$, adds $h$, and moves every element
above $h$ up by one; so the holes $c_l+e_l$ with $l<i$ stay, and the top
moves from $M+e_\nu$ to $M+e_\nu+1$.  An integer $p>h+1$ lies in the new
set exactly when $p-1$ lies in $Z_0^E$, so the new holes above $h+1$ are
the $c_l+e_l+1$ with $l>i$.  The integer $h+1$ lies in the new set exactly
when $h\in Z_0^E$, which is false, so $h+1=c_i+e_i+1$ is a hole.  The second
formula is the definition of $\Pi_h$, the holes at or above $h$ being the
$c_l+e_l$ with $l\ge i$.
\end{proof}

For $r>2$ and $1\le i\le\nu$ put
\[
  \pi_i(x)=\prod_{l=i}^{\nu}\Bigl(1+\frac1{c_l+x}\Bigr),\qquad
  x_i^*=\min\{x\in\N:\ \pi_i(x)\le r-1\}.
\]
These are finite because $\pi_i(x)\to1<r-1$, and $x_1^*\ge x_2^*\ge\cdots$
because $\pi_{i+1}\le\pi_i$.  Let $F_r(Z_0)$ be the set of nondecreasing
$E$ with
\[
  e_i\le\max(e_{i-1},x_i^*)\qquad(1\le i\le\nu,\ e_0=0).
\]
By induction $e_i\le x_1^*$, so $F_r(Z_0)$ is finite; it is $\{0\}$ when
$\Pi(Z_0)=\pi_1(0)\le r-1$.

\begin{theorem}[Rows from finitely many shifted sets]\label{thm:rowsfinite}
Let $r>2$, $n\ge1$, and let $Z_0\subset\N_{>0}$ be finite with
$|Z_0|\le n-1$ and
\[
  \Phi_r(Z_0^E)\le\Phi_r(Z_0)\qquad\text{for every }E\in F_r(Z_0).
\]
Then every monic multiple $F$ of $(x-r)^L$, $L\ge n$, satisfies
$b_{L-n+1}(F)\ge1/\Phi_r(Z_0)$.  If moreover $\beta_r(n)=1/\Phi_r(Z_0)$,
then $V_r(L,L-n+1)=\beta_r(n)$ for every $L\ge n$.
\end{theorem}

\begin{proof}
\emph{The chain.}  Let $k=L-n+1$, let $S=\{s_1<\cdots<s_{k-1}\}$ be the
exempted set of Proposition~\ref{prop:rowvalue}(a), and build
$Z_1,\ldots,Z_{k-1}$ from $Z_0$ by the three rules of the proof of
Theorem~\ref{thm:rowspi}.  Up to the first step of the third kind every
$Z_j$ has the form $Z_0^{D}$: a step of the first kind changes nothing,
and a step of the second kind inserts at a hole $c_i+d_i$ of $Z_0^D$, which
gives $Z_0^{D+\mathbf 1_{\ge i}}$ by Lemma~\ref{lem:holestates}.  After
that, every step adds an element above the current maximum, which does not
increase $\Phi_r$ by Lemma~\ref{lem:confdelr}, as $r\ge2$.  So
$Z_{k-1}\supseteq S$, $|Z_{k-1}|\le L-1$, and
$\Phi_r(Z_{k-1})\le\Phi_r(Z_0^D)$ for some nondecreasing $D\ge0$.  By
Proposition~\ref{prop:rowvalue}(a) it suffices to show
$\Phi_r(Z_0^D)\le\Phi_r(Z_0)$ for every such $D$.

\emph{A path to $D$.}  Put $d_0=0$ and, for $1\le i\le\nu$ and
$d_{i-1}\le x\le d_i$,
\[
  E^{(i,x)}=(d_1,\ldots,d_{i-1},x,\ldots,x).
\]
Then $E^{(1,0)}=0$, $E^{(i,d_i)}=E^{(i+1,d_i)}$, $E^{(\nu,d_\nu)}=D$, and
$E^{(i,x+1)}=E^{(i,x)}+\mathbf 1_{\ge i}$.  Walking through
$i=1,\ldots,\nu$ and $x=d_{i-1},\ldots,d_i-1$ gives a sequence of steps from
$0$ to $D$, and by Lemma~\ref{lem:holestates} the step $(i,x)$ is the
insertion at the hole $c_i+x$ of $Z_0^{E^{(i,x)}}$, whose hole product at
and above that hole is $\pi_i(x)$.  Call the step \emph{free} if
$x\ge x_i^*$.  Then $a\pi_i(x)\le1$, and Lemma~\ref{lem:insertpi} shows that
a free step does not increase $\Phi_r$.

\emph{After a free step all steps are free.}  Let $(i,x)$ be free.  The
later steps of stage $i$ have larger $x$.  A step $(i',x')$ with $i'>i$ has
$x'\ge d_{i'-1}\ge d_i\ge x+1>x_i^*\ge x_{i'}^*$.

\emph{The stopping state.}  Let $E$ be the state just before the first free
step, or $E=D$ if no step is free.  Then $\Phi_r(Z_0^D)\le\Phi_r(Z_0^E)$,
and we show $E\in F_r(Z_0)$, which finishes the proof.  For every stage $l$
that contains only steps that are not free, either $d_l=d_{l-1}$ or the last
step $(l,d_l-1)$ has $d_l-1<x_l^*$; in both cases $d_l\le\max(d_{l-1},x_l^*)$.
If no step is free this applies to every $l$ and $E=D\in F_r(Z_0)$.
Otherwise let the first free step be $(i,x)$.  All stages $l<i$ contain only
steps that are not free, and the steps $(i,x')$ with $d_{i-1}\le x'<x$
are not free, so $x=\max(d_{i-1},x_i^*)$.  So
$E=(d_1,\ldots,d_{i-1},x,\ldots,x)$ satisfies the defining inequalities at
the positions $l<i$, at $i$ with equality, and trivially at the positions
$l>i$.

The last claim follows with Proposition~\ref{prop:rowvalue}(c).
\end{proof}

When $\Pi(Z_0)\le r-1$ the set $F_r(Z_0)$ is $\{0\}$ and this is
Theorem~\ref{thm:rowspi} (for $r>2$).  The next lemma turns a check at finitely
many roots into a check on intervals.

\begin{lemma}[Monotonicity in the root]\label{lem:rmono}
For every finite $Z$, $r\mapsto\Phi_r(Z)$ is decreasing on $(1,\infty)$.
For fixed $Z_0$ each $x_i^*$ is nonincreasing in $r>2$, so
$F_r(Z_0)\subseteq F_{r'}(Z_0)$ for $r\ge r'>2$.  Consequently, if
$2<A\le B$ and
\[
  \Phi_A(Z_0^E)\le\Phi_B(Z_0)\qquad\text{for every }E\in F_A(Z_0),
\]
then the hypothesis of Theorem~\ref{thm:rowsfinite} holds for every
$r\in[A,B]$.
\end{lemma}

\begin{proof}
$\Phi_r(Z)=\sum_{s\ge1}|r_Z(s)|r^{-s}$ with $r_Z$ independent of $r$ and
not identically zero on $\N_{>0}$.  The thresholds $x_i^*$ are defined by
$\pi_i(x)\le r-1$ with $\pi_i$ independent of $r$, and the defining
inequalities of $F_r(Z_0)$ weaken as the $x_i^*$ grow.  For $r\in[A,B]$ and
$E\in F_r(Z_0)\subseteq F_A(Z_0)$,
$\Phi_r(Z_0^E)\le\Phi_A(Z_0^E)\le\Phi_B(Z_0)\le\Phi_r(Z_0)$.
\end{proof}

\paragraph{Cutting the chain.}
Near $r=2$ the shifted sets $Z_0^E$ with $E=(x,\ldots,x)$ can exceed
$\Phi_r(Z_0)$ (see the remarks after
Corollary~\ref{cor:rowsfiniteranges}).  Such a state is reached by
inserting a block of exempted positions at the lowest hole, and there the
chain may re-optimize the zeros above the block.  To state this, let $W$ be
finite and $t\ge0$, and let $c_1<\cdots<c_\nu$ be the holes of $W$ above
$t$.  Define $W^E$, $\pi_i$ and $x_i^*$ as above with these holes, the holes
of $W$ at most $t$ staying in place, and call the sets
$E^{(i,x)}=(e_1,\ldots,e_{i-1},x,\ldots,x)$, $x\ge e_{i-1}$, the
\emph{states}; the \emph{step} $(i,x)$ leads from $E^{(i,x)}$ to
$E^{(i,x+1)}$, and it is free if $x\ge x_i^*$.  Put
$\lambda_i(x)=\max\{1,a\pi_i(x)\}$ and define $\Lambda_i(x)$ for
$1\le i\le\nu+1$ and $x\ge0$ by $\Lambda_{\nu+1}\equiv1$, $\Lambda_i(x)=1$
for $x\ge x_i^*$, and
\[
  \Lambda_i(x)=\max\{\Lambda_{i+1}(x),\ \lambda_i(x)\Lambda_i(x+1)\}
  \qquad(x<x_i^*).
\]
A state $E'$ lies \emph{beyond} the state $E=E^{(i,x)}$ if $e'_l=e_l$ for
$l<i$ and $e'_i\ge x$; these are the states whose path, in the stage order
of the proof of Theorem~\ref{thm:rowsfinite}, passes through $E$.

\begin{lemma}[Loss along a path]\label{lem:lossbound}
Let $r>2$.  If $E'$ lies beyond $E=E^{(i,x)}$, then
$\Phi_r(W^{E'})\le\Lambda_i(x)\,\Phi_r(W^E)$.
\end{lemma}

\begin{proof}
The path from $E$ to $E'$ is a sequence of steps and stage advances.  By
Lemmas~\ref{lem:holestates} and~\ref{lem:insertpi}, a step $(i',x')$
multiplies $\Phi_r$ by at most $\lambda_{i'}(x')$, which is $1$ when the
step is free.  After a free step all later steps are free, as in the proof
of Theorem~\ref{thm:rowsfinite}.  So the product of the factors along the
path is at most the largest such product over all continuations from
$(i,x)$, which is $\Lambda_i(x)$ by its recursion.
\end{proof}

Call $(W,t)$ \emph{settled} if for every finite $S'\subset(t,\infty)$ some
$W'\supseteq(W\cap[1,t])\cup S'$ has $|W'|\le|W|+|S'|$ and
$\Phi_r(W')\le\Phi_r(Z_0)$.

\begin{theorem}[Recursive certificate]\label{thm:rowsrec}
Let $r>2$, $n\ge1$, let $Z_0\subset\N_{>0}$ be finite with $|Z_0|\le n-1$,
and let $\mathcal P$ be a finite set of pairs $(W,t)$ containing $(Z_0,0)$.
Suppose that every $(W,t)\in\mathcal P$ has $\Phi_r(W)\le\Phi_r(Z_0)$ and
comes with two sets $\mathcal C$ (cuts) and $\mathcal G$ (secured states)
of states, each reached by a step that is not free, such that
\begin{itemize}
\item every $E=E^{(i,x)}\in\mathcal G$ has
$\Lambda_i(x)\Phi_r(W^E)\le\Phi_r(Z_0)$;
\item for every $E\in\mathcal C$, reached by the step $(i,x-1)$, and
$t'=c_i+x-1$, some $(W',t')\in\mathcal P$ has
$W'\cap[1,t']=W^E\cap[1,t']$ and $|W'|=|W^E|$;
\item $\Phi_r(W^E)\le\Phi_r(Z_0)$ for every state $E$ reached by a path of
steps that are not free, and lying beyond no element of
$\mathcal C\cup\mathcal G$.
\end{itemize}
Then $(Z_0,0)$ is settled, and every monic multiple $F$ of $(x-r)^L$,
$L\ge n$, satisfies $b_{L-n+1}(F)\ge1/\Phi_r(Z_0)$.
\end{theorem}

\begin{proof}
Since $t'=c_i+x-1>t$, we may show that every $(W,t)\in\mathcal P$ is
settled by induction downward in $t$.  Let $S'\subset(t,\infty)$ be finite,
and run the chain of the proof of Theorem~\ref{thm:rowsfinite} from $W$ on
the elements of $S'$ in increasing order.  By Lemma~\ref{lem:holestates},
which applies verbatim to the holes above $t$, up to the first step of the
third kind the chain visits sets $W^E$.  Two consecutive insertions happen
at increasing positions, so the insertions occur in stage order: the hole
index never decreases and, within a stage, the level increases.  Let $D$ be
the final shift, and follow the path to $D$.

If the path passes through some $E\in\mathcal C$, reached by the step
$(i,x-1)$, then right after the insertion at $t'=c_i+x-1$ the chain holds
$W^E$, and every element of $S'$ processed so far lies in $W^E\cap[1,t']$.
Take $(W',t')\in\mathcal P$ as in the hypothesis.  It is settled by
induction, so $S''=S'\cap(t',\infty)$ yields
$W''\supseteq(W^E\cap[1,t'])\cup S''\supseteq(W\cap[1,t])\cup S'$ with
$\Phi_r(W'')\le\Phi_r(Z_0)$ and $|W''|\le|W^E|+|S''|\le|W|+|S'|$, since
$|W^E|$ exceeds $|W|$ by the number of insertions so far.  Here
$W^E\cap[1,t]=W\cap[1,t]$ because only holes above $t$ move.

Otherwise, if the first state on the path that lies in
$\mathcal C\cup\mathcal G$ exists, it lies in $\mathcal G$, and
Lemma~\ref{lem:lossbound} gives $\Phi_r(W^D)\le\Phi_r(Z_0)$.  If there is
no such state, then the stopping state of the proof of
Theorem~\ref{thm:rowsfinite} is reached by a path of steps that are not free
and lies beyond no element of $\mathcal C\cup\mathcal G$, so
$\Phi_r(W^D)\le\Phi_r(W^{E_{\rm stop}})\le\Phi_r(Z_0)$.  In both cases the
elements appended afterwards do not increase $\Phi_r$
(Lemma~\ref{lem:confdelr}), and the final set contains
$(W\cap[1,t])\cup S'$ and has at most $|W|+|S'|$ elements.

For $(Z_0,0)$ and $S'=S$, the exempted set of
Proposition~\ref{prop:rowvalue}(a), this gives a witness with at most
$n-1+|S|=L-1$ elements.
\end{proof}

Theorem~\ref{thm:rowsfinite} is the case $\mathcal P=\{(Z_0,0)\}$,
$\mathcal C=\mathcal G=\varnothing$: the states reached by paths of steps
that are not free are exactly the elements of $F_r(Z_0)$.  The secured
states only shorten the computation.

\begin{corollary}[Rows equal the top row, finite certificates]\label{cor:rowsfiniteranges}
For every $L\ge n$, $V_r(L,L-n+1)=\beta_r(n)$ in each of the following cases:
\begin{enumerate}
\item[(a)] $n=3$ and $r\ge\frac{103}{50}$; $n=4$ and $r\ge\frac{209}{100}$;
$n=5$ and $r\ge\frac{213}{100}$; $n=6$ and $r\ge\frac{43}{20}$; $n=7$ and $r\ge\frac{11}5$;
\item[(b)] $r=\frac{11}5$ and $n\le10$; $r=\frac94$ and $n\le13$;
$r=\frac{23}{10}$ and $n\le15$; $r=\frac{12}5$ and $n\le18$; $r=\frac52$ and
$n\le25$; $r=\frac{11}4$ and $n\le18$; $r=3$ and $n\le45$; $r=\frac72$ and
$n\le 45$; $r=4$ and $n\le 40$;
\item[(c)] $r=\frac{21}{10}$ and $n\le7$; $r=\frac{11}5$ and $11\le n\le 14$.
\end{enumerate}
\end{corollary}

\begin{proof}
For $n\le2$ this is Corollary~\ref{cor:tworowsr}.  In each remaining case
we exhibit finitely many sets $Z_0$ with $|Z_0|=n-1$ and check, for every
$r$ in the range, that one of them satisfies the hypothesis of
Proposition~\ref{prop:toprowcert} and the hypothesis of
Theorem~\ref{thm:rowsfinite}; in case (c) the hypothesis of
Theorem~\ref{thm:rowsrec} replaces the latter, with $\mathcal P$ consisting of
$(Z_0,0)$ and at most one pair obtained by a cut at the lowest hole.  For (b) and
(c) these are finite exact computations at a point.  For (a) the first hypothesis is checked on the whole
half-line exactly as in the proof of Corollary~\ref{cor:rowspiranges}
(without the condition on $\Pi$).  This cuts the half-line at finitely many
algebraic points into pieces, each with a certified $Z_0$.  On each piece,
enclosed in a rational interval $[A,B]$, the second hypothesis is checked by
Lemma~\ref{lem:rmono}, after bisecting $[A,B]$ where needed; on the last,
unbounded piece it is automatic once $r-1\ge\Pi(Z_0)$.  At a cut point one
uses a set that is certified and checked on an adjacent piece and whose
determinant does not vanish at the point.  All comparisons of $\Phi$ are
exact, with $\Phi_r(Z)$ written as a rational number: the finitely many
terms $s\le\max Z$ are summed directly, and the rest by the finite-difference
formula $\sum_{s>M}r_Z(s)y^s=y^{M+1}\sum_j\Delta^jr_Z(M+1)\,y^j/(1-y)^{j+1}$,
$y=1/r$, using that $r_Z$ has one sign beyond $M=\max Z$.
\end{proof}

\paragraph{What the computations show.}
The largest ratios $\Phi_r(Z_0^E)/\Phi_r(Z_0)$ over $F_r(Z_0)$ occur mostly
at uniform shifts $E=(x,\ldots,x)$, that is, for an exempted block starting
at the lowest hole of $Z_0$, and every ratio above $1$ that we met occurs
there.  At $r=3$ the ratio is
below $0.41$ for $n\le45$, and at $r=\frac52$ it is below $0.64$ for
$n\le25$.  It rises as $r$ decreases: it is $0.71$ at $r=\frac94$, $0.78$
at $r=\frac{11}5$ and $n=10$, and $1.03$ at $r=\frac{11}5$ and $n=11$.  At
$r=\frac{21}{10}$ it is $1.12$ and $1.26$ for $n=5,6$, and at
$r=\frac{41}{20}$ it is already $1.03$ for $n=3$.  At $r=\frac{21}{10}$
($5\le n\le7$) and at $r=\frac{11}5$ ($11\le n\le14$),
Theorem~\ref{thm:rowsrec} succeeds with at most one cut, at the lowest hole;
this gives (c).  After the cut the re-optimized witness has much smaller
$\Phi_r$, and the secured states keep the search to at most a few million
states.

The hole product of the certified optima grows slowly.  At $r=3$ it
exceeds $r-1=2$ for $n=25,26,27$ and $31\le n\le45$ (it is $2.06$ at
$n=45$), so Theorem~\ref{thm:rowspi} does not apply there.  Even so,
$|F_3(Z_0)|\le3$ for these $n$.  At $r=\frac72$ and $r=4$ the hole
product stays below $1.80$ ($n\le45$) and $1.64$ ($n\le40$), well below
$r-1$, so there Theorem~\ref{thm:rowspi} applies directly; if this persists
for all $n$, the rows equal the top row for every $n$ at these roots.  At $r=\frac52$, $\Pi(Z_0)\approx2.4$
against $r-1=\frac32$ for $19\le n\le25$, and $|F_{5/2}(Z_0)|$ grows to
about $2\cdot10^5$.

Near $r=2$ translating the holes is the wrong move.  An exempted block
$[c_1,c_1+x-1]$ acts like the prefix weight $\binom{s-1}{x}r^{-s}$ of
Theorem~\ref{thm:prefixrows}, under which good zero sets are dilated rather
than translated.  At $r=2$ every $x_i^*$ is infinite, and the chain fails
outright.  For $n=2$, $Z_0=\{2\}$ has $\Phi_2(Z_0)=\frac12$, but
Lemma~\ref{lem:holeintr} gives
\[
  \Phi_2\bigl(Z_0^{(x)}\bigr)=\Phi_2\bigl([1,x]\cup\{x+2\}\bigr)
  =\frac{x+2^{-x}}{x+2}\ \longrightarrow\ 1,
\]
although $V_2(L,L-1)=2$ by Corollary~\ref{cor:tworowsr}.  At $r=2$ the
certificates below therefore treat one row at a time.

\paragraph{Far exempted positions.}
For a fixed number of exempted positions there is a second reduction, which
does not use Lemma~\ref{lem:insertpi} and so works for every $r>1$.

\begin{lemma}[Appending far zeros]\label{lem:farappend}
Let $r>1$, let $W\subset\N_{>0}$ be finite and nonempty with $d=|W|$ and
$x=\max W$, let $u\ge1$, and let $D>d/(r-1)$ be an integer.  Put
$\rho=(1+d/D)/r<1$ and $s_0=x+D$.  Suppose that $\tau\ge s_0$ satisfies
$y:=\rho(1+1/\tau)^u<1$ and
\begin{equation}\tag{FA}\label{eq:farappend}
  \frac y{1-y}\ \le\ \sum_{j=0}^{\tau-s_0}\rho^{-j}\Bigl(1-\frac j\tau\Bigr).
\end{equation}
Then \eqref{eq:farappend} holds for every larger $\tau$, and
$\Phi_r(W\cup U)\le\Phi_r(W)$ for every finite $U$ with $\min U\ge\tau$ and
$|U|\le u$.
\end{lemma}

\begin{proof}
As $\tau$ grows the left side of \eqref{eq:farappend} decreases, and the right
side gains a nonnegative term while each $1-j/\tau$ grows.

Put $w(s)=|r_W(s)|r^{-s}$ and $\pi(s)=\prod_{v\in U}(1-s/v)$, so that
$\Phi_r(W\cup U)=\sum_sw(s)|\pi(s)|$.  Let $\tau'=\min U\ge\tau$.  For
$s\le\tau'$ every factor of $\pi(s)$ lies in $[0,1]$ and one of them is
$1-s/\tau'$, so $|\pi(s)|\le1-s/\tau'$.  For $s>\tau'$ every factor has
absolute value at most $\max(1,s/v-1)\le s/\tau'$, so
$|\pi(s)|\le(s/\tau')^u$.  Hence
\[
  \Phi_r(W\cup U)-\Phi_r(W)\le-\sum_{s\le\tau'}w(s)\frac s{\tau'}
     +\sum_{s>\tau'}w(s)\Bigl(\frac s{\tau'}\Bigr)^u .
\]
For $s>x$,
\[
  \frac{w(s+1)}{w(s)}=\frac1r\prod_{z\in W}\Bigl(1+\frac1{s-z}\Bigr)
  \le\frac1r\prod_{i=1}^{d}\Bigl(1+\frac1{s-x-1+i}\Bigr)
  =\frac1r\cdot\frac{s-x+d}{s-x},
\]
because $1+1/(s-z)$ increases with $z<s$ and the $i$-th largest element of
$W$ is at most $x-i+1$.  So $w(s+1)\le\rho\,w(s)$ for $s\ge s_0$, and
$w(\tau')>0$.  With $(1+j/\tau')^u\le(1+1/\tau')^{uj}\le(1+1/\tau)^{uj}$,
\[
  \sum_{s>\tau'}w(s)\Bigl(\frac s{\tau'}\Bigr)^u\le w(\tau')\sum_{j\ge1}y^j
  =w(\tau')\frac y{1-y},\qquad
  \sum_{s\le\tau'}w(s)\frac s{\tau'}\ge w(\tau')
     \sum_{j=0}^{\tau'-s_0}\rho^{-j}\Bigl(1-\frac j{\tau'}\Bigr),
\]
the second because $w(\tau'-j)\ge\rho^{-j}w(\tau')$ for $\tau'-j\ge s_0$.
By \eqref{eq:farappend} at $\tau'$ the difference is at most $0$.
\end{proof}

\begin{theorem}[A finite tree for a fixed row]\label{thm:rowstree}
Let $r>1$, $n,k\ge1$ and $\varphi>0$.  Let $\mathcal T$ be a family of finite
sets $T\subset\N_{>0}$ with $|T|\le k-1$ and $\varnothing\in\mathcal T$.
Suppose that every $T\in\mathcal T$ carries a set $W_T\supseteq T$ with
$|W_T|\le n-1+|T|$ and $\Phi_r(W_T)\le\varphi$, and, if $|T|<k-1$, an
integer $\theta_T$ such that
\begin{itemize}
\item Lemma~\ref{lem:farappend} applies to $W_T$, $u=k-1-|T|$ and
$\tau=\theta_T$, and
\item $T\cup\{\sigma\}\in\mathcal T$ for every integer $\sigma$ with
$\max T<\sigma<\theta_T$ ($\max\varnothing=0$).
\end{itemize}
Then every monic multiple $F$ of $(x-r)^{n+k-1}$ has $b_k(F)\ge1/\varphi$,
that is, $V_r(n+k-1,k)\ge1/\varphi$.
\end{theorem}

\begin{proof}
Let $S=\{s_1<\cdots<s_{k-1}\}$ be the exempted set of
Proposition~\ref{prop:rowvalue}(a), $T_j=\{s_1,\ldots,s_j\}$, and let $j$
be largest with $T_0,\ldots,T_j\in\mathcal T$.  If $j=k-1$, then $W_{T_j}$
contains $S$ and has at most $n+k-2$ elements.  Otherwise
$s_{j+1}\ge\theta_{T_j}$, since else $T_{j+1}\in\mathcal T$.  Then
$U=\{s_{j+1},\ldots,s_{k-1}\}$ has $\min U\ge\theta_{T_j}>\max W_{T_j}$ and
$|U|=k-1-j$, so $W_{T_j}\cup U\supseteq S$ has at most $n+k-2$ elements
and, by Lemma~\ref{lem:farappend},
$\Phi_r(W_{T_j}\cup U)\le\Phi_r(W_{T_j})\le\varphi$.  In both cases
Proposition~\ref{prop:rowvalue}(a) gives $b_k(F)\ge1/\varphi$.
\end{proof}

The tree is finite: $\theta_T$ depends only on $W_T$ and $u$, and each chain
$\varnothing\subset T_1\subset\cdots$ in $\mathcal T$ has at most $k$
members.  The natural target is $\varphi=\max_{i\le k}\nu_i(n)$, the value
conjectured after Corollary~\ref{cor:belowtwo}; the matching upper bound
needs the exact value of one $\nu_i(n)$, which the construction of
Proposition~\ref{prop:toprowcert} gives with exempted positions.

\begin{proposition}[An exact certificate with exempted positions]\label{prop:toprowcertS}
Let $r>1$, $L\ge2$, and let $S\subseteq Z\subset\N_{>0}$ with $|Z|=L-1$ and
$M=\max Z$.  Suppose that some real $p$ with $p(0)=1$ and
$\deg p\le M+1-L$ has the following property: with
$N(y)=p(y)(1-ry)^L=\sum_jN_jy^j$ and $g_z=\sum_{j\le z}N_j$, we have
$|g_z|\le1/\Phi_r(Z)$ for every $z\in[1,M+1]\setminus S$.  Then
$\mu_r(S,L)=\Phi_r(Z)$ and $V_r(L,|S|+1)\le1/\Phi_r(Z)$.
\end{proposition}

\begin{proof}
The construction in the proof of Proposition~\ref{prop:toprowgen} uses only
$N(0)=1$ and the divisibility of $N$ by $(1-ry)^L$.  For each $K\ge M+1$ it
gives a monic multiple of $(x-r)^L$ whose nonleading coefficient at backward
distance $s$ is $g_s$ for $s\le K$ (with $g_s=g_{M+1}$ for $s>M+1$), and
at most $\beta+C/K$ in absolute value beyond, where
$\beta=\max_{z\notin S}|g_z|\le1/\Phi_r(Z)$.  The $|S|$ positions at the
distances in $S$ can carry at most the $|S|$ largest magnitudes, so
$b_{|S|+1}\le\beta+C/K$ for this multiple, and $V_r(L,|S|+1)\le1/\Phi_r(Z)$.
For any real $q$ with $\deg q\le L-1$, $q(0)=1$ and $q|_S=0$, the identity
of Proposition~\ref{prop:rowvalue}(a) for this multiple gives
$1\le(\beta+C/K)\Phi_r(q)$, so $\mu_r(S,L)\ge1/\beta\ge\Phi_r(Z)$; and
$\mu_r(S,L)\le\Phi_r(r_Z)=\Phi_r(Z)$.
\end{proof}

As for Proposition~\ref{prop:toprowcert}, $p$ is found from the
complementary slackness equations $g_s=-\beta\sgn r_Z(s)$ for
$s\in[1,M+1]\setminus Z$, whose solution has $\beta=1/\Phi_r(Z)$; the
positions in $S$ are not constrained.  For $S=\varnothing$ this is
Proposition~\ref{prop:toprowcert}.

\begin{corollary}[Single rows, including $r\le2$]\label{cor:rowstreevals}
Write $n=L-k+1$.  The conjectured value
$V_r(L,k)=\min_{1\le i\le k}1/\nu_i(n)$ holds in each of the following
cases, and the minimum is attained at the index $i^*$ stated.
\begin{enumerate}
\item[(a)] $r=2$ with $k=2$, $n\le14$; $k=3$, $n\le12$; $k=4$, $n\le8$;
$k=5$, $n\le5$; $k=6$, $n\le4$.  Here $i^*=1$, so
$V_2(L,k)=\beta_2(L-k+1)$; for instance every monic multiple of $(x-2)^4$
has $b_2\ge\frac{40}{11}$, and every monic multiple of $(x-2)^5$ has
$b_2\ge\frac{672}{115}$ and $b_3\ge\frac{40}{11}$.
\item[(b)] $r\in\{\frac{41}{20},\frac{21}{10}\}$ with $k=2$, $n\le12$ and
$k=3$, $n\le9$; here $i^*=1$.
\item[(c)] $r\in\{\frac{19}{10},\frac74,\frac32\}$ with $k=2$, $n\le12$;
$k=3$, $n\le9$; $k=4$, $n\le6$.  Here $i^*=1$, except $i^*=k$ for $n=2$ at
$r=\frac74$, and at $r=\frac32$ for $n\le4$ ($k=2$) and $n\le5$ ($k=3,4$).
\item[(d)] $r=\frac54$ with $k=2$, $n\le12$ ($i^*=2$ for $n\le10$ and
$i^*=1$ for $n=11,12$); $k=3$, $n\le9$ ($i^*=3$); $k=4$, $n\le4$ ($i^*=4$).
\item[(e)] $r=\frac{11}{10}$ with $k=2$, $n\le12$ ($i^*=2$); $k=3$,
$n\le 7$ ($i^*=3$).
\end{enumerate}

\end{corollary}

\begin{proof}
For each case we exhibit the data of Theorem~\ref{thm:rowstree} with
$\varphi=\Phi_r(Z^*)$, where $Z^*\supseteq[1,i^*-1]$ has $|Z^*|=n+i^*-2$ and
satisfies the hypothesis of Proposition~\ref{prop:toprowcertS} with
$S=[1,i^*-1]$ and $L=n+i^*-1$ (for $i^*=1$ this is
Proposition~\ref{prop:toprowcert}).  So $\varphi=\nu_{i^*}(n)$, and
Theorem~\ref{thm:rowstree} gives $V_r(L,k)\ge1/\nu_{i^*}(n)$, while
Theorem~\ref{thm:prefixrows} gives $V_r(L,k)\le\min_{i\le k}1/\nu_i(n)$.
Hence $V_r(L,k)=1/\nu_{i^*}(n)=\min_i1/\nu_i(n)$.  The root witness is a
top-row optimum $Z_0$ with $\Phi_r(Z_0)\le\varphi$; the witnesses $W_T$ are
found by a local search and every inequality $\Phi_r(W_T)\le\varphi$ is
checked in exact arithmetic, as are the certificates and the thresholds
$\theta_T$ (condition~\eqref{eq:farappend} at $\tau=\theta_T$, for some
$D$).  The largest tree has about $5.4\cdot10^5$ nodes ($r=2$, $k=6$,
$n=4$).
\end{proof}

At $r=2$ these rows sharpen \cite[Theorem~4.16]{coefficient-mass}, since
$\beta_2(n)>n$ for $n\ge3$.  At $r=\frac32$ and $L=5$ all five rows are now
exact: the approximate values $2.98,\,1.99,\,0.80,\,0.20,\,\frac1{32}$
quoted in the paragraph on roots between one and two are
\[
  \frac{215233605}{72256243},\quad \frac{172186884}{86336203},\quad
  \frac{2657205}{3318724},\quad \frac{2302911}{11419568},\quad \frac1{32},
\]
the first by Proposition~\ref{prop:toprowcert} with $Z_0=\{2,5,9,16\}$, the
last by Corollary~\ref{cor:lastrowr}, and the others by (c), with
$i^*=k$ and $Z^*=\{1,4,9,16\}$, $\{1,2,8,15\}$, $\{1,2,3,13\}$.  In the cases computed at $r<2$ the minimizing
index is always $1$ or $k$: the worst exempted set is either the prefix
$[1,k-1]$ or a far block, as observed numerically after
Corollary~\ref{cor:belowtwo}, and the far block wins once $n$ is large
compared with $k$ (at $r=\frac{11}{10}$ this has not yet happened for
$n\le12$).  Theorem~\ref{thm:rowstree} treats one row at a time,
because Lemma~\ref{lem:farappend} needs a bound on the number of appended
zeros; for $r>2$, Theorems~\ref{thm:rowsfinite} and~\ref{thm:rowsrec} give
all $L$ at once.

\section{Monotone prefix values}\label{sec:prefixmono}

The conjectured value of the rows is $\min_{i\le k}1/\nu_i(n)$ (after
Corollary~\ref{cor:belowtwo}).  This section shows that the prefix values
$\nu_i(n)$ decrease geometrically in $i$, for every $r>1$ and every $n$
(Theorem~\ref{thm:prefixmono}).  So at $r\ge2$ the conjectured value is the
top row for every $n$, and the conjecture there says exactly that rows equal
the top row (Corollary~\ref{cor:prefixmono}).  The tool is a positive kernel
that \emph{dilates} a zero set instead of translating it.  This is the move
suggested at the end of Section~\ref{sec:rowsfinite}, where translating the
zeros past a block of exempted positions at the lowest hole fails near
$r=2$.  Lemma~\ref{lem:dilation} gives the exact cost of translation for
comparison.

Write $x^{\overline m}=x(x+1)\cdots(x+m-1)$ for the rising factorial, with
$x^{\overline0}=1$.  For each $d$, the polynomials $x^{\overline m}$ with
$m\le d$ form a basis of the real polynomials of degree at most $d$, and
$0^{\overline m}=0$ for $m\ge1$.

\begin{lemma}[A dilation kernel]\label{lem:polyakernel}
Let $i\ge1$.  For integers $s\ge i+1$ and $i\le u\le s-1$ put
\[
  K_i(s,u)=\binom{u-1}{i-1}\Big/\binom{s-1}{i} .
\]
Then, for every $m\ge0$,
\[
  \sum_{u=i}^{s-1}K_i(s,u)\,u^{\overline m}=\frac i{i+m}\,s^{\overline m} .
\]
In particular $K_i(s,\cdot)$ is a probability distribution on $[i,s-1]$.
Consequently, if $p=\sum_mc_mx^{\overline m}$ is a real polynomial and
$\hat p=\sum_m\frac i{i+m}\,c_mx^{\overline m}$, then $\deg\hat p\le\deg p$,
$\hat p(0)=p(0)$, and $\hat p(s)=\sum_{u=i}^{s-1}K_i(s,u)\,p(u)$ for every
$s\ge i+1$.
\end{lemma}

\begin{proof}
Since $u^{\overline m}=(u+m-1)!/(u-1)!$,
\[
  \binom{u-1}{i-1}u^{\overline m}
  =\frac{(u+m-1)!}{(i-1)!\,(u-i)!}
  =\frac{(i+m-1)!}{(i-1)!}\binom{u+m-1}{i+m-1}.
\]
By the hockey-stick identity,
$\sum_{u=i}^{s-1}\binom{u+m-1}{i+m-1}=\binom{s+m-1}{i+m}$.  Moreover
\[
  \binom{s+m-1}{i+m}\Big/\binom{s-1}i=\frac{(s+m-1)!}{(s-1)!}\cdot
  \frac{i!}{(i+m)!}=s^{\overline m}\,\frac{i!}{(i+m)!},
\]
and $\frac{(i+m-1)!}{(i-1)!}\cdot\frac{i!}{(i+m)!}=\frac i{i+m}$.  This is the
identity.  The case $m=0$ says that $K_i(s,\cdot)$ has total mass $1$, and it
is nonnegative.  The last claim follows by linearity, and
$\hat p(0)=c_0=p(0)$ because $0^{\overline m}=0$ for $m\ge1$.
\end{proof}

In the negative-binomial language of the paragraph on roots between one
and two, $K_i(s,\cdot)$ is the law of the time of the $i$-th success given
that the $(i+1)$-st success occurs at time $s$.  In that language the next
theorem says that the minimal absolute moment $m_i(n)=\inf_p\mathbb
E|p(X_i)|$ does not increase with $i$.

\begin{theorem}[Prefix values decrease]\label{thm:prefixmono}
Let $r>1$ and $n\ge1$.  Then for every $i\ge1$
\[
  \nu_{i+1}(n)\le a\,\nu_i(n),\qquad\text{equivalently}\qquad
  m_{i+1}(n)\le m_i(n).
\]
Hence $\nu_i(n)\le a^{i-1}\nu_1(n)=a^{i-1}/\beta_r(n)$ for all $i\ge1$.
The constant $a$ is attained at $n=1$.
\end{theorem}

\begin{proof}
Let $p$ be real with $\deg p\le n-1$ and $p(0)=1$, and let $\hat p$ be as in
Lemma~\ref{lem:polyakernel}.  Then $\deg\hat p\le n-1$ and $\hat p(0)=1$, so
$\hat p$ is admissible for $\nu_{i+1}(n)$.  The weight
$\binom{s-1}ir^{-s}$ vanishes for $s\le i$.  For $s\ge i+1$ the lemma gives
$|\hat p(s)|\le\sum_{u=i}^{s-1}K_i(s,u)|p(u)|$.  Since
$\binom{s-1}iK_i(s,u)=\binom{u-1}{i-1}$, exchanging the order of summation
(all terms are nonnegative) gives
\[
  \sum_{s\ge1}\binom{s-1}ir^{-s}|\hat p(s)|
  \le\sum_{u\ge i}\binom{u-1}{i-1}|p(u)|\sum_{s>u}r^{-s}
  =a\sum_{u\ge1}\binom{u-1}{i-1}r^{-u}|p(u)| ,
\]
because $\sum_{s>u}r^{-s}=r^{-u}/(r-1)$.  The left side is at least
$\nu_{i+1}(n)$.  Taking the infimum over $p$ gives
$\nu_{i+1}(n)\le a\nu_i(n)$, and with $\nu_i=a^im_i$ this is
$m_{i+1}(n)\le m_i(n)$.  Iterating from $\nu_1(n)=1/\beta_r(n)$ gives the
second claim.  For $n=1$ the only admissible polynomial is $p=1$, and
$\nu_i(1)=\sum_s\binom{s-1}{i-1}r^{-s}=a^i$.
\end{proof}

\begin{corollary}[Consequences for the rows]\label{cor:prefixmono}
Let $r>1$, $n,k\ge1$ and $L=n+k-1$.
\begin{enumerate}
\item[(a)] If $r\ge2$, then $\min_{1\le i\le k}1/\nu_i(n)=\beta_r(n)$.  So for
$r\ge2$ the conjecture stated after Corollary~\ref{cor:belowtwo} is exactly
the statement $V_r(L,k)=\beta_r(n)$ for all $L$ and $k$.  In particular the
index $i^*=1$ in Corollary~\ref{cor:rowstreevals}(a),(b) is forced.
\item[(b)] $\beta_r(n+j)\ge(r-1)^j\beta_r(n)$ for all $j\ge0$.
\item[(c)] Let $Z\supseteq[1,\ell]$ be finite with $|Z|\le n-1$.  Let $F$ be a
monic multiple of $(x-r)^L$, and suppose that its $k-1$ largest nonleading
magnitudes can be placed at backward distances forming a set
$S\subseteq[1,\ell+k-1]$, as in Proposition~\ref{prop:rowvalue}(a).  Then
\[
  b_k(F)\ \ge\ \frac{(r-1)^{k-1}}{\Phi_r(Z)} .
\]
In particular, if $\Phi_r(Z)=1/\beta_r(n)$, then
$b_k(F)\ge(r-1)^{k-1}\beta_r(n)$ for every such $F$.
\end{enumerate}
\end{corollary}

\begin{proof}
(a) For $r\ge2$ we have $a\le1$, so Theorem~\ref{thm:prefixmono} gives
$\nu_i(n)\le\nu_1(n)$ for every $i$, and the minimum of $1/\nu_i(n)$ is
$1/\nu_1(n)=\beta_r(n)$.  In Corollary~\ref{cor:rowstreevals} the value
$1/\nu_{i^*}(n)$ is this minimum.

(b) A polynomial of degree at most $n$ that vanishes at $1$ is admissible
for $\mu_r(\varnothing,n+1)$.  So
$1/\beta_r(n+1)=\mu_r(\varnothing,n+1)\le\mu_r(\{1\},n+1)=\nu_2(n)\le
a\nu_1(n)=a/\beta_r(n)$.  Iterate.

(c) We have $\ell\le|Z|\le n-1$.  Write $Z=[1,\ell]\cup Y$ with $Y\subset(\ell,\infty)$, so
$|Y|\le n-1-\ell$.  For $s\ge1$, $|r_Z(s)|=\binom{s-1}\ell|r_Y(s)|$.  Here $r_Y$
has degree at most $n-1-\ell$ and $r_Y(0)=1$.  So
$\Phi_r(Z)\ge\nu_{\ell+1}(n-\ell)$, and Theorem~\ref{thm:prefixmono} gives
$\nu_{\ell+k}(n-\ell)\le a^{k-1}\nu_{\ell+1}(n-\ell)\le a^{k-1}\Phi_r(Z)$.
Now let $p$ be real with $\deg p\le n-\ell-1$ and $p(0)=1$, and put
$q=r_{[1,\ell+k-1]}\,p$.  Then $\deg q\le L-1$, $q(0)=1$, and $q$ vanishes on
$[1,\ell+k-1]\supseteq S$.  Also
$\Phi_r(q)=\sum_s\binom{s-1}{\ell+k-1}r^{-s}|p(s)|$.  Proposition~\ref{prop:rowvalue}(a)
gives $b_k(F)\,\Phi_r(q)\ge1$.  Taking the infimum over $p$ gives
$b_k(F)\ge1/\nu_{\ell+k}(n-\ell)\ge(r-1)^{k-1}/\Phi_r(Z)$.
\end{proof}

Part (a) turns into a theorem the remark after
Corollary~\ref{cor:belowtwo} that, at $r\ge2$, the conjecture puts the
minimum at $i=1$.  For every $n$ and $k$, the prefix sets
$S=[1,i-1]$ completed by a far block (the sets of the upper bound in
Theorem~\ref{thm:prefixrows}) are never worse for the multiple than the far
block alone.  Part (c) says the same for every $F$ whose large coefficients
all sit among the top $\ell+k-1$ positions, where $\ell$ is the initial run
of any good top-row set.  For $r>2$ these $F$ even satisfy
$b_k(F)\ge(r-1)^{k-1}\beta_r(n)$, far above the row value
$V_r(L,k)\le\beta_r(n)$.  So the hard exempted sets are those that meet
the holes of the good zero sets.  For $1<r<2$ the theorem gives
$\nu_k(n)\le a^{k-1}\nu_1(n)$, and it gives that $m_k(n)$ does not increase
in $k$.  In Corollary~\ref{cor:belowtwo} this means that $(r-1)^{-k}$ times
the leading term $(r-1)^k/m_k(n)$ of $V_r(L,k)$ does not decrease in $k$.

The dilation should be compared with the translation that the chain of
Theorem~\ref{thm:rowsfinite} performs at the lowest hole.

\begin{lemma}[Translating past the lowest hole]\label{lem:dilation}
Let $W=[1,\ell]\cup Y$ be finite with $Y\subset[\ell+2,\infty)$, put
$w(t)=|r_W(t)|r^{-t}$, and for $x\ge0$ let $W^{(x)}=[1,\ell+x]\cup(Y+x)$.
Then
\[
  \Phi_r\bigl(W^{(x)}\bigr)=\frac{r^{-x}}{\prod_{z\in W}(1+x/z)}
  \sum_{t\ge1}\binom{t+x-1}x\,w(t).
  \tag{D}\label{eq:dilation}
\]
In particular, with $h=\ell+1$,
\[
  \Phi_r\bigl(\mathrm{ins}(W,h)\bigr)=
  \frac{\Phi_r(W)+\Phi_r(W\cup\{h\})}{r\prod_{y\in Y}(1+1/y)} .
\]
\end{lemma}

\begin{proof}
$r_{W^{(x)}}$ vanishes on $[1,\ell+x]$.  Let $t\ge\ell+1$.  Then
$|r_{[1,\ell+x]}(t+x)|=\binom{t+x-1}{\ell+x}$, and expanding the factorials
gives
\[
  \binom{t+x-1}{\ell+x}\binom{\ell+x}x=\binom{t-1}\ell\binom{t+x-1}x
  =\frac{(t+x-1)!}{\ell!\,x!\,(t-1-\ell)!} .
\]
Also $|r_{Y+x}(t+x)|=\prod_{y\in Y}|y-t|/(y+x)=|r_Y(t)|\prod_{y\in
Y}y/(y+x)$, and $\binom{\ell+x}x=\prod_{z=1}^{\ell}(1+x/z)$.  So
\[
  |r_{W^{(x)}}(t+x)|=\frac{\binom{t+x-1}x}{\prod_{z\in W}(1+x/z)}\,
  |r_W(t)|,
\]
using $|r_W(t)|=\binom{t-1}\ell|r_Y(t)|$.  Multiply by $r^{-t-x}$ and sum
over $t\ge\ell+1$.  The terms $t\le\ell$ of \eqref{eq:dilation} vanish because
$w(t)=0$ there.

For $x=1$ the set $W^{(1)}$ is $\mathrm{ins}(W,h)$.  Moreover
$\prod_{z\in W}(1+1/z)=h\prod_{y\in Y}(1+1/y)$.  Since $w(t)=0$ for
$t<h$ and $|1-t/h|=t/h-1$ for $t\ge h$,
\[
  \Phi_r(W\cup\{h\})=\sum_tw(t)\Bigl(\frac th-1\Bigr)
  =\frac1h\sum_ttw(t)-\Phi_r(W).
\]
Insert this into \eqref{eq:dilation}.
\end{proof}

For $W=Z_0$ the sets $W^{(x)}$ are the uniform shifts $Z_0^{(x,\ldots,x)}$
of Lemma~\ref{lem:holestates}.  They are the states reached when the
exempted positions begin with a block at the lowest hole.  By \eqref{eq:dilation},
translating the zeros costs the $w$-average of
$\binom{t+x-1}xr^{-x}/\prod_{z}(1+x/z)$.  This factor is large when the
mass of $w$ sits far to the right of the zeros.  At $r=2$, $Z_0=\{2\}$ it
reproduces $\Phi_2([1,x]\cup\{x+2\})\to1$ from
Section~\ref{sec:rowsfinite}.

Re-optimizing costs at most $a^x$.  Let $n-1=|Z_0|$ and
$Z_0=[1,\ell]\cup Y$ with $Y\subset[\ell+2,\infty)$.  The sets $W'=[1,\ell+x]\cup Y'$ with
$Y'\subset(\ell+x,\infty)$ and $|Y'|=n-1-\ell$ have infimum
$\nu_{\ell+x+1}(n-\ell)$ of $\Phi_r$ (Lemma~\ref{lem:intzeros}, with the zeros at least $\ell+x+1$).  By Theorem~\ref{thm:prefixmono}, this infimum is at most
$a^x\nu_{\ell+1}(n-\ell)\le a^x\Phi_r(Z_0)$.

So for $r>2$ a cut of Theorem~\ref{thm:rowsrec} at a uniform shift, at
$t'=\ell+x$, always admits a set $W'$ with
$W'\cap[1,t']=Z_0^{(x,\ldots,x)}\cap[1,t']$, $|W'|=|Z_0^{(x,\ldots,x)}|$ and
$\Phi_r(W')<\Phi_r(Z_0)$.  Only the recursive certification of $(W',t')$
remains to be done.  This agrees with the observation, after
Corollary~\ref{cor:rowsfiniteranges}, that the re-optimized witness after
such a cut has much smaller $\Phi_r$.

The dilation does not extend to holes above the lowest one.  The part of
the zero set below such a hole is not a full run, and the kernel of
Lemma~\ref{lem:polyakernel} does not apply.  One could hope for another
kernel that works for a general exempted set $S$ and all $n$ at once.
Suppose $K\ge0$ is a kernel from $\N_{>0}\setminus S$ to $\N_{>0}$ such
that, for every polynomial $p$, the function $s\mapsto\sum_uK(s,u)p(u)$
agrees on $\N_{>0}\setminus S$ with a polynomial $\hat p$ with
$\deg\hat p\le\deg p$ and $\hat p(0)=p(0)$.  Suppose also that its
push-forward of the weight $|r_S(s)|r^{-s}$ is at most $r^{-u}$.  Then, as in
the proof of Theorem~\ref{thm:prefixmono}, $\mu_r(S,n+|S|)\le1/\beta_r(n)$
for every $n$.  No such kernel exists when $S\ne\varnothing$ and
$1\notin S$.  Indeed, $K$ is stochastic, since $\hat1=1$.  The polynomial
$\hat x$ has degree at most $1$ and vanishes at $0$, so it is $\lambda x$.
At $s=1$ we get $\lambda=\sum_uK(1,u)u\ge1$.  By Jensen's inequality,
$\sum_uK(s,u)\theta^u\ge\theta^{\lambda s}\ge\theta^s$ for $\theta>1$.
Integrating against the weight gives
\[
  \sum_{s\notin S}|r_S(s)|y^s\ \le\ \sum_{u\ge1}y^u,\qquad y=\theta/r<1 .
\]
Since $|r_S(s)|\ge c\,s^{|S|}$ for large $s$ and some $c>0$, the left side
is at least a constant times $(1-y)^{-|S|-1}$ as $y\to1$, while the right
side is $y/(1-y)$.  This is a contradiction.  So for a general exempted set
a dilation must depend on $n$.

\section{Crossing an exempted position}\label{sec:crossing}

By Corollary~\ref{cor:prefixmono}(a), for $r\ge2$ the conjecture on the
rows is the statement $V_r(L,k)=\beta_r(n)$.  The chains of
Theorems~\ref{thm:rowspi}, \ref{thm:rowsfinite} and~\ref{thm:rowsrec} insert
a zero at every exempted position that is a hole, and each insertion has to
be paid for.  This section shows that the last exempted position costs at
most one insertion, made at a place that does not depend on that position.
Two facts give this.  The polynomials admissible for
Proposition~\ref{prop:rowvalue}(a) with $\Phi_r\le\varphi$ form a convex set,
so a position is covered as soon as two admissible polynomials take opposite
signs there (Lemma~\ref{lem:crossing}).  And an insertion reverses the sign
of $r_W$ at every later position that lies in neither set
(Lemma~\ref{lem:insflip}).

The consequences are these.
\begin{itemize}
\item One inequality covers all later positions at once
(Theorem~\ref{thm:crossing}), for every $r>1$.
\item The second row follows from a single inequality for each $n$.  This
gives $V_r(n+1,2)=\beta_r(n)$ in many new cases, among them $r=2$ for every
$n\le30$ (Corollary~\ref{cor:secondrow}).
\item The conjecture at $r\ge2$ reduces to one insertion at constrained
minimizers, made at the first gap above the exempted positions only
(Corollary~\ref{cor:crossreduce}).
\item The finite trees of Theorem~\ref{thm:rowstree} lose their last level
(Theorem~\ref{thm:crosstree}).  At $r=2$ this gives many new rows
(Corollary~\ref{cor:crosstreevals}).
\item Combined with the dilation of Section~\ref{sec:prefixmono}, it gives a
bound valid for every $n$ and $k$ for the exempted sets all of whose elements
but the largest lie in $[1,\ell+k-2]$ (Proposition~\ref{prop:dilatecross}).
\end{itemize}

For a finite $W\subset\N_{>0}$ and an integer $t\ge0$ put
\[
  h_t(W)=\min\{x\in\N:\ x>t,\ x\notin W\},
\]
the first gap of $W$ above $t$.  Every integer in $(t,h_t(W))$ lies in $W$.
If $h>\max W$ then $\mathrm{ins}(W,h)=W\cup\{h\}$.

\begin{lemma}[Crossing]\label{lem:crossing}
Let $r>1$, $L\ge1$, $\varphi>0$, let $T\subset\N_{>0}$ be finite and
$\sigma\in\N_{>0}\setminus T$.
\begin{enumerate}
\item[(a)] Let $q_1,q_2$ be real polynomials with $\deg q_i\le L-1$,
$q_i(0)=1$, $q_i|_T=0$ and $\Phi_r(q_i)\le\varphi$.  If
$q_1(\sigma)q_2(\sigma)\le0$, then $|T|+1<L$ and
$\mu_r(T\cup\{\sigma\},L)\le\varphi$.
\item[(b)] For every real polynomial $Q$ and all reals $0<y\le\sigma'$,
\[
  \Phi_r\bigl(Q\cdot(1-s/\sigma')\bigr)\le
  \max\bigl\{\Phi_r(Q),\ \Phi_r\bigl(Q\cdot(1-s/y)\bigr)\bigr\}.
\]
\end{enumerate}
\end{lemma}

\begin{proof}
(a) If $q_1(\sigma)=0$ put $q=q_1$, and if $q_2(\sigma)=0$ put $q=q_2$.
Otherwise $\lambda=q_2(\sigma)/(q_2(\sigma)-q_1(\sigma))$ lies in $(0,1)$;
put $q=\lambda q_1+(1-\lambda)q_2$.  Then $q(\sigma)=0$, $q(0)=1$, $q|_T=0$
and $\deg q\le L-1$.  Moreover
$|q(s)|\le\lambda|q_1(s)|+(1-\lambda)|q_2(s)|$ for every $s$, so
$\Phi_r(q)\le\varphi$.  The polynomial $q$ is nonzero and vanishes at the
$|T|+1$ points of $T\cup\{\sigma\}$, so $|T|+1\le L-1$, and $q$ is admissible
for $\mu_r(T\cup\{\sigma\},L)$.

(b) With $\lambda=y/\sigma'\in(0,1]$ we have
$1-s/\sigma'=(1-\lambda)+\lambda(1-s/y)$.  So
$|Q(s)(1-s/\sigma')|\le(1-\lambda)|Q(s)|+\lambda|Q(s)(1-s/y)|$ for every $s$.
Multiply by $r^{-s}$ and sum.
\end{proof}

\begin{lemma}[An insertion reverses the later signs]\label{lem:insflip}
Let $W\subset\N_{>0}$ be finite, let $h\in\N_{>0}\setminus W$ and
$W'=\mathrm{ins}(W,h)$.  Then:
\begin{itemize}
\item $W'\cap[1,h-1]=W\cap[1,h-1]$, $h\in W'$ and $|W'|=|W|+1$;
\item for every integer $\sigma>h$ with $\sigma\notin W\cup W'$,
$r_W(\sigma)\,r_{W'}(\sigma)<0$.
\end{itemize}
\end{lemma}

\begin{proof}
The first three claims are the definition of $\mathrm{ins}$ (as $h\notin W$,
the elements of $W$ that are at least $h$ exceed $h$).  For a finite $Z$ and
an integer $\sigma\notin Z$, the factor $1-\sigma/z$ of $r_Z(\sigma)$ is
negative exactly when $z<\sigma$, so $r_Z(\sigma)$ has the sign
$(-1)^{|Z\cap[1,\sigma-1]|}$.

Let $\sigma>h$ with $\sigma\notin W\cup W'$, and put
$A=|W\cap[1,h-1]|$ and $B=|W\cap[h+1,\sigma-2]|$.  The elements of $W'$
above $h$ are the $w+1$ with $w\in W$, $w>h$.  Those below $\sigma$ come
from $w\le\sigma-2$, so $|W'\cap[1,\sigma-1]|=A+1+B$.  Next,
$|W\cap[1,\sigma-1]|=A+B+[\sigma-1\in W]$, since $h\notin W$.  If
$\sigma-1\in W$, then $\sigma-1\ne h$, so $\sigma-1>h$ and $\sigma\in W'$,
which is excluded.  So $|W\cap[1,\sigma-1]|=A+B$, and the two counts differ
by one.
\end{proof}

\begin{theorem}[One insertion covers every later position]\label{thm:crossing}
Let $r>1$, let $t\ge0$ be an integer, let $W\subset\N_{>0}$ be finite, let
$T\subseteq W\cap[1,t]$ and $h=h_t(W)$.  Then for every integer $\sigma>t$,
\[
  \mu_r\bigl(T\cup\{\sigma\},\,|W|+2\bigr)\ \le\
  \max\bigl\{\Phi_r(W),\ \Phi_r(\mathrm{ins}(W,h))\bigr\}.
\]
Consequently, let $F$ be a monic multiple of $(x-r)^L$, let $S$ be the
exempted set of Proposition~\ref{prop:rowvalue}(a) with $|S|=k-1\ge1$, and
let $t=\max(S\setminus\{\max S\})$ ($t=0$ if $k=2$).  If $W\supseteq S\setminus\{\max S\}$
and $|W|\le L-2$, then
\[
  b_k(F)\ \ge\ \frac1{\max\{\Phi_r(W),\,\Phi_r(\mathrm{ins}(W,h_t(W)))\}} .
\]
\end{theorem}

\begin{proof}
Put $W'=\mathrm{ins}(W,h)$.  The polynomials $r_W$ and $r_{W'}$ have degree
at most $|W|+1$, take the value $1$ at $0$, and vanish on $T$, because
$T\subseteq W\cap[1,t]\subseteq W\cap[1,h-1]\subseteq W'$.  If $\sigma<h$,
then $\sigma\in(t,h)\subseteq W$ and $r_W$ vanishes at $\sigma$.  If
$\sigma=h$, then $r_{W'}$ vanishes at $\sigma$.  If $\sigma>h$ and
$\sigma\in W\cup W'$, one of the two vanishes at $\sigma$.  Otherwise
$r_W(\sigma)r_{W'}(\sigma)<0$ by Lemma~\ref{lem:insflip}.  In every case
Lemma~\ref{lem:crossing}(a), with $L=|W|+2$ and
$\varphi=\max\{\Phi_r(W),\Phi_r(W')\}$, gives the first claim.

For the second, apply the first with $T=S\setminus\{\max S\}$ and
$\sigma=\max S$.  This gives polynomials $q$ with $\deg q\le|W|+1\le L-1$,
$q(0)=1$, $q|_S=0$ and $\Phi_r(q)$ arbitrarily close to
$\mu_r(S,|W|+2)$.  Proposition~\ref{prop:rowvalue}(a) gives
$b_k(F)\Phi_r(q)\ge1$ for each of them.
\end{proof}

No condition on $r$ is needed.  When $h>\max W$ the second set is
$W\cup\{h\}$, and for $r\ge2$ Lemma~\ref{lem:confdelr} gives
$\Phi_r(W\cup\{h\})\le\Phi_r(W)$.  So a set $W$ without gaps in
$(t,\max W]$ handles every later position for free, which is the third rule
of the chain of Theorem~\ref{thm:rowspi}.  The theorem lets that chain stop
one step early.  Instead of inserting at the last exempted position, it
inserts once at the first gap above the previous ones, and for all later
positions at the same time.

\paragraph{The second row.}
For $k=2$ there is nothing before the exempted position, and the theorem
reduces to one inequality per $n$.  Its right-hand side is explicit by
Lemma~\ref{lem:dilation}.

\begin{corollary}[The second row from one insertion]\label{cor:secondrow}
Let $r>1$ and $n\ge1$.  Let $Z=[1,\ell]\cup Y$ be finite with
$Y\subset[\ell+2,\infty)$ and $|Z|\le n-1$, and put
$\Pi_Y=\prod_{y\in Y}(1+1/y)$ and
$\rho=\Phi_r(Z\cup\{\ell+1\})/\Phi_r(Z)$.  Then
\[
  \Phi_r\bigl(\mathrm{ins}(Z,\ell+1)\bigr)=\frac{1+\rho}{r\,\Pi_Y}\,\Phi_r(Z)
  \qquad\text{and}\qquad
  V_r(n+1,2)\ \ge\ \frac1{\Phi_r(Z)\max\{1,(1+\rho)/(r\Pi_Y)\}} .
\]
In particular, if $\Phi_r(Z)=1/\beta_r(n)$ and $1+\rho\le r\Pi_Y$, then
$V_r(n+1,2)=\beta_r(n)$.
\end{corollary}

\begin{proof}
The identity is the second formula of Lemma~\ref{lem:dilation} with $W=Z$.
The first gap of $Z$ above $0$ is $h_0(Z)=\ell+1$, and $|Z|\le n-1=L-2$ for
$L=n+1$.  So Theorem~\ref{thm:crossing} with $k=2$ gives the lower bound for
every monic multiple of $(x-r)^{n+1}$.  The upper bound
$V_r(n+1,2)\le\beta_r(n)$ is Proposition~\ref{prop:rowvalue}(c).
\end{proof}

The condition $1+\rho\le r\Pi_Y$ says that the mean of $s$ under the weight
$|r_Z(s)|r^{-s}$ is at most $r(\ell+1)\Pi_Y$, since
$\rho=\sum_s(s/(\ell+1)-1)|r_Z(s)|r^{-s}/\Phi_r(Z)$.  We checked it in exact
arithmetic at top-row optima $Z_0$ certified by
Proposition~\ref{prop:toprowcert}.  This gives $V_r(n+1,2)=\beta_r(n)$ for
\[
  r=2,\ n\le30;\quad r=\tfrac94,\ n\le32;\quad r=\tfrac52,\ n\le36;\quad
  r=3,\ n\le42;\quad r=4,\ n\le50;\quad r=\tfrac{19}{10},\ n\le19;
\]
and for $r=\frac74$ with $n\in\{4,6,7,8,10,11,12,13,16,17\}$.  At $r=2$ the ratio
$(1+\rho)/(2\Pi_Y)$ is $1$ at $n=2$, where $Z_0=\{2\}$ and
$\Phi_2(\{1,3\})=\Phi_2(\{2\})=\frac12$.  It is $0.93$ at $n=4$ and at most
$0.78$ for $5\le n\le30$.  At $r=3$ it is at most $0.47$ for $4\le n\le42$.  For
$r=2$ this extends Corollary~\ref{cor:rowstreevals}(a) with $k=2$ from
$n\le14$ to $n\le30$, and it needs one exact inequality per $n$ instead of a
tree.  Below $r=2$ the criterion fails at some $n$: at $r=\frac74$ for
$n\in\{2,3,5,9,14,15\}$, at $r=\frac32$ for every $n\le12$, and at
$r=\frac54$ for every $n\le8$.  There the append steps are not free.  At
$r=\frac74$ with $n=2$, at $r=\frac32$ with $n\le4$ and at $r=\frac54$
with $n\le8$, the minimizing index of Corollary~\ref{cor:rowstreevals} is
$i^*=2$, so no criterion with target $\beta_r(n)$ can hold there.

\paragraph{The conjecture reduces to the first gap.}
The re-optimizing chain replaces, after each exempted position, the current
set by a minimizer for the positions seen so far.  By
Theorem~\ref{thm:crossing} it needs only the following inequality at each
step.

\begin{corollary}[Reduction to the first gap]\label{cor:crossreduce}
Let $r>1$ and $n\ge1$.  Suppose that every finite $T\subset\N_{>0}$ admits
a set $W_T\supseteq T$ with $|W_T|\le|T|+n-1$,
$\Phi_r(W_T)=\mu_r(T,|T|+n)$ and
\[
  \Phi_r\bigl(\mathrm{ins}(W_T,h_{\max T}(W_T))\bigr)\le\Phi_r(W_T)
  \qquad(\max\varnothing=0).
\]
Then $\mu_r(S,|S|+n)\le1/\beta_r(n)$ for every finite $S$, and
$V_r(L,L-n+1)=\beta_r(n)$ for every $L\ge n$.
\end{corollary}

\begin{proof}
Let $S\ne\varnothing$ be finite, $\sigma=\max S$ and $T=S\setminus\{\sigma\}$.
Theorem~\ref{thm:crossing} with $W=W_T$ and $t=\max T$ gives
$\mu_r(S,|W_T|+2)\le\Phi_r(W_T)=\mu_r(T,|T|+n)$.  Since
$|W_T|+2\le|S|+n$ and $\mu_r(S,\cdot)$ does not increase when more degree is
allowed, $\mu_r(S,|S|+n)\le\mu_r(T,|T|+n)$.  By induction on $|S|$,
$\mu_r(S,|S|+n)\le\mu_r(\varnothing,n)=1/\beta_r(n)$.
Proposition~\ref{prop:rowvalue}(b) with $k=|S|+1$ then gives
$V_r(L,k)\ge\beta_r(n)$, and (c) gives the reverse inequality.
\end{proof}

When $h_{\max T}(W_T)>\max W_T$ the hypothesis is $\Phi_r(W_T\cup\{h\})\le\Phi_r(W_T)$,
which holds for $r\ge2$ by Lemma~\ref{lem:confdelr}.  So at $r\ge2$ only
constrained minimizers with a gap above $\max T$ matter, and only their first
gap.  This is much weaker than the hypothesis of
Proposition~\ref{prop:insertion}, which asks for every $Z$ and every $\tau$
and fails at $r=2$, $\frac94$ and $\frac{23}{10}$ (the remarks after that
proposition).  It is also weaker than asking for the insertion at every hole
of $W_T$ above $\max T$.  When $T=[1,j]$, $W_T$ is a minimizer for
$\nu_{j+1}(n)$ of the form $[1,\ell]\cup Y$ with $\ell\ge j$ and
$Y\subset[\ell+2,\infty)$, and the hypothesis is the inequality
$1+\rho\le r\Pi_Y$ of Corollary~\ref{cor:secondrow} for $Z=W_T$.  In
floating point, with minimizers found by exhaustive search in a window, the
largest ratio $\Phi_r(\mathrm{ins}(W_T,h))/\Phi_r(W_T)$ was:
\begin{itemize}
\item $r=2$: $0.940$ over all $T\subseteq[1,9]$ with $|T|\le3$ at $n=4$
(attained at $T=\{2,3\}$, $W_T=\{1,2,3,7,11\}$); $0.859$ over all
$T\subseteq[1,7]$ with $|T|\le3$ at $n=5$;
\item $r=\frac52$: $0.64$, and $r=3$: $0.60$, at $n=5$.
\end{itemize}
No minimizer reached the edge of its window.  At $r=2$ the worst cases have $W_T\cap[1,h-1]$ an initial run.
Lemma~\ref{lem:optins} proves the hypothesis in a truncated form at every
$r\ge2$, and Theorem~\ref{thm:allrowstwo} deduces the conclusion there.

\paragraph{Trees without their last level.}
Theorem~\ref{thm:crossing} also shortens the trees of
Theorem~\ref{thm:rowstree}.  Fix $r>1$ and $\varphi>0$.  For a finite $W$
and $j\ge1$, call an integer $\theta>\max W$ a \emph{$j$-threshold} of $W$
if $\Phi_r(W\cup U)\le\Phi_r(W)$ for every finite $U\subset[\theta,\infty)$
with $|U|\le j$.  For $r\ge2$, $\theta=\max W+1$ is a $j$-threshold for
every $j$, by Lemma~\ref{lem:confdelr}.  For every $r>1$,
Lemma~\ref{lem:farappend} with $u=j$ gives one.  Consider pairs $(W,t)$ of
a finite $W\subset\N_{>0}$ with $\Phi_r(W)\le\varphi$ and an integer $t\ge0$,
and define safety by induction on $j$.
\begin{itemize}
\item $(W,t)$ is \emph{$1$-safe} if
$\Phi_r(\mathrm{ins}(W,h_t(W)))\le\varphi$.
\item For $j\ge2$, $(W,t)$ is \emph{$j$-safe} if $W$ has a $j$-threshold
$\theta$ such that every integer $x$ with $t<x<\theta$ satisfies the
following.  If $x\in W$, the pair $(W,x)$ is $(j-1)$-safe.  If $x\notin W$,
some $(j-1)$-safe pair $(W_x,x)$ has $|W_x|\le|W|+1$ and
$(W\cap[1,t])\cup\{x\}\subseteq W_x$.
\end{itemize}

\begin{theorem}[Crossing trees]\label{thm:crosstree}
Let $r>1$, $n\ge1$, $k\ge2$ and $\varphi>0$, and let $Z_0\subset\N_{>0}$
with $|Z_0|\le n-1$ be such that $(Z_0,0)$ is $(k-1)$-safe.  Then every
monic multiple $F$ of $(x-r)^{n+k-1}$ has $b_k(F)\ge1/\varphi$, that is,
$V_r(n+k-1,k)\ge1/\varphi$.
\end{theorem}

\begin{proof}
We show by induction on $j$ that, if $(W,t)$ is $j$-safe, then
\[
  \mu_r\bigl((W\cap[1,t])\cup U,\ |W|+j+1\bigr)\le\varphi
\]
for every $U\subset(t,\infty)$ with $|U|=j$.  For $j=1$ this is
Theorem~\ref{thm:crossing} with $T=W\cap[1,t]$.  Let $j\ge2$, let $\theta$
be as in the definition, and let $x=\min U$ and $U'=U\setminus\{x\}$.
\begin{itemize}
\item If $x\ge\theta$, then $U\subset[\theta,\infty)$ and
$\Phi_r(W\cup U)\le\Phi_r(W)\le\varphi$, with $|W\cup U|=|W|+j$.
\item If $x<\theta$ and $x\in W$, the induction hypothesis for $(W,x)$ and
$U'\subset(x,\infty)$ bounds $\mu_r((W\cap[1,x])\cup U',|W|+j)$ by
$\varphi$.  The set $(W\cap[1,x])\cup U'$ contains $(W\cap[1,t])\cup U$.
\item If $x<\theta$ and $x\notin W$, the induction hypothesis for $(W_x,x)$
bounds $\mu_r((W_x\cap[1,x])\cup U',|W_x|+j)$ by $\varphi$.  Here
$W_x\cap[1,x]\supseteq(W\cap[1,t])\cup\{x\}$ and $|W_x|+j\le|W|+j+1$.
\end{itemize}
In each case the claim follows, since $\mu_r(S,L)$ can only decrease when
points are removed from $S$ or when $L$ grows.  For
$(Z_0,0)$ and $U=S$, the exempted set of Proposition~\ref{prop:rowvalue}(a)
with $|S|=k-1$, this gives $\mu_r(S,n+k-1)\le\varphi$, since
$|Z_0|+k\le n+k-1$.  Proposition~\ref{prop:rowvalue}(a) gives the bound.
\end{proof}

The recursion is finite: each level branches over the finitely many $x$ in
$(t,\theta)$, and there are $k-1$ levels.  The last level is a single
inequality, where Theorem~\ref{thm:rowstree} branches over every position
below a threshold.  For $r\ge2$ with $\theta=\max W+1$, a pair with
$t\ge\max W$ is $j$-safe for every $j$: for $j\ge2$ there is no $x$ to check,
and for $j=1$ the inserted set is $W\cup\{t+1\}$.  For $k=2$ the theorem is
Corollary~\ref{cor:secondrow}.

\begin{corollary}[Rows at the root $2$ from crossing trees]\label{cor:crosstreevals}
$V_2(L,k)=\beta_2(L-k+1)$ for every $L$ with $n=L-k+1$ in the following
ranges: $k=3$, $n\le30$; $k=4$, $n\le22$; $k=5$, $n\le18$; $k=6$, $n\le13$;
$k=7$, $n\le10$; $k=8$, $n\le9$; $k=9$, $n\le8$; $k=10$, $n\le7$.
\end{corollary}

\begin{proof}
Take $\varphi=\Phi_2(Z_0)$ for a top-row optimum $Z_0$ certified by
Proposition~\ref{prop:toprowcert}, so that $1/\varphi=\beta_2(n)$, and
$\theta=\max W+1$.  For $x\notin W$ the set $W_x$ is $\mathrm{ins}(W,x)$ or,
if that fails, a local re-optimization of $\mathrm{ins}(W,x)$ with the elements
of $(W\cap[1,t])\cup\{x\}$ held fixed.  Every inequality
$\Phi_2(W_x)\le\varphi$ and every $1$-safety inequality was checked in exact
arithmetic.  Theorem~\ref{thm:crosstree} gives $V_2(L,k)\ge\beta_2(n)$, and
Proposition~\ref{prop:rowvalue}(c) gives the reverse inequality.  The
largest recursion ($k=10$, $n=7$) has about $5\cdot10^4$ nodes.
\end{proof}

This extends Corollary~\ref{cor:rowstreevals}(a).  There the ranges were
$k=3$, $n\le12$; $k=4$, $n\le8$; $k=5$, $n\le5$; $k=6$, $n\le4$, and the
largest tree had about $5.4\cdot10^5$ nodes.

\paragraph{Dilate, then cross.}
Combined with the dilation kernel of Lemma~\ref{lem:polyakernel},
Theorem~\ref{thm:crossing} admits one arbitrary exempted position on top of
an initial block, for every $n$.  For $i\ge1$ write
$\omega_i(s)=\binom{s-1}{i-1}r^{-s}$, so that
$\Phi_r([1,i-1]\cup Y)=\sum_s\omega_i(s)|r_Y(s)|$ for
$Y\subset[i,\infty)$.  The first step is to track the weight with one binomial
index more through the kernel.

\begin{lemma}[The excess weight under the kernel]\label{lem:excess}
Let $r>1$, $i\ge1$, let $p$ be a real polynomial and let $\hat p$ be as in
Lemma~\ref{lem:polyakernel}.  Then
\[
  \sum_{s\ge1}\omega_{i+2}(s)|\hat p(s)|\ \le\
  \frac a{i+1}\sum_{u\ge1}\bigl(i\,\omega_{i+1}(u)+a\,\omega_i(u)\bigr)|p(u)| .
\]
\end{lemma}

\begin{proof}
$\omega_{i+2}(s)=0$ for $s\le i+1$.  For $s\ge i+2$,
Lemma~\ref{lem:polyakernel} gives
$|\hat p(s)|\le\sum_{u=i}^{s-1}K_i(s,u)|p(u)|$, and
\[
  \omega_{i+2}(s)K_i(s,u)=\frac{\binom{s-1}{i+1}}{\binom{s-1}i}
  \binom{u-1}{i-1}r^{-s}=\frac{s-1-i}{i+1}\binom{u-1}{i-1}r^{-s} .
\]
This vanishes at $s=i+1$, so the sum may run over all $s>u$.  For $u\ge i$,
$\sum_{s>u}(s-1-i)r^{-s}=r^{-u}\sum_{m\ge1}(u-i-1+m)r^{-m}=ar^{-u}(u-i+a)$,
using $\sum_{m\ge1}r^{-m}=a$ and $\sum_{m\ge1}mr^{-m}=a(1+a)$.  All terms are
nonnegative, so exchanging the sums gives
\[
  \sum_s\omega_{i+2}(s)|\hat p(s)|\le\frac a{i+1}\sum_{u\ge i}
  \binom{u-1}{i-1}(u-i+a)r^{-u}|p(u)| .
\]
Finally $\binom{u-1}{i-1}(u-i)=i\binom{u-1}i$.
\end{proof}

\begin{proposition}[Dilate, then cross]\label{prop:dilatecross}
Let $r>1$, $n\ge1$, $k\ge2$ and $m=k-2$.  Let $Z=[1,\ell]\cup Y$ with
$Y\subset[\ell+2,\infty)$ and $|Z|\le n-1$, and put
$\rho=\Phi_r(Z\cup\{\ell+1\})/\Phi_r(Z)$.  Let $F$ be a monic multiple of
$(x-r)^{n+k-1}$ whose exempted set $S$
(Proposition~\ref{prop:rowvalue}(a)) satisfies
$S\setminus\{\max S\}\subseteq[1,\ell+m]$.  Then
\[
  b_k(F)\ \ge\ \frac{(r-1)^m}{\Phi_r(Z)}\cdot
  \frac1{\max\bigl\{1,\ ((\ell+1)\rho+ma)/(\ell+m+1)\bigr\}}
  \ \ge\ \frac{(r-1)^{k-2}}{\max\{1,a,\rho\}\,\Phi_r(Z)} .
\]
\end{proposition}

\begin{proof}
Put $p_0=r_Y$ and, for $0\le j<m$, let $p_{j+1}=\hat p_j$ be given by
Lemma~\ref{lem:polyakernel} with index $i=\ell+1+j$.  Each $p_j$ has degree
at most $|Y|$ and $p_j(0)=1$.  Put
$A_j=\sum_s\omega_{\ell+1+j}(s)|p_j(s)|$ and
$B_j=\sum_s\omega_{\ell+2+j}(s)|p_j(s)|$.  Then $A_0=\Phi_r(Z)$ and
$B_0=\Phi_r(Z\cup\{\ell+1\})=\rho A_0$, because $\ell+1\notin Y$.  The
proof of Theorem~\ref{thm:prefixmono} gives $A_{j+1}\le aA_j$, hence
$A_j\le a^jA_0$.  Lemma~\ref{lem:excess} with $i=\ell+1+j$ gives
\[
  B_{j+1}\le\frac a{\ell+2+j}\bigl((\ell+1+j)B_j+aA_j\bigr).
\]
By induction on $j$ this yields
$B_j\le a^j\bigl((\ell+1)B_0+jaA_0\bigr)/(\ell+1+j)$: inserting the bound for
$B_j$ and $A_j\le a^jA_0$ into the right side gives
$a^{j+1}((\ell+1)B_0+(j+1)aA_0)/(\ell+2+j)$.

Let $J=\ell+m+1$ and $Q=r_{[1,J-1]}\,p_m$.  Since
$|r_{[1,J-1]}(s)|=\binom{s-1}{J-1}$,
\[
  \Phi_r(Q)=A_m\le a^m\Phi_r(Z),\qquad
  \Phi_r\bigl(Q\cdot(1-s/J)\bigr)=B_m\le a^m\,\frac{(\ell+1)\rho+ma}{J}\,\Phi_r(Z).
\]
Also $\deg Q\le\ell+m+|Y|\le n+m-1$ and $Q(0)=1$.  Let $\sigma=\max S$.  If
$\sigma\le J-1$, then $Q$ vanishes on $S$.  Otherwise
$Q\cdot(1-s/\sigma)$ vanishes on $S$ and has degree at most $n+m=L-1$, where
$L=n+k-1$.  By Lemma~\ref{lem:crossing}(b) with $y=J\le\sigma$, its
$\Phi_r$ is at most $\max\{A_m,B_m\}$.  Proposition~\ref{prop:rowvalue}(a)
gives the first bound.  For the second, $((\ell+1)\rho+ma)/(\ell+m+1)$ is an
average of $\rho$ and $a$.
\end{proof}

For $m=0$ the set $S$ is a single arbitrary position.  The proposition then
gives $b_2\ge1/(\max\{1,\rho\}\,\Phi_r(Z))$, and
Corollary~\ref{cor:secondrow} improves the maximum to
$\max\{1,(1+\rho)/(r\Pi_Y)\}$.  For $m\ge1$ the proposition
extends Corollary~\ref{cor:prefixmono}(c) at the price of one position of the
window.  There, all $k-1$ positions lie in $[1,\ell+k-1]$ and the bound is
$(r-1)^{k-1}/\Phi_r(Z)$.  Here $k-2$ positions lie in $[1,\ell+k-2]$, the
last one is arbitrary, and the loss is the factor
$(r-1)\max\{1,a,\rho\}$ at most.  The bound holds for every $n$ and every
$Z$ of this form, optimal or not.  With a top-row optimum it reads as
follows.
\begin{itemize}
\item \emph{All $n$, $r>2$.}  If $Z_0=[1,\ell]\cup Y$ is a top-row optimum
and $\rho=\rho(Z_0)$, every such $F$ has
$b_k(F)\ge(r-1)^{k-2}\beta_r(n)/\max\{1,\rho\}$.  This is at least
$\beta_r(n)$ as soon as $(r-1)^{k-2}\ge\rho$.
\item \emph{All $k$ at once.}  If $r\ge2$ and $\rho(Z_0)\le r-1$, then
$b_k(F)\ge\beta_r(n)$ for every $k\ge2$ and every $F$ with
$|S\setminus[1,\ell+k-2]|\le1$.  For $k\ge3$ this is because
$a^m((\ell+1)\rho+ma)\le a^{m-1}(\ell+1)+ma^{m+1}\le\ell+1+m$.  For $k=2$ it
is Corollary~\ref{cor:secondrow}, since $1+\rho\le r\le r\Pi_Y$.
\end{itemize}
In exact arithmetic at the certified optima of Corollary~\ref{cor:secondrow},
$\max_n\rho(Z_0)$ is $1.41$ at $r=\frac52$ ($n\le36$), $1.02$ at $r=3$ ($n\le42$)
and $0.70$ at $r=4$ ($n\le50$).  So $\rho(Z_0)\le r-1$ holds throughout these
ranges.  At $r=\frac94$ it fails for most $n$ (the maximum is $1.65$).  At
$r=2$, $\rho(Z_0)$ lies between $1.27$ and $2$ for $n\le30$, and
$\rho\le r-1$ fails.  A proof of $\rho(Z_0)\le r-1$ for all $n$ would give, for every
$n$, both the second row and the exempted sets of the proposition.  So would
a proof of the weaker $1+\rho(Z_0)\le r\Pi_Y$ for the second row alone.
Both are statements about the top-row optimum only.  As in
Section~\ref{sec:prefixmono}, they seem to need control of the zero
distribution of $Z_0$ that we do not have.

\section{Every row is a top row at roots at least \texorpdfstring{$2$}{2}}\label{sec:allrowstwo}

Sections~\ref{sec:bigroot}--\ref{sec:rootbelowfour} prove that rows equal the
top row by bounding the hole product of a top-row optimum.  That route stops
near $r=3.55$.  This section takes another route and proves the
conjecture stated after Corollary~\ref{cor:belowtwo} at every root $r\ge2$
(Theorem~\ref{thm:allrowstwo}).  The hole product plays no role.  We verify
the hypothesis of Corollary~\ref{cor:crossreduce}, in a truncated form, by a
first-order optimality argument.

The idea is as follows.  Insert at a gap $c$ of a minimizer $W$.  The zeros of $W$
above $c$ move up by one, and above $c$ the new polynomial is essentially
the old one shifted by one position.  Measured against the first moment
$\sum_ss|r_W(s)|r^{-s}$, the shift costs exactly the factor $1/r$.  The
first moment is in turn bounded by optimality, in the direction that trades
the leading coefficient of $r_{\mathrm{ins}(W,c)}$ against that of $s\,r_W$.

For $D\in\N_{>0}\cup\{\infty\}$ and a real polynomial $q$ put
\[
  \Phi_{r,D}(q)=\sum_{s=1}^{D}|q(s)|\,r^{-s},\qquad
  \Phi_{r,D}^{<c}(q)=\sum_{1\le s<c,\ s\le D}|q(s)|\,r^{-s},\qquad
  \Phi_{r,D}^{>c}(q)=\sum_{c<s\le D}|q(s)|\,r^{-s},
\]
and let $\mu_{r,D}(S,L)$ be the infimum of $\Phi_{r,D}(q)$ over real $q$ with
$\deg q\le L-1$, $q(0)=1$ and $q|_S=0$.  This is the truncated quantity of the
proof of Proposition~\ref{prop:rowvalue}(b), and
$\Phi_{r,\infty}=\Phi_r$, $\mu_{r,\infty}=\mu_r$.  For a finite
$W\subset\N_{>0}$ and an integer $c\ge1$ with $c\notin W$ put
$A_c=r_{W\cap[1,c-1]}$ and $B_c=r_{W\cap(c,\infty)}$, so $r_W=A_cB_c$.  Since
$\mathrm{ins}(W,c)=(W\cap[1,c-1])\cup\{c\}\cup\{w+1:w\in W,\ w>c\}$ and
$1-s/(b+1)=(1-(s-1)/b)\,b/(b+1)$,
\[
  r_{\mathrm{ins}(W,c)}(s)=A_c(s)\Bigl(1-\frac sc\Bigr)\frac{B_c(s-1)}{B_c(-1)},
  \qquad B_c(-1)=\prod_{w\in W,\ w>c}\Bigl(1+\frac1w\Bigr).
\]

\begin{lemma}[Vertex minimizers of truncated problems]\label{lem:truncvertex}
Let $r>1$, let $S\subset\N_{>0}$ be finite, let $L>|S|$, and let $D$ be an
integer with $D\ge L+\max(S\cup\{0\})$.  Then some $W$ with
$S\subseteq W\subseteq[1,D]$ and $|W|=L-1$ satisfies
$\Phi_{r,D}(r_W)=\mu_{r,D}(S,L)$.
\end{lemma}

\begin{proof}
This is the argument of Lemma~\ref{lem:intzeros}.  Put $d=L-1-|S|$ and
$N=[1,D]\setminus S$, so $|N|\ge D-|S|\ge L>d$.  Every admissible $q$
factors as $q=r_Sp$ with $\deg p\le d$ and $p(0)=1$.  Moreover
$\Phi_{r,D}(q)=f(p):=\sum_{s\in N}\omega(s)|p(s)|$ with
$\omega(s)=|r_S(s)|r^{-s}>0$ on $N$.  If $d=0$, then $p=1$ and $W=S$.  Let
$d\ge1$.  For each sign vector $(\epsilon_s)_{s\in N}$, the cell
$\{p:\deg p\le d,\ p(0)=1,\ \epsilon_sp(s)\ge0\ (s\in N)\}$ contains no line.
Indeed, a direction $e$ of a line has $e(0)=0$, $\deg e\le d$ and $e|_N=0$,
so $e=0$.  On the cell $f$ is affine and nonnegative, so its minimum there is
attained at a vertex.  A vertex satisfies $d$ linearly independent equations
$p(y)=0$ with $y\in N$.  So $p=\prod_{y\in Y}(1-s/y)$ for a set $Y\subseteq N$
of $d$ integers.  The finitely many cells cover all admissible $p$.  So the
minimum of $f$ is attained at such a $p$, and $W=S\cup Y$ works.
\end{proof}

\begin{lemma}[Inserting at a gap of a minimizer]\label{lem:optins}
Let $r>1$ and $D\in\N_{>0}\cup\{\infty\}$.  Let $T\subseteq W$ be finite
subsets of $\N_{>0}$, put $L=|W|+1$, and suppose that
$\varphi:=\Phi_{r,D}(r_W)=\mu_{r,D}(T,L)$.  Let $c\ge1$ be an integer with
$c\notin W$ and $c>\max T$ ($\max\varnothing=0$), and put
$Q=r_{\mathrm{ins}(W,c)}$.  Then
\[
  \Phi_{r,D}(Q)\le\frac{\varphi+(r-2)\,\Phi_{r,D}^{<c}(Q)}{r-1},
  \qquad
  \Phi_{r,D}^{<c}(Q)\le\Phi_{r,D}^{<c}(r_W)\le\varphi .
\]
In particular $\Phi_{r,D}(Q)\le\max\{1,a\}\,\varphi$, and
$\Phi_{r,D}(Q)\le\varphi$ if $r\ge2$.
\end{lemma}

\begin{proof}
Write $P=r_W$, $A=A_c$, $B=B_c$ and $\kappa=1/(cB(-1))>0$.  The zeros of $A$
lie in $[1,c-1]$, and those of $B$ are integers greater than $c$.  So $B>0$ on
$[-1,c)$, and $B(s-1)B(s)\ge0$ for every integer $s$, since $B$ does not
change sign between consecutive integers.

\emph{Step 1: signs.}  Let $1\le s<c$.  Then $P(s)Q(s)=A(s)^2(1-s/c)B(s)B(s-1)/B(-1)\ge0$.  If
$A(s)\ne0$, then
\[
  \frac{|Q(s)|}{|P(s)|}=\frac{c-s}c\prod_{w\in W,\ w>c}g(w),\qquad
  g(w)=\frac{(w+1-s)\,w}{(w-s)(w+1)}=1+\frac s{(w-s)(w+1)} .
\]
The factor $g$ exceeds $1$ and decreases in $w$.  The $m$ elements of $W$
above $c$ are at least $c+1,\ldots,c+m$, so the product is at most
$\prod_{j=1}^mg(c+j)=\frac{(c+m+1-s)(c+1)}{(c+1-s)(c+m+1)}\le\frac{c+1}{c+1-s}$.
Since $(c-s)(c+1)\le c(c+1-s)$, we get $|Q(s)|\le|P(s)|$.  If $A(s)=0$, then
$P(s)=Q(s)=0$.  So $Q$ has the sign of $P$ below $c$, and
$\Phi_{r,D}^{<c}(Q)\le\Phi_{r,D}^{<c}(P)\le\varphi$.  Next, $Q(c)=0$.  For
$s>c$ we have $P(s)Q(s)=A(s)^2(1-s/c)B(s)B(s-1)/B(-1)\le0$.

\emph{Step 2: the shift.}  Let $s>c$ and $u=s-1\ge c$.  Every factor
$(u+1-z)/(u-z)$ with $z\le c-1$ exceeds $1$, so
\[
  \frac{|A(u+1)|}{|A(u)|}=\prod_{z\in W\cap[1,c-1]}\frac{u+1-z}{u-z}
  \le\prod_{z=1}^{c-1}\frac{u+1-z}{u-z}=\frac u{u+1-c} .
\]
Hence $|A(s)|(s-c)\le(s-1)|A(s-1)|$, and
$|Q(s)|r^{-s}\le\frac\kappa r(s-1)|P(s-1)|r^{-(s-1)}$.  Put
$M=\sum_{1\le u\le D}u|P(u)|r^{-u}$, which is finite, and split it as
$M=M_<+M_\ge$ according to $u<c$ or $u\ge c$.  Summing over $c<s\le D$ gives
\[
  \Phi_{r,D}^{>c}(Q)\le\frac\kappa r\,M_\ge .
\]

\emph{Step 3: optimality.}  The leading coefficient of $Q$ is $-\kappa$ times
that of $P$, and $\deg Q=\deg P+1=L$.  So $E=P-Q-\kappa sP$ has
$\deg E\le L-1$ and $E(0)=0$.  It vanishes on $T$, because $T\subseteq W\cap[1,c-1]$
consists of zeros of $A$.  Hence, for $\varepsilon>0$, the polynomial
\[
  q_\varepsilon=P+\varepsilon E=(1+\varepsilon-\varepsilon\kappa s)P-\varepsilon Q
\]
is admissible for $\mu_{r,D}(T,L)$, and $\Phi_{r,D}(q_\varepsilon)\ge\varphi$.
Let $\varepsilon\le1/(\kappa c)$.  For $1\le s<c$ we have
$1+\varepsilon-\varepsilon\kappa s\ge\varepsilon$, and by Step~1
$|q_\varepsilon(s)|=(1+\varepsilon-\varepsilon\kappa s)|P(s)|-\varepsilon|Q(s)|$.
For $s\ge c$, $|q_\varepsilon(s)|\le|1+\varepsilon-\varepsilon\kappa s||P(s)|+\varepsilon|Q(s)|$.
Moreover
$|1+\varepsilon-\varepsilon\kappa s|\le1+\varepsilon-\varepsilon\kappa s+2\varepsilon\kappa s\,[\varepsilon\kappa s>1]$.
Summing over $s\le D$ gives
\[
  \varphi\le(1+\varepsilon)\varphi-\varepsilon\kappa M
  -\varepsilon\Phi_{r,D}^{<c}(Q)+\varepsilon\Phi_{r,D}^{>c}(Q)
  +2\varepsilon\kappa\!\!\sum_{s>1/(\varepsilon\kappa),\ s\le D}\!\!s|P(s)|r^{-s}.
\]
Divide by $\varepsilon$ and let $\varepsilon\to0$.  The last sum tends to $0$
(it is empty for small $\varepsilon$ when $D<\infty$), so
\[
  \kappa M+\Phi_{r,D}^{<c}(Q)\le\varphi+\Phi_{r,D}^{>c}(Q).
\]

\emph{Step 4.}  By Step~2 the last inequality gives
$\kappa M_<+\kappa M_\ge(1-1/r)\le\varphi-\Phi_{r,D}^{<c}(Q)$.  So
$\frac\kappa rM_\ge\le(\varphi-\Phi_{r,D}^{<c}(Q))/(r-1)$.  Since $Q(c)=0$,
\[
  \Phi_{r,D}(Q)=\Phi_{r,D}^{<c}(Q)+\Phi_{r,D}^{>c}(Q)
  \le\Phi_{r,D}^{<c}(Q)+\frac{\varphi-\Phi_{r,D}^{<c}(Q)}{r-1}
  =\frac{\varphi+(r-2)\Phi_{r,D}^{<c}(Q)}{r-1}.
\]
For $r\ge2$ use $\Phi_{r,D}^{<c}(Q)\le\varphi$.  For $r<2$ drop the
nonpositive term.
\end{proof}

Only the optimality of $W$ in the single direction $E$ is used, and the
lemma fails without it, already at the first gap.  For
$Z=\{1,2,4,5,7,8,11,13\}$ at $r=2$ an exact computation gives
$\Phi_2(\mathrm{ins}(Z,3))>\frac75\Phi_2(Z)$.  The example after
Proposition~\ref{prop:insertion} at $r=\frac{23}{10}$ also inserts at the first gap.

\begin{theorem}[All rows at every root at least $2$]\label{thm:allrowstwo}
Let $r>1$ and $n\ge1$.  For every finite $S\subset\N_{>0}$,
\[
  \mu_r\bigl(S,|S|+n\bigr)\le\frac{\max\{1,a\}^{|S|}}{\beta_r(n)} .
\]
Consequently $V_r(L,k)\ge\min\{1,r-1\}^{k-1}\beta_r(L-k+1)$ for all
$1\le k\le L$.  In particular, for every $r\ge2$, every $n\ge1$ and every
$L\ge n$,
\[
  V_r(L,L-n+1)=\beta_r(n).
\]
\end{theorem}

\begin{proof}
Put $\lambda=\max\{1,a\}$.  We show by induction on $|S|$ that
\[
  \mu_{r,D}\bigl(S,|S|+n\bigr)\le\lambda^{|S|}\mu_{r,D}(\varnothing,n)
  \qquad\text{for every integer }D\ge|S|+n+\max(S\cup\{0\}).
\]
For $S=\varnothing$ there is nothing to show.  Let $S\ne\varnothing$, and put
$\sigma=\max S$, $T=S\setminus\{\sigma\}$ and $t=\max(T\cup\{0\})$.  Then $D$
is admissible for $T$ as well.  By Lemma~\ref{lem:truncvertex}, some $W$ with
$T\subseteq W\subseteq[1,D]$ and $|W|=|T|+n-1$ has
$\Phi_{r,D}(r_W)=\mu_{r,D}(T,|T|+n)$.  Let $h=h_t(W)$ and
$W'=\mathrm{ins}(W,h)$.  The proof of Theorem~\ref{thm:crossing} uses only
the pointwise inequality $|q(s)|\le\lambda'|q_1(s)|+(1-\lambda')|q_2(s)|$ of
Lemma~\ref{lem:crossing}(a).  So it applies to the truncated sums, and gives
\[
  \mu_{r,D}\bigl(T\cup\{\sigma\},|W|+2\bigr)\le
  \max\bigl\{\Phi_{r,D}(r_W),\ \Phi_{r,D}(r_{W'})\bigr\}.
\]
Here $|W|+2=|S|+n$.  Lemma~\ref{lem:optins} applies with this $T$ and $W$
and with $c=h$, since $h\notin W$ and $h>t$.  It gives
$\Phi_{r,D}(r_{W'})\le\lambda\,\Phi_{r,D}(r_W)$.  Hence
$\mu_{r,D}(S,|S|+n)\le\lambda\,\mu_{r,D}(T,|T|+n)$, and the induction
hypothesis concludes.

Truncated sums are at most full ones, so
$\mu_{r,D}(\varnothing,n)\le\mu_r(\varnothing,n)=1/\beta_r(n)$.  The proof of
Proposition~\ref{prop:rowvalue}(b) shows that $\mu_{r,D}(S,L)\to\mu_r(S,L)$
as $D\to\infty$.  This gives the first claim.  Now let $1\le k\le L$ and
$n=L-k+1$.  Proposition~\ref{prop:rowvalue}(b) gives
$V_r(L,k)=\inf_{|S|=k-1}1/\mu_r(S,L)\ge\lambda^{-(k-1)}\beta_r(n)$, and
$\lambda^{-1}=\min\{1,r-1\}$.  For $r\ge2$ this is $\beta_r(n)$, and
Proposition~\ref{prop:rowvalue}(c) gives the reverse inequality.
\end{proof}

By Corollary~\ref{cor:prefixmono}(a), the theorem proves the conjecture stated
after Corollary~\ref{cor:belowtwo} at every root $r\ge2$.  So the rows at a
single root $r\ge2$ are
\[
  V_r(L,k)=\beta_r(L-k+1)\qquad(1\le k\le L),
\]
the top row of a multiple of $(x-r)^{L-k+1}$.  At $r=2$ this sharpens
\cite[Theorem~4.16]{coefficient-mass} to $b_k(F)\ge\beta_2(L-k+1)$, as anticipated in
the paragraph on rows equal to the top row.  The theorem contains
Theorems~\ref{thm:bigroot}, \ref{thm:rootfour} and~\ref{thm:rootbelowfour}.
It also contains every identity $V_r(L,L-n+1)=\beta_r(n)$ with $r\ge2$ in
Corollaries~\ref{cor:rowstopranges}, \ref{cor:rowspiranges},
\ref{cor:rowsfiniteranges}, \ref{cor:rowstreevals}, \ref{cor:crosstreevals},
\ref{cor:bigrootfinite} and~\ref{cor:rootfourfinite}, and in the ranges
listed after Corollary~\ref{cor:secondrow}.  What remains open at a single root $r\ge2$ is
the value of $\beta_r(n)$ itself.  It is known in closed form where
Theorem~\ref{thm:rowstop} applies, and otherwise through optima certified by
Proposition~\ref{prop:toprowcert}.

For $1<r<2$ the theorem gives $V_r(L,k)\ge(r-1)^{k-1}\beta_r(L-k+1)$.  By
Corollary~\ref{cor:lastrowr} this is an equality for $k=L$, where both sides
are $(r-1)^L$.  The conjectured value $\min_{i\le k}1/\nu_i(n)$ is at least
$(r-1)^{k-1}\beta_r(n)$ by Theorem~\ref{thm:prefixmono}.
Section~\ref{sec:onepoly} proves the conjecture for $1<r<2$ whenever
$\beta_r(n)\le1$, and Section~\ref{sec:longdiag} whenever $k$ is large
compared with $n$.

Lemma~\ref{lem:optins} with $D=\infty$ and $T=\varnothing$ applies to every
top-row optimum $Z_0$ (Lemma~\ref{lem:topzero}(a)) and every hole $c$ of
$Z_0$.  For $r\ge2$ it gives $\Phi_r(\mathrm{ins}(Z_0,c))\le\Phi_r(Z_0)$ at
every hole, whatever $\Pi(Z_0)$ is.  At the first hole $c=\ell+1$ we have
$Z_0\supseteq[1,\ell]$, so $\Phi_r^{<c}(r_{\mathrm{ins}(Z_0,c)})=0$ and
$\Phi_r(\mathrm{ins}(Z_0,\ell+1))\le\Phi_r(Z_0)/(r-1)$.  In the notation of
Corollary~\ref{cor:secondrow} this reads
\[
  \frac{1+\rho}{r\,\Pi_Y}\le\frac1{r-1}
\]
for every $r>1$.  The ratios reported there satisfy this bound.  At $r=2$,
$n=2$ the ratio is $1$, so the bound is attained.  At $r=3$ the ratio is at
most $0.47<\frac12$ for $4\le n\le42$.  The same lemma also recovers
Theorem~\ref{thm:prefixmono}.  Take $T=[1,i-1]$, a set $W$ given by
Lemma~\ref{lem:truncvertex} for $\mu_{r,D}([1,i-1],i-1+m)$, and $c=h_{i-1}(W)$.  Then
$W\supseteq[1,c-1]$, so $\Phi_{r,D}^{<c}(r_{\mathrm{ins}(W,c)})=0$.  Also
$\mathrm{ins}(W,c)\supseteq[1,i]$.  So
$\mu_{r,D}([1,i],i+m)\le a\,\mu_{r,D}([1,i-1],i-1+m)$, and letting
$D\to\infty$ gives $\nu_{i+1}(m)\le a\,\nu_i(m)$.

\section{Rows below the root \texorpdfstring{$2$}{2}}\label{sec:belowtwo}

For $1<r<2$ the rows are not top rows, and Theorem~\ref{thm:allrowstwo}
gives only $V_r(L,k)\ge(r-1)^{k-1}\beta_r(n)$.  The conjectured value is
$\min_{i\le k}1/\nu_i(n)$ (the conjecture stated after
Corollary~\ref{cor:belowtwo}).  This section proves it in two regimes: for
every $k$ when $\beta_r(n)\le1$ (Section~\ref{sec:onepoly}), and for every
$n$ once $k$ is large compared with $n$ (Section~\ref{sec:longdiag}).  In
both the minimum is at $i=k$.

\subsection{One polynomial for all exempted sets}\label{sec:onepoly}

This subsection proves the conjecture for $1<r<2$ whenever $\beta_r(n)\le1$, for every $k$
(Theorem~\ref{thm:onepolyrows}).  For each $n$ this is an interval of roots
$1<r\le r_n$, and $r_n\ge1+\frac1n$ for $n\le10$
(Corollary~\ref{cor:onepolyranges}).  There the minimum is at $i=k$: the
prefix $[1,k-1]$ is the worst exempted set.

Two facts give this.  First, Lemma~\ref{lem:crossing}(b) pushes every
exempted position down into an initial block while a fixed polynomial factor
rides along.  So one polynomial $p$ serves all exempted sets at once, at the
cost of the largest of its prefix sums (Theorem~\ref{thm:onepoly}).  Second,
these prefix sums grow at least by the factor $\nu_1(n)=1/\beta_r(n)$ from
one index to the next (Lemma~\ref{lem:prefixshift}).  Nothing is inserted,
and the only optimality used is that of $p$ for the last index.

Keep $\omega_i(s)=\binom{s-1}{i-1}r^{-s}$ from the paragraph on dilating
and crossing, let $\mathcal Q_n$ be the set of real polynomials $p$ with
$\deg p\le n-1$ and $p(0)=1$, and put
\[
  A_i(p)=\sum_{s\ge1}\omega_i(s)\,|p(s)|=\Phi_r\bigl(r_{[1,i-1]}\,p\bigr)
  \qquad(i\ge1),
\]
for real polynomials $p$, so that $\nu_i(n)=\inf_{p\in\mathcal Q_n}A_i(p)$.
The second expression holds because $|r_{[1,i-1]}(s)|=\binom{s-1}{i-1}$
for $s\ge1$.  For finite $S\subset\N_{>0}$ let $\ell(S)$ be the largest
$\ell\ge0$ with $[1,\ell]\subseteq S$.

\begin{lemma}[Pushing exempted positions into the prefix]\label{lem:prefixpush}
Let $r>1$, let $p$ be a real polynomial and let $S\subset\N_{>0}$ be finite.
Then
\[
  \Phi_r(r_S\,p)\le\max_{1\le i\le|S|+1}A_i(p).
\]
\end{lemma}

\begin{proof}
Induct on $|S|$ and, for fixed $|S|$, on $d(S)=|S|-\ell(S)$.  If $d(S)=0$,
then $S=[1,|S|]$ and $\Phi_r(r_Sp)=A_{|S|+1}(p)$.  Let $d(S)>0$, put
$\ell=\ell(S)$ and $\sigma=\max S$.  Then $\ell+1\notin S$, so
$\sigma\ge\ell+2$; otherwise $S\subseteq[1,\ell]$ and $d(S)=0$.  Put
$U=S\setminus\{\sigma\}$ and $Q=r_Up$, so that $r_Sp=Q\cdot(1-s/\sigma)$.
Lemma~\ref{lem:crossing}(b) with $y=\ell+1\le\sigma$ gives
\[
  \Phi_r(r_Sp)\le\max\bigl\{\Phi_r(r_Up),\ \Phi_r(r_{U\cup\{\ell+1\}}\,p)\bigr\},
\]
because $\ell+1\notin U$.  Here $|U|=|S|-1$, so the induction hypothesis
bounds the first term by $\max_{i\le|S|}A_i(p)$.  The set
$S'=U\cup\{\ell+1\}$ has $|S'|=|S|$ and $[1,\ell+1]\subseteq S'$, so
$d(S')<d(S)$, and the induction hypothesis bounds the second term by
$\max_{i\le|S|+1}A_i(p)$.
\end{proof}

\begin{theorem}[One polynomial for all exempted sets]\label{thm:onepoly}
Let $r>1$, $1\le k\le L$, $n=L-k+1$ and $p\in\mathcal Q_n$.  Every monic
multiple $F$ of $(x-r)^L$ has
\[
  b_k(F)\ \ge\ \frac1{\max_{1\le i\le k}A_i(p)} .
\]
Hence, with $\Lambda_k(n)=\inf_{p\in\mathcal Q_n}\max_{1\le i\le k}A_i(p)$,
\[
  \frac1{\Omega_k(n)}\ \le\ \frac1{\Lambda_k(n)}\ \le\ V_r(L,k)\ \le\
  \min_{1\le i\le k}\frac1{\nu_i(n)} .
\]
\end{theorem}

\begin{proof}
Let $S$ be the exempted set of Proposition~\ref{prop:rowvalue}(a), so
$|S|=k-1$.  The polynomial $q=r_Sp$ has $\deg q\le k-1+n-1=L-1$, $q(0)=1$
and $q|_S=0$.  So $b_k(F)\,\Phi_r(q)\ge1$, and Lemma~\ref{lem:prefixpush}
bounds $\Phi_r(q)$.  The upper bound is Theorem~\ref{thm:prefixrows}.  For
the first inequality, $W_k(s)=\max_{1\le i\le k}\binom{s-1}{i-1}$ in the
notation of Theorem~\ref{thm:prefixrows}, so
$\sum_sW_k(s)r^{-s}|p(s)|\ge\max_{i\le k}A_i(p)$ for every $p$.
\end{proof}

The theorem moves the maximum over $i$ out of the sum in
Theorem~\ref{thm:prefixrows}, and this loses nothing when one $p$ is nearly
optimal for all prefix problems at once.  The conjecture stated after
Corollary~\ref{cor:belowtwo} holds at $(r,n,k)$ whenever
$\Lambda_k(n)=\max_{i\le k}\nu_i(n)$.  For $n=1$ we have
$\mathcal Q_1=\{1\}$ and $A_i(1)=a^i$.  So the theorem gives
$V_r(L,L)\ge r-1$ for $r\ge2$ and $V_r(L,L)\ge(r-1)^L$ for $r\le2$, and
Lemma~\ref{lem:prefixpush} with $p=1$ gives $\Phi_r(S)\le\max\{a,a^{|S|+1}\}$.
These are the lower bounds of Corollary~\ref{cor:lastrowr}, obtained here
from convexity alone, without Lemma~\ref{lem:confdelr} and the tail theorem
behind it.

\begin{lemma}[Shifting the prefix index]\label{lem:prefixshift}
Let $r>1$, $n\ge1$, $i,m\ge1$, and let $p$ be real with $\deg p\le n-1$.
Then
\[
  A_{i+m}(p)\ \ge\ \nu_m(n)\,A_i(p).
\]
Consequently $\nu_{i+m}(n)\ge\nu_i(n)\,\nu_m(n)$, and
$\nu_i(n)\ge\nu_1(n)^i=\beta_r(n)^{-i}$.
\end{lemma}

\begin{proof}
Counting the $(i+m-1)$-element subsets of $[1,s-1]$ by their $i$-th
smallest element $u$ gives
$\sum_{u=1}^{s-1}\binom{u-1}{i-1}\binom{s-u-1}{m-1}=\binom{s-1}{i+m-1}$, that is,
$\omega_{i+m}(s)=\sum_{u+t=s,\ u,t\ge1}\omega_i(u)\,\omega_m(t)$.  All terms
being nonnegative,
\[
  A_{i+m}(p)=\sum_{u\ge1}\omega_i(u)\sum_{t\ge1}\omega_m(t)\,|p(u+t)| .
\]
If $p(u)\ne0$, the polynomial $t\mapsto p(u+t)/p(u)$ lies in $\mathcal Q_n$,
so the inner sum is at least $\nu_m(n)|p(u)|$; if $p(u)=0$ this is trivial.
This proves the inequality.  Taking the infimum over $p\in\mathcal Q_n$
gives $\nu_{i+m}(n)\ge\nu_m(n)\nu_i(n)$, and $\nu_1(n)=1/\beta_r(n)$.
\end{proof}

Together with Theorem~\ref{thm:prefixmono},
$\nu_1(n)\,\nu_i(n)\le\nu_{i+1}(n)\le a\,\nu_i(n)$, and both inequalities are
equalities at $n=1$, where $\nu_i(1)=a^i$.

\begin{theorem}[All rows when the top row is at most one]\label{thm:onepolyrows}
Let $r>1$ and $n\ge1$ with $\beta_r(n)\le1$.  Then
$1\le\nu_1(n)\le\nu_2(n)\le\cdots$, and for every $k\ge1$ and $L=n+k-1$,
\[
  V_r(L,k)=\frac1{\nu_k(n)}=\min_{1\le i\le k}\frac1{\nu_i(n)} .
\]
So the conjecture stated after Corollary~\ref{cor:belowtwo} holds at $(r,n)$
for every $k$, with the minimum at $i=k$.
\end{theorem}

\begin{proof}
By Lemma~\ref{lem:prefixshift} with $m=1$ and $\nu_1(n)=1/\beta_r(n)\ge1$,
$A_{i+1}(p)\ge A_i(p)$ for every $p\in\mathcal Q_n$ and every $i$, and
$\nu_{i+1}(n)\ge\nu_i(n)$.  Given $\varepsilon>0$ choose $p\in\mathcal Q_n$
with $A_k(p)\le\nu_k(n)+\varepsilon$.  Then
$\max_{i\le k}A_i(p)=A_k(p)$, and Theorem~\ref{thm:onepoly} gives
$V_r(L,k)\ge1/(\nu_k(n)+\varepsilon)$.  The upper bound
$V_r(L,k)\le\min_i1/\nu_i(n)\le1/\nu_k(n)$ is Theorem~\ref{thm:prefixrows}.
\end{proof}

Since $\nu_1(n+1)\le\nu_1(n)$, we have $\beta_r(n)\ge\beta_r(1)=r-1$.  So the
hypothesis forces $r\le2$, and at $r=2$ it holds only for $n=1$.  For each
$q\in\mathcal Q_n$ the sum $\Phi_r(q)$ decreases in $r$, as in
Lemma~\ref{lem:rmono}, so $\nu_1(n)$ does not increase with $r$.  Hence
$\{r>1:\beta_r(n)\le1\}$ is an interval with left end $1$, it shrinks as $n$
grows, and for $n=1$ it is $(1,2]$.  Let $r_n$ be its right end.

\begin{corollary}[Explicit ranges]\label{cor:onepolyranges}
Let $1\le n\le10$ and $1<r\le1+\frac1n$.  Then
$V_r(L,L-n+1)=1/\nu_{L-n+1}(n)$ for every $L\ge n$.
\end{corollary}

\begin{proof}
For $n=1$, $\beta_r(1)=r-1\le1$.  For $2\le n\le10$ put $\rho=1+\frac1n$.
Proposition~\ref{prop:toprowcert} with the set $Z_0$ of
Table~\ref{tab:onepolyranges} gives $\beta_\rho(n)=1/\Phi_\rho(Z_0)$, and
exact arithmetic gives $\Phi_\rho(Z_0)>1$.  By the monotonicity just noted,
$\beta_r(n)\le\beta_\rho(n)<1$ for $1<r\le\rho$, and
Theorem~\ref{thm:onepolyrows} applies.
\end{proof}

\begin{table}[ht]
\centering
\small
\begin{tabular}{rll}
$n$ & $Z_0$ & $\beta_{1+1/n}(n)$\\\hline
2 & $\{4\}$ & $0.9153$\\
3 & $\{4,11\}$ & $0.8767$\\
4 & $\{4,11,21\}$ & $0.8526$\\
5 & $\{4,10,20,35\}$ & $0.8380$\\
6 & $\{4,10,19,33,52\}$ & $0.8271$\\
7 & $\{4,10,19,32,48,73\}$ & $0.8195$\\
8 & $\{4,10,19,31,47,68,97\}$ & $0.8135$\\
9 & $\{4,10,19,31,46,65,90,125\}$ & $0.8084$\\
10 & $\{4,10,19,30,45,64,87,116,156\}$ & $0.8042$
\end{tabular}
\caption{Top-row optima at $r=1+\frac1n$, certified by
Proposition~\ref{prop:toprowcert} in exact arithmetic (values rounded).}
\label{tab:onepolyranges}
\end{table}

In floating point the right ends are $r_2\approx1.5405$, $r_3\approx1.3738$,
$r_4\approx1.2865$, $r_5\approx1.2321$, $r_6\approx1.1951$,
$r_7\approx1.1685$ and $r_8\approx1.1482$, so $n(r_n-1)$ grows from $1.08$
to $1.19$.  The values in Table~\ref{tab:onepolyranges} decrease with $n$,
and we expect $\beta_{1+1/n}(n)<1$ for every $n$, but have no proof.  The
corollary agrees with Corollary~\ref{cor:rowstreevals}, where every case
with $n\ge2$ and $r\le1+\frac1n$ has $i^*=k$.  It extends that corollary: at
$r=\frac{11}{10}$ it gives $i^*=k$ for all $k$ and $n\le10$, at
$r=\frac54$ for $n\le4$, and at $r=\frac32$ for $n\le2$.

Theorem~\ref{thm:onepoly} also works beyond $\beta_r(n)\le1$, one row at a
time and with a finite check in place of a tree.

\begin{corollary}[Rows from one certified set]\label{cor:onepolyexact}
Let $r>1$, $n\ge1$, $k\ge2$ and $L=n+k-1$.  Let $Z^*=[1,k-1]\cup Y$ with
$|Y|=n-1$ satisfy the hypothesis of Proposition~\ref{prop:toprowcertS} with
$S=[1,k-1]$.  If $\Phi_r([1,i-1]\cup Y)\le\Phi_r(Z^*)$ for $1\le i<k$, then
\[
  V_r(L,k)=\frac1{\Phi_r(Z^*)}=\frac1{\nu_k(n)}=\min_{1\le i\le k}\frac1{\nu_i(n)} .
\]
This holds in exact arithmetic for $r=\frac32$ with $n=3$, $2\le k\le10$, and
with $n=4$, $5\le k\le10$; for $r=\frac54$ with $n=5$, $2\le k\le10$; and for
$r=\frac74$ with $n=2$, $5\le k\le12$.
\end{corollary}

\begin{proof}
Proposition~\ref{prop:toprowcertS} gives
$\nu_k(n)=\mu_r([1,k-1],L)=\Phi_r(Z^*)$ and $V_r(L,k)\le1/\Phi_r(Z^*)$.
Since $Z^*$ is a set, $Y\subset[k,\infty)$, so $p=r_Y\in\mathcal Q_n$ has
$A_i(p)=\Phi_r([1,i-1]\cup Y)$ for $i\le k$, and $A_k(p)=\Phi_r(Z^*)$.
Theorem~\ref{thm:onepoly} gives $V_r(L,k)\ge1/\Phi_r(Z^*)$, and
Theorem~\ref{thm:prefixrows} gives the rest.  In the listed cases $Y$ is,
for $k=2,3,\ldots$ in turn:
\begingroup\sloppy
\begin{itemize}
\item $r=\frac32$, $n=3$: $\{5,11\}$, $\{8,15\}$, $\{11,18\}$, $\{13,22\}$,
$\{16,25\}$, $\{19,28\}$, $\{21,31\}$, $\{24,35\}$, $\{27,38\}$;
\item $r=\frac32$, $n=4$, from $k=5$: $\{12,18,26\}$, $\{14,21,30\}$,
$\{17,24,33\}$, $\{19,27,37\}$, $\{22,30,40\}$, $\{24,33,44\}$;
\item $r=\frac54$, $n=5$: $\{6,13,22,35\}$, $\{10,17,27,42\}$,
$\{13,22,33,48\}$, $\{17,26,38,55\}$, $\{21,31,44,61\}$, $\{24,36,49,67\}$,
$\{28,40,55,74\}$, $\{32,45,60,80\}$, $\{36,50,65,86\}$;
\item $r=\frac74$, $n=2$, from $k=5$: $\{12\}$, $\{15\}$, $\{17\}$, $\{19\}$,
$\{22\}$, $\{24\}$, $\{26\}$, $\{29\}$.
\end{itemize}
\endgroup
The certificates of Proposition~\ref{prop:toprowcertS} and the inequalities
$\Phi_r([1,i-1]\cup Y)\le\Phi_r(Z^*)$ were checked in exact arithmetic.
\end{proof}

Here $i^*=k$ throughout.  Corollary~\ref{cor:rowstreevals} has these rows
only for $k\le4$ ($r=\frac32$ and $r=\frac74$) and for $k\le3$ ($r=\frac54$, $n=5$).  The
check fails with the certified set at $r=\frac32$, $n=4$, $k\le4$ and at
$r=\frac74$, $n=2$, $2\le k\le4$, although $i^*=k$ there as well.  Where
$\max_i\nu_i(n)$ sits at $i=1<k$, the theorem cannot give the conjecture in
general, because its bound covers real exempted positions as well.
For $k=2$ and any $p\in\mathcal Q_n$, Lemma~\ref{lem:crossing}(b) gives
$\Phi_r(p\cdot(1-s/x))\le\max\{A_1(p),A_2(p)\}$ for every real $x\ge1$, so
$\Lambda_2(n)$ is at least the infimum of $\Phi_r(q)$ over $q$ with
$\deg q\le n$, $q(0)=1$ and $q(x)=0$.  At $r=\frac74$, $n=6$ and
$x=\frac{16}9$ this infimum is at least $0.10845$.  The certificate is a set
of weights $z_s\in[-1,1]$, $1\le s\le40$, with
$\sum_sz_s|1-s/x|r^{-s}s^j=0$ for $1\le j\le5$ and
$\sum_sz_s|1-s/x|r^{-s}\ge0.10845$.  For $q=p\cdot(1-s/x)$ with
$p\in\mathcal Q_6$ it gives
$\Phi_r(q)\ge\sum_{s\le40}z_s|1-s/x|r^{-s}p(s)=\sum_sz_s|1-s/x|r^{-s}$.
But $\nu_1(6)\le\Phi_r(\{1,3,6,9,15\})<0.09940$ and
$\nu_2(6)\le\Phi_r(\{1,3,5,8,12,18\})<0.05975$.  So
$\Lambda_2(6)>1.09\max\{\nu_1(6),\nu_2(6)\}$ there, while
Corollary~\ref{cor:rowstreevals}(c) gives $V_{7/4}(7,2)=\beta_{7/4}(6)$.

Two other routes fail below $2$.  First, Lemma~\ref{lem:optins} does not
carry over.  At $r=\frac74$ and $n=5$ the top-row optimum is
$Z_0=\{1,3,7,11\}$, and it is the only polynomial in $\mathcal Q_5$ with
$\Phi_r\le\nu_1(5)$: the weights $z_s=\sgn r_{Z_0}(s)$ for $s\notin Z_0$,
completed on $Z_0$ so that $\sum_sz_s\,s^j r^{-s}=0$ for $1\le j\le4$, have
$|z_s|<1$ on $Z_0$.  So every $q\in\mathcal Q_5$ has
$\Phi_r(q)\ge\sum_sz_sq(s)r^{-s}=\Phi_r(Z_0)$, with equality only if $q$
vanishes on $Z_0$.  Moreover $\nu_2(5)\le\Phi_r(\{1,3,6,9,15\})<\nu_1(5)$,
yet $\Phi_r(\mathrm{ins}(Z_0,2))>1.17\,\Phi_r(Z_0)$.  So
Theorem~\ref{thm:crossing} with $k=2$ gives no bound as large as
$\beta_{7/4}(5)=V_{7/4}(6,2)$ from any set $W$.  Second, the argument of
Lemma~\ref{lem:prefixpush} would prove the conjecture if each push kept the
best polynomial, that is, if
$\mu_r(U\cup\{x\},|U|+n+1)\le\max\{\mu_r(U,|U|+n),\,\mu_r(U\cup\{h\},|U|+n+1)\}$
for $x>\max U$ and $h$ the first gap of $U$.  This fails: at $r=\frac32$,
$n=6$, $U=\{3,4\}$ and $x=5$ the left side is at least $0.17071$ (a dual
certificate on $s\le80$), and the right side is at most $0.16846$, by
the sets $\{1,3,4,9,13,20,28\}$ and $\{1,3,4,8,12,17,24,33\}$.  The
conjecture itself is not affected there: it only asks for
$\max_{i\le4}\nu_i(6)$ on the right, and in floating point
$\nu_1(6)\approx0.2571$.
For $1<r<2$ the conjecture remains open where $\beta_r(n)>1$ and no single
polynomial serves all prefixes, in particular wherever the maximum of
$\nu_i(n)$ sits at $i=1<k$.  Section~\ref{sec:longdiag} shows that a single
polynomial does serve all prefixes once $k$ is large compared with $n$.

\subsection{Long diagonals}\label{sec:longdiag}

Theorem~\ref{thm:onepoly} gives the conjectured value $1/\nu_k(n)$ as soon
as one polynomial that is nearly optimal for $\nu_k(n)$ has its $k$-th
prefix sum $A_k$ at least as large as the earlier ones.
Theorem~\ref{thm:onepolyrows} obtains this from $\beta_r(n)\le1$.  This
subsection obtains it for every $1<r<2$, with no condition on $\beta_r(n)$,
once $k$ is large compared with $n$ (Theorem~\ref{thm:longdiag} and
Proposition~\ref{prop:longdiagexplicit}).  So on every diagonal
$L-k=n-1$ the conjecture can fail in at most finitely many rows
(Corollary~\ref{cor:longdiagfinite}).  Together with the exact certificates
of Section~\ref{sec:onepoly}, this settles whole diagonals at
$r=\frac54,\frac32,\frac74$ (Corollary~\ref{cor:completediag}).

The mechanism is simple.  By Lemma~\ref{lem:intzeros} the polynomials for
$\nu_k(n)$ may be taken with all zeros at least $k$, and such a polynomial is
bounded by $1$ on $[1,2k-3]$.  There the weights $\omega_i$, $i<k$, exceed
$\omega_k$ by an amount that is exponentially small compared with $\nu_k(n)$
when $r<2$.  From $2k-2$ on, $\omega_k$ dominates every $\omega_i$, $i<k$,
with a margin that is a fixed fraction of $\omega_k$ near the mode of
$\omega_k$, where every polynomial in $\mathcal Q_n$ has some mass.

For $k\ge2$ put
\[
  w_k(s)=\omega_k(s)-\omega_{k-1}(s),\qquad
  \tau_k(n)=\inf_{p\in\mathcal Q_n}\sum_{s\ge2k-2}w_k(s)\,|p(s)|,
\]
\[
  N_k(r)=\max_{1\le i<k}\sum_{s=1}^{2k-3}\bigl(\omega_i(s)-\omega_k(s)\bigr)^+,
\]
where $x^+=\max\{x,0\}$ and $\omega_i(s)=0$ for $s<i$.  For $s\ge k$ we have
$\omega_{k-1}(s)/\omega_k(s)=(k-1)/(s-k+1)$, so
\begin{equation}\label{eq:wkratio}
  w_k(s)=\frac{s-2k+2}{s-k+1}\,\omega_k(s)\qquad(s\ge k).
\end{equation}
In particular $w_k\ge0$ on $[2k-2,\infty)$, and $\tau_k(n)\ge0$.

\begin{lemma}[Far zeros keep the last prefix sum largest]\label{lem:farprefix}
Let $r>1$, $k\ge2$ and $n\ge1$.  Let $Y$ be a set of $n-1$ integers, each at
least $k$, and $p=r_Y$.  Then for $1\le i<k$,
\[
  A_k(p)-A_i(p)\ \ge\ \tau_k(n)-N_k(r).
\]
\end{lemma}

\begin{proof}
Let $1\le s\le2k-3$.  Every $y\in Y$ has $2y\ge2k>s$, so $|1-s/y|\le1$ and
$|p(s)|\le1$.  Hence
$(\omega_k(s)-\omega_i(s))|p(s)|\ge-(\omega_i(s)-\omega_k(s))^+$, and the
sum of these terms over $s\le2k-3$ is at least $-N_k(r)$.

Let $s\ge2k-2$.  The coefficients $\binom{s-1}j$ increase in $j$ for
$j\le(s-1)/2$, and $i-1\le k-2\le(s-1)/2$, so $\omega_i(s)\le\omega_{k-1}(s)$
and $(\omega_k(s)-\omega_i(s))|p(s)|\ge w_k(s)|p(s)|$.  Since $p\in\mathcal Q_n$,
the sum of these terms over $s\ge2k-2$ is at least $\tau_k(n)$.
\end{proof}

\begin{theorem}[Long diagonals]\label{thm:longdiag}
Let $r>1$, $n\ge1$, $k\ge2$ and $L=n+k-1$.  If $N_k(r)\le\tau_k(n)$, then
$\nu_i(n)\le\nu_k(n)$ for $1\le i\le k$, and
\[
  V_r(L,k)=\frac1{\nu_k(n)}=\min_{1\le i\le k}\frac1{\nu_i(n)} .
\]
So the conjecture stated after Corollary~\ref{cor:belowtwo} holds at
$(r,L,k)$, with the minimum at $i=k$.
\end{theorem}

\begin{proof}
Let $\varepsilon>0$.  Lemma~\ref{lem:intzeros} with $\omega=\omega_k$ gives a
set $Y$ of $n-1$ integers, each at least $k$, such that $p=r_Y$ has
$A_k(p)\le\nu_k(n)+\varepsilon$.  By Lemma~\ref{lem:farprefix} and the
hypothesis, $A_i(p)\le A_k(p)$ for $i<k$.  Hence
$\nu_i(n)\le A_i(p)\le\nu_k(n)+\varepsilon$ for $i\le k$, and
$\max_{i\le k}A_i(p)=A_k(p)$.  Theorem~\ref{thm:onepoly} gives
$V_r(L,k)\ge1/(\nu_k(n)+\varepsilon)$.  Let $\varepsilon\to0$.
Theorem~\ref{thm:prefixrows} gives
$V_r(L,k)\le\min_i1/\nu_i(n)\le1/\nu_k(n)$.
\end{proof}

The hypothesis compares a quantity that is exponentially small for $r<2$
with one that is not.  The next proposition makes this explicit.  It uses
the mode of $\omega_k$ and Lagrange interpolation, as in
Lemma~\ref{lem:nbmass}.

\begin{proposition}[An explicit form]\label{prop:longdiagexplicit}
Let $1<r<2$, $\theta=4(r-1)/r^2$ and $n\ge1$.  For $k\ge2$ put
$M_k=\lfloor(k-1)r/(r-1)\rfloor+1$ and
\[
  F(k)=\theta^k\,(2r)^{n-1}\binom{M_k+n-1}{n-1}
  \Bigl(\frac{4\sqrt{kr}}{r-1}+1\Bigr).
\]
\begin{enumerate}
\item[(a)] If $F(k)\le3(2-r)^2/r^2$, then $N_k(r)\le\tau_k(n)$, so
Theorem~\ref{thm:longdiag} applies.
\item[(b)] Put $d=\lceil r/(r-1)\rceil$ and
$\rho_k=\theta\,(1+(n-1)/(M_k+1))^d\sqrt{1+1/k}$.  Then $F(k+1)\le\rho_kF(k)$,
and $\rho_k$ does not increase with $k$.  So if $F(k_1)\le3(2-r)^2/r^2$ and
$\rho_{k_1}\le1$, then (a) applies to every $k\ge k_1$.
\item[(c)] Let $\alpha\ge1$ satisfy
$\rho(\alpha):=\theta\,\bigl(2re(\alpha r/(r-1)+1)\bigr)^{1/\alpha}<1$.  If
$n-1\le k/\alpha$, then
$F(k)\le\rho(\alpha)^k(4\sqrt{kr}/(r-1)+1)$.  Consequently there is $K$ such
that $V_r(n+k-1,k)=1/\nu_k(n)$ for every $n\ge1$ and every
$k\ge\max\{K,\alpha(n-1)\}$.
\end{enumerate}
Since $\bigl(2re(\alpha r/(r-1)+1)\bigr)^{1/\alpha}\to1$ as $\alpha\to\infty$
and $\theta<1$, admissible $\alpha$ exist for every $1<r<2$.  For instance,
$(\alpha,K)=(30,21)$ works at $r=\frac54$, $(100,149)$ at $r=\frac32$ and
$(1000,761)$ at $r=\frac74$.
\end{proposition}

\begin{proof}
(a) First, $(\omega_i(s)-\omega_k(s))^+\le\omega_i(s)\le
\binom{s-1}{\lfloor(s-1)/2\rfloor}r^{-s}$, so $N_k(r)\le E_k(r)$, and the
proof of Corollary~\ref{cor:belowtwo} gives
$E_k(r)\le\frac{r^2}{4(2-r)}(4/r^2)^k=\frac{r^2}{4(2-r)}\,\theta^ka^k$.

Next let $p\in\mathcal Q_n$.  The ratio $\omega_k(s+1)/\omega_k(s)=s/(r(s-k+1))$
is at least $1/r$, and it is at least $1$ exactly when $s\le(k-1)r/(r-1)$.  So
$M_k$ maximizes $\omega_k$: it is a mode of the law $P$ of
Lemma~\ref{lem:nbmass}, where $\omega_k=a^kP$.  That lemma gives
$\omega_k(M_k)\ge a^k\cdot3/(4(4\sqrt{kr}/(r-1)+1))$, and
$\omega_k(M_k+l)\ge r^{-l}\omega_k(M_k)$ for $l\ge0$.  Since
$M_k>(k-1)r/(r-1)\ge2k-2$, the nodes $t_l=M_k+l$, $0\le l<n$, lie in
$[2k-2,\infty)$.  By~\eqref{eq:wkratio} the ratio $w_k(s)/\omega_k(s)=1-(k-1)/(s-k+1)$
increases in $s$, and at $s=t_0>(k-1)r/(r-1)$ it exceeds $1-(r-1)=2-r$.  So
$w_k(t_l)\ge(2-r)\,r^{-(n-1)}\omega_k(M_k)$ for every $l$.  Finally
$1=p(0)=\sum_l\ell_lp(t_l)$ with the Lagrange coefficients
$|\ell_l|=\prod_{l'\ne l}t_{l'}/|l-l'|\le\prod_{j=1}^{n-1}(M_k+j)/(l!\,(n-1-l)!)$,
so $\sum_l|\ell_l|\le2^{n-1}\binom{M_k+n-1}{n-1}$ and
$\max_l|p(t_l)|\ge1/(2^{n-1}\binom{M_k+n-1}{n-1})$.  Altogether
\[
  \sum_{s\ge2k-2}w_k(s)|p(s)|\ge\sum_lw_k(t_l)|p(t_l)|
  \ge\frac{3(2-r)\,a^k}{4\,(4\sqrt{kr}/(r-1)+1)\,(2r)^{n-1}\binom{M_k+n-1}{n-1}} .
\]
This lower bound for $\tau_k(n)$ is at least the upper bound for $E_k(r)$
exactly when $F(k)\le3(2-r)^2/r^2$.

(b) We have $M_{k+1}-M_k\le d$, so
$\binom{M_{k+1}+n-1}{n-1}/\binom{M_k+n-1}{n-1}=\prod_{j=M_k+1}^{M_{k+1}}(j+n-1)/j
\le(1+(n-1)/(M_k+1))^d$.  Also
$(4\sqrt{(k+1)r}/(r-1)+1)/(4\sqrt{kr}/(r-1)+1)\le\sqrt{(k+1)/k}$.  This gives
$F(k+1)\le\rho_kF(k)$.  Since $M_k$ does not decrease, neither does
$1/\rho_k$, and induction on $k$ gives the rest.

(c) For $n=1$ the claim is $\theta\le\rho(\alpha)$.  Let $m=n-1\ge1$.  Then
$\binom{M_k+m}m\le(e(M_k+m)/m)^m$, and the right side increases with $m$,
since the derivative of $m\log(e(M_k+m)/m)$ is $\log(1+M_k/m)+m/(M_k+m)>0$.
With $m\le k/\alpha$ and $M_k\le(k-1)r/(r-1)+1\le kr/(r-1)$ this gives
$\binom{M_k+m}m\le(e(\alpha r/(r-1)+1))^{k/\alpha}$.  Also
$(2r)^m\le(2r)^{k/\alpha}$, which proves the bound on $F(k)$.  Its right side
tends to $0$ as $k\to\infty$ and decreases once $k\ge1/(2\log(1/\rho(\alpha)))$.
So some $K$ has $\rho(\alpha)^k(4\sqrt{kr}/(r-1)+1)\le3(2-r)^2/r^2$ for all
$k\ge K$, and (a) and Theorem~\ref{thm:longdiag} apply whenever also
$k\ge\alpha(n-1)$.  The listed pairs were checked with $40$-digit
arithmetic.
\end{proof}

\begin{corollary}[Every diagonal is eventually determined]\label{cor:longdiagfinite}
Let $1<r<2$ and $n\ge1$.  Then $V_r(L,L-n+1)=1/\nu_{L-n+1}(n)$ for all but
finitely many $L\ge n$.  Moreover the sequence $\nu_k(n)$ is eventually
nondecreasing in $k$, and $\nu_k(n)=\max_{i\le k}\nu_i(n)$ for all large $k$.
\end{corollary}

\begin{proof}
Apply Proposition~\ref{prop:longdiagexplicit}(c) with $k=L-n+1\ge\max\{K,\alpha(n-1)\}$.
Theorem~\ref{thm:longdiag} also gives $\nu_i(n)\le\nu_k(n)$ for all
$i\le k$ and all such $k$.
\end{proof}

This sharpens Corollary~\ref{cor:belowtwo}, which gives
$V_r(L,k)=(1+o(1))/\nu_k(n)$ as $k\to\infty$ with $n$ fixed: the factor
$1+o(1)$ is exactly $1$ once $k\ge\max\{K,\alpha(n-1)\}$, uniformly in $n$.
For small $n$ the explicit thresholds are much smaller than
$\alpha(n-1)$.  The hypothesis $N_k(r)\le\tau_k(n)$ of
Theorem~\ref{thm:longdiag} can be checked row by row in exact arithmetic.  A
set $Y$ of $n-1$ integers in $[2k-2,D]$ and the weights
$z_s=w_k(s)\sgn r_Y(s)$ for $s\in[2k-2,D]\setminus Y$, completed on $Y$ so
that $\sum_sz_ss^j=0$ for $1\le j\le n-1$, give
$\tau_k(n)\ge\sum_{s=2k-2}^Dz_s$ whenever $|z_y|\le w_k(y)$ on $Y$: for
$p\in\mathcal Q_n$, $\sum_{s\ge2k-2}w_k|p|\ge\sum_{s=2k-2}^Dz_sp(s)=\sum_sz_s$.
We found $Y$ by linear programming and then compared $\sum_sz_s$ with
$E_k(r)$, or with $N_k(r)$ when $E_k(r)$ was too large.

\begin{corollary}[Complete diagonals]\label{cor:completediag}
For every $L\ge n$ and $k=L-n+1$,
\[
  V_r(L,k)=\frac1{\nu_k(n)}=\min_{1\le i\le k}\frac1{\nu_i(n)}
\]
in each of the cases $r=\frac54$ with $n\le9$; $r=\frac32$ with $n\le4$;
and $r=\frac74$ with $n\le2$.  So the conjecture stated after
Corollary~\ref{cor:belowtwo} holds on these diagonals, with the minimum at
$i=k$ in every row.
\end{corollary}

\begin{proof}
The row $k=1$ is $V_r(n,1)=\beta_r(n)=1/\nu_1(n)$.  For $n=1$,
$\nu_i(1)=a^i$ and Corollary~\ref{cor:lastrowr} gives $V_r(L,L)=(r-1)^L$.
For $r=\frac54$ with $n\le4$ and $r=\frac32$ with $n=2$ the claim is
Corollary~\ref{cor:onepolyranges}.  For the remaining diagonals and $k\ge2$
we use four sources:
(I)~Corollary~\ref{cor:rowstreevals}(c),(d), whose rows have $i^*=k$ in
the cases used here;
(II)~single certified sets as in Corollary~\ref{cor:onepolyexact};
(III)~the exact row-by-row check of $N_k(r)\le\tau_k(n)$ just described;
(IV)~Proposition~\ref{prop:longdiagexplicit}(b) from $k_1$ on, with
$F(k_1)\le3(2-r)^2/r^2$ and $\rho_{k_1}\le1$ checked in exact arithmetic.
The rows are covered as follows.
\begingroup\sloppy
\begin{itemize}
\item $r=\frac54$, $n=5$: (III) for $2\le k\le63$, (IV) with $k_1=64$.
\item $r=\frac54$, $n=6$: (I) for $k\le3$, (III) for $4\le k\le77$, (IV)
with $k_1=78$.
\item $r=\frac54$, $n=7$: (I) for $k\le3$, (II) for $k=4,5$, (III) for
$6\le k\le91$, (IV) with $k_1=92$.
\item $r=\frac54$, $n=8$: (I) for $k\le3$, (II) for $4\le k\le7$, (III) for
$8\le k\le105$, (IV) with $k_1=106$.
\item $r=\frac54$, $n=9$: (I) for $k\le3$, (II) for $4\le k\le10$, (III) for
$11\le k\le118$, (IV) with $k_1=119$.
\item $r=\frac32$, $n=3$: Corollary~\ref{cor:onepolyexact} for $k\le10$,
(III) for $11\le k\le169$, (IV) with $k_1=170$.
\item $r=\frac32$, $n=4$: (I) for $k\le4$, Corollary~\ref{cor:onepolyexact}
for $5\le k\le10$, (II) for $11\le k\le19$, (III) for $20\le k\le231$, (IV)
with $k_1=232$.
\item $r=\frac74$, $n=2$: (I) for $k\le4$, Corollary~\ref{cor:onepolyexact}
for $5\le k\le12$, (II) for $13\le k\le47$, (III) for $48\le k\le820$, (IV)
with $k_1=821$.
\end{itemize}
\endgroup
The sets $Y$ of source (II), with $Z^*=[1,k-1]\cup Y$ certified by
Proposition~\ref{prop:toprowcertS} and $\Phi_r([1,i-1]\cup Y)\le\Phi_r(Z^*)$
for $i<k$, are, for $k=4,5,\ldots$ in turn:
\begingroup\sloppy
\begin{itemize}
\item $r=\frac54$, $n=7$: $\{11,17,25,35,48,65\}$, $\{14,21,30,41,54,72\}$;
\item $r=\frac54$, $n=8$: $\{10,16,23,31,42,55,73\}$,
$\{13,20,27,36,48,62,80\}$, $\{16,23,32,41,53,68,87\}$,
$\{19,27,36,46,59,74,94\}$;
\item $r=\frac54$, $n=9$: $\{9,15,21,29,38,49,63,82\}$,
$\{12,18,25,33,43,55,69,89\}$, $\{15,22,29,38,48,61,76,96\}$,
$\{18,25,33,43,54,66,82,103\}$, $\{21,29,38,47,59,72,88,110\}$,
$\{24,33,42,52,64,78,95,116\}$, $\{28,37,46,57,69,84,101,123\}$;
\item $r=\frac32$, $n=4$, from $k=11$: $\{27,36,47\}$, $\{29,39,51\}$,
$\{32,42,54\}$, $\{35,45,57\}$, $\{37,48,61\}$, $\{40,51,64\}$,
$\{43,54,67\}$, $\{45,57,71\}$, $\{48,60,74\}$;
\item $r=\frac74$, $n=2$, from $k=13$: $Y=\{\lfloor7k/3\rfloor+1\}$ for
$13\le k\le47$ (the sets of Corollary~\ref{cor:onepolyexact} for
$5\le k\le12$ have the same form).
\end{itemize}
\endgroup
In every case $i^*=k$, so the rows equal $1/\nu_k(n)$.
\end{proof}

The certified margins in source~(III) grow quickly with $k$: the ratio of the
lower bound for $\tau_k(n)$ to $E_k(r)$ or $N_k(r)$ exceeds $10^6$ at the
end of each range.  Source~(IV) is crude, and so is the lemma behind it: it
holds for all sets $Y\subset[k,\infty)$, not only for the optimal ones.  In
floating point the first $k$ at which $N_k(r)\le\tau_k(n)$ holds is $k=2,4,6,8,11,14$
for $n=5,\ldots,10$ at $r=\frac54$, $k=9$ for $n=3$ and $k=20$ for $n=4$ at
$r=\frac32$, and $k=48$ for $n=2$ at $r=\frac74$.  The one-polynomial
condition $\max_{i<k}A_i(p)\le A_k(p)$ for an optimal $p$, which is what
Theorem~\ref{thm:onepoly} needs, holds much earlier, from about $k=2n$ at
$r=\frac32$.

What remains open for $1<r<2$ are the rows with $\beta_r(n)>1$ and $k$ below
these thresholds.  Many of them are out of reach of every one-polynomial
argument, because $\Lambda_k(n)>\max_{i\le k}\nu_i(n)$ there.  In floating
point this happens already for $n=2$: at $r=\frac{17}{10}$ for $k=2,3$, at
$r=\frac95$ for $2\le k\le8$, and at $r=\frac{19}{10}$ for every $k\le30$
(at $r=\frac{19}{10}$ the maximum of $\nu_i(2)$ over $i\le k$ sits at $i=1$
for $k\le11$).  In the second row with $i^*=1$ the one-point exempted sets
still have room.  At $r=\frac32,\frac74$ and $n=5,6$, where
$\nu_1(n)>\nu_2(n)$, floating point gives
$\mu_r(\{\sigma\},n+1)<0.97\,\nu_1(n)$ for $2\le\sigma\le120$, and the value
tends to $\nu_1(n)$ as $\sigma\to\infty$.  For all large $\sigma$ the
polynomial $r_{Z_0}(s)(1-s/\sigma)$, with $Z_0$ a top-row optimum, has
$\Phi_r<\nu_1(n)$, because $\Phi_r(r_{Z_0}(s)(1-us))$ is convex in $u$ with
negative right derivative at $u=0$.  For small $\sigma$ the optimal sets are of
a different kind (at $r=\frac74$, $n=5$ and $\sigma=2$ floating point gives
$\{1,2,5,9,14\}$).  So a proof has to switch polynomials between ranges of
$\sigma$, as the trees of Theorem~\ref{thm:rowstree} do for small $n$.

\paragraph{Data availability.}
No datasets were generated or analyzed in this study.

\appendix

\section{Hole products of top-row optima}\label{sec:holeproducts}

Before Theorem~\ref{thm:allrowstwo}, rows equal to the top row were proved
through Theorem~\ref{thm:rowspi}: it suffices that some top-row optimum $Z_0$
has hole product $\Pi(Z_0)\le r-1$.  This appendix proves that bound for every
$n$ at $r\ge\frac92$, $r\ge4$ and $r\ge\frac{25}7$, and the second-row
criterion of Corollary~\ref{cor:secondtail} below $4$.  Theorem~\ref{thm:allrowstwo}
supersedes all of these statements about rows.  We keep the appendix for its
information about top-row optima: the balance of the largest zero
(Lemma~\ref{lem:topzero}), the first hole (Lemma~\ref{lem:firsthole}), the
Jensen bound $\beta_r(n)\ge r^n/(8n)$ (Lemma~\ref{lem:jensentop}), block and
tail bounds on the hole product (Lemmas~\ref{lem:holeblock}
and~\ref{lem:tailmean}), and the explicit bounds on $\Pi(Z_0)$, such as
$\Pi(Z_0)\le\frac{18}7$ at $r\ge\frac{25}7$ (Theorem~\ref{thm:rootbelowfour}).

\subsection{The hole product at roots at least \texorpdfstring{$\frac92$}{9/2}}\label{sec:bigroot}

Theorem~\ref{thm:rowspi} shows that rows equal the top row as soon as some
top-row optimum has hole product at most $r-1$.  So far this was checked one
$n$ at a time (Corollary~\ref{cor:rowspiranges}).  This subsection proves it
for every $n$ when $r\ge\frac92$ (Theorem~\ref{thm:bigroot}).  So the
conjecture stated after Corollary~\ref{cor:belowtwo} holds at every root
$r\ge\frac92$.

Three crude facts about a top-row optimum suffice.
\begin{itemize}
\item Its largest zero balances the weight below it against the weight above
it (Lemma~\ref{lem:topzero}).  This bounds the largest zero.
\item Its first hole carries weight at most $1/\beta_r(n)$
(Lemma~\ref{lem:firsthole}).
\item Jensen's formula gives $\beta_r(n)\ge r^n/(8n)$
(Lemma~\ref{lem:jensentop}).  With the previous item this bounds the initial
run from below.
\end{itemize}
For large $n$ they give an initial run longer than $0.52n$ and a largest zero
below $2.2n$, so the hole product is below $3.32$.  For $3\le r<\frac92$ the
same facts give finite ranges (Corollary~\ref{cor:bigrootfinite}).

Throughout, a \emph{top-row optimum} for $(r,n)$ is a set $Z_0$ of $n-1$
positive integers with $\Phi_r(Z_0)=1/\beta_r(n)$.  For $n\ge2$ and $x\ge1$
put
\[
  T_r(n,x)=\sum_{s\ge x}s\binom{s-1}{n-2}r^{-s} .
\]

\begin{lemma}[The largest zero]\label{lem:topzero}
Let $r>1$ and $n\ge2$.
\begin{enumerate}
\item[(a)] A top-row optimum for $(r,n)$ exists.
\item[(b)] Let $Z_0$ be a top-row optimum with a hole, $x=\max Z_0$,
$W=Z_0\setminus\{x\}$, and let $\ell$ be the largest integer with
$[1,\ell]\subseteq Z_0$.  Then
\[
  \sum_{s<x}s\,|r_W(s)|\,r^{-s}\ \le\ \sum_{s\ge x}s\,|r_W(s)|\,r^{-s}
  \qquad\text{and}\qquad
  \frac{\ell+1}{\beta_r(n-1)}\ \le\ 2\,T_r(n,x) .
\]
\end{enumerate}
\end{lemma}

\begin{proof}
Let $\mathcal A$ be the set of real $p$ with $\deg p\le n-1$ and $p(0)=1$.
For $D\in\{n+1,n+2,\ldots\}\cup\{\infty\}$ put
$f_D(p)=\sum_{s=1}^D|p(s)|r^{-s}$, so $f_\infty=\Phi_r$.

\emph{Step 1.}  Let $Z$ be a set of $n-1$ positive integers with
$x=\max Z\le D$ such that $r_Z$ minimizes $f_D$ on $\mathcal A$.  Put
$W=Z\setminus\{x\}$ and $w(s)=|r_W(s)|r^{-s}$.  For $u\ge0$ the polynomial
$r_W(s)(1-us)$ lies in $\mathcal A$, so
$F(u)=\sum_{s\le D}w(s)|1-us|$ is at least $F(1/x)=f_D(r_Z)$.  The function
$F$ is convex.  Its right derivative at $1/x$ is
$\sum_{x\le s\le D}sw(s)-\sum_{s<x}sw(s)$.  For $D=\infty$ the difference
quotients are dominated by $sw(s)$, and $\sum_ssw(s)<\infty$.  So
\[
  \sum_{s<x}sw(s)\le\sum_{x\le s\le D}sw(s).
\]
For $s\ge x$ every factor $1-s/z$ of $r_W(s)$ is negative.  If
$z_1<\cdots<z_{n-2}$ are the elements of $W$, then $z_i\ge i$, so
$|r_W(s)|=\prod_i(s-z_i)/z_i\le\prod_i(s-i)/i=\binom{s-1}{n-2}$.  Let $m$ be
the least positive integer outside $W$.  Then $w$ vanishes on $[1,m-1]$, and
\[
  f_D(r_W)=\sum_{s\le D}w(s)\le\frac1m\sum_{s\le D}sw(s)
  \le\frac2m\sum_{x\le s\le D}sw(s)\le\frac2m\,T_r(n,x).
\]

\emph{Step 2: (a).}  By the proof of Lemma~\ref{lem:intzeros} (with
$\omega(s)=r^{-s}$ and $i=1$), for each finite $D\ge n+1$ some set $Z_D$ of
$n-1$ integers in $[1,D]$ has $r_{Z_D}$ minimizing $f_D$ on $\mathcal A$,
and $\min_{\mathcal A}f_D$ increases to $\mu_r(\varnothing,n)$.  Let $c$ be the
minimum of $f_{n+1}$ over the real $q$ with $\deg q\le n-2$ and $q(0)=1$.
It is attained and positive, because such a $q$ cannot vanish at the $n+1$
points of $[1,n+1]$.  Step 1 with $m\ge1$ gives
$c\le f_{n+1}(r_{W_D})\le f_D(r_{W_D})\le2T_r(n,\max Z_D)$.  Since
$T_r(n,x)\to0$ as $x\to\infty$, all the sets $Z_D$ lie in one interval
$[1,X]$.  So one set $Z_0$ equals $Z_D$ for infinitely many $D$.  For these
$D$, $f_D(r_{Z_0})=\min_{\mathcal A}f_D$.  Letting $D\to\infty$ gives
$\Phi_r(Z_0)=\mu_r(\varnothing,n)=1/\beta_r(n)$.

\emph{Step 3: (b).}  Apply Step 1 with $D=\infty$ and $Z=Z_0$.  Since $Z_0$
has a hole, $x\ge\ell+2$, so $[1,\ell]\subseteq W$ and $m=\ell+1$.  This gives
the first inequality and $\Phi_r(W)\le2T_r(n,x)/(\ell+1)$.  Finally $r_W$ has
degree $n-2$ and value $1$ at $0$, so
$\Phi_r(W)\ge\mu_r(\varnothing,n-1)=1/\beta_r(n-1)$.
\end{proof}

So the largest zero of a top-row optimum sits at a weighted median: at least
half of the weight $s|r_W(s)|r^{-s}$ lies at or above it.  Above the largest
zero this weight is at most $s\binom{s-1}{n-2}r^{-s}$, which is tiny far out.
So the largest zero cannot be far out.

\begin{lemma}[The first hole]\label{lem:firsthole}
Let $r>1$, and let $Z$ be a set of $n-1$ positive integers with a hole.  Put
$x=\max Z$, let $C=[1,x]\setminus Z$ be its set of holes, and let
$\ell+1=\min C$.  Then
\[
  r^{\ell+1}\,\Phi_r(Z)\ \ge\ |r_Z(\ell+1)|
  =\binom x{\ell+1}^{-1}\prod_{c\in C,\ c>\ell+1}\frac c{c-\ell-1}
  \ \ge\ \binom n{\ell+1}^{-1},
\]
and $\Pi(Z)\le(x-n+\ell+2)/(\ell+1)$.
\end{lemma}

\begin{proof}
The first inequality is the term $s=\ell+1$ of $\Phi_r(Z)$.  Write
$Z=[1,\ell]\cup Y$ with $Y=[\ell+2,x]\setminus C$.  Then
$|r_{[1,\ell]}(\ell+1)|=1$ and $|r_Y(\ell+1)|=\prod_{y\in Y}(y-\ell-1)/y$.  The
product of $(y-\ell-1)/y$ over all $y\in[\ell+2,x]$ is
$(x-\ell-1)!\,(\ell+1)!/x!=\binom x{\ell+1}^{-1}$.  Dividing by the factors at
the holes $c>\ell+1$ gives the identity.  For the lower bound, list $Y$ as
$y_1<\cdots<y_{n-1-\ell}$.  Then $y_i\ge\ell+1+i$, and $(y-\ell-1)/y$
increases with $y$.  So the product is at least
$\prod_ii/(\ell+1+i)=\binom n{\ell+1}^{-1}$.

There are $|C|=x-n+1$ holes, all at least $\ell+1$.  The product of
$1+1/c$ over them is largest when they are $\ell+1,\ldots,\ell+|C|$, where it
equals $(\ell+|C|+1)/(\ell+1)$.
\end{proof}

\begin{lemma}[Jensen's bound for the top row]\label{lem:jensentop}
Let $r>1$, $n\ge1$ and $1/r<\rho<1$.  Then
$\beta_r(n)\ge((\rho r)^n-1)(1-\rho)/\rho$.  If $r\ge3$ and $n\ge2$, then
$\beta_r(n)\ge r^n/(8n)$.
\end{lemma}

\begin{proof}
Let $F$ be a monic multiple of $(x-r)^n$ of degree $D$, and let
$F^*(t)=t^DF(1/t)=1+\sum_{j\ge1}f_jt^j$, where $f_j$ is the coefficient at
backward distance $j$.  As in the proof of \cite[Proposition~4.3]{coefficient-mass},
$F^*$ has a zero of order at least $n$ at $1/r$, and Jensen's formula on
$|t|\le\rho$ gives
\[
  (\rho r)^n\le\max_{|t|=\rho}|F^*(t)|\le1+b_1(F)\frac\rho{1-\rho}.
\]
Take the infimum over $F$; by Proposition~\ref{prop:rowvalue}(c) it is
$\beta_r(n)$.  For the second claim take $\rho=1-1/n$, which exceeds $1/r$.
Since $(1-1/n)^n\ge\frac14$,
$\beta_r(n)\ge(r^n/4-1)/(n-1)\ge(r^n/4-1)/n\ge r^n/(8n)$, the last step because
$r^n\ge8$.
\end{proof}

The polynomial $(x-r)^n$ itself gives the trivial upper bound
$\beta_r(n)\le\max_j\binom nj r^j\le(r+1)^n$.

\begin{theorem}[All rows at large roots]\label{thm:bigroot}
Let $r\ge\frac92$ and $n\ge1$.  Every top-row optimum $Z_0$ for $(r,n)$
satisfies $\Pi(Z_0)\le r-1$.  Consequently $V_r(L,L-n+1)=\beta_r(n)$ for
every $L\ge n$.
\end{theorem}

\begin{proof}
For $n=1$, $Z_0=\varnothing$ and $\Pi(Z_0)=1$.  Let $n\ge2$.  If $Z_0$ has no
hole, $\Pi(Z_0)=1$.  Otherwise let $x=\max Z_0$ and let $\ell$ be as in
Lemma~\ref{lem:firsthole}, so $0\le\ell\le n-2$ and $x\ge n$.  Put
$r_0=\frac92$.  We collect three inequalities.
\begin{itemize}
\item[(i)] $8n\binom n{\ell+1}\ge r_0^{\,n-\ell-1}$.  Indeed, by
Lemmas~\ref{lem:jensentop} and~\ref{lem:firsthole},
$r^n/(8n)\le\beta_r(n)=1/\Phi_r(Z_0)\le r^{\ell+1}\binom n{\ell+1}$, and
$r\ge r_0$.
\item[(ii)] $1\le2(r_0+1)^{n-1}T_{r_0}(n,x)$.  Indeed,
Lemma~\ref{lem:topzero}(b) and $\beta_r(n-1)\le(r+1)^{n-1}$ give
$1\le2(r+1)^{n-1}T_r(n,x)$.  Every term of $T_r(n,x)$ has $s\ge x>n-1$, and
$r\mapsto(r+1)^{n-1}r^{-s}$ decreases, since its logarithmic derivative is
$(n-1)/(r+1)-s/r<0$.
\item[(iii)] $(\ell+1)(r_0-1)\le2r_0^{\ell+1}\binom n{\ell+1}T_{r_0}(n,x)$.
Indeed, Corollary~\ref{cor:prefixmono}(b) and Lemma~\ref{lem:firsthole} give
$\beta_r(n-1)\le\beta_r(n)/(r-1)\le r^{\ell+1}\binom n{\ell+1}/(r-1)$.  With
Lemma~\ref{lem:topzero}(b) this is (iii) at $r$.  Divided by $r-1$, its right
side is a sum of terms $s\binom{s-1}{n-2}\binom n{\ell+1}r^{\ell+1-s}/(r-1)$
with $s>\ell+1$, and each decreases in $r$.
\end{itemize}
By Lemma~\ref{lem:firsthole},
$\Pi(Z_0)\le(x-n+\ell+2)/(\ell+1)=1+(x-n+1)/(\ell+1)$.  We show that this is
at most $r_0-1=\frac72$.

\emph{Large $n$.}  Let $n\ge301$, $m=n-\ell-1$ and $\mu=m/n$.  With
$h(t)=-t\log t-(1-t)\log(1-t)$ we have $\binom nm\le e^{nh(\mu)}$, so (i) gives
$n\psi(\mu)\le\log(8n)$ for $\psi(\mu)=\mu\log r_0-h(\mu)$.  The function
$\psi$ is convex with $\psi'(\mu)=\log(r_0\mu/(1-\mu))$, which is positive for
$\mu>\frac2{11}$.  If $\mu\ge0.48$, then $\psi(\mu)\ge\psi(0.48)>0.0296$, and
$0.0296\,n>\log(8n)$ for $n\ge301$.  So $\mu<0.48$, that is, $\ell+1>0.52n$.

Next suppose $x\ge2.2n$, and put $t_s=s\binom{s-1}{n-2}r_0^{-s}$.  For
$s\ge x$,
\[
  \frac{t_{s+1}}{t_s}=\frac{s+1}{r_0(s-n+2)}\le\frac{2.2n+1}{r_0(1.2n+2)}
  <\frac{11}{6r_0}<0.41 .
\]
So $T_{r_0}(n,x)\le t_{x_0}/0.59$ with $x_0=\lceil2.2n\rceil$.  Here
$t_{x_0}=(n-1)\binom{x_0}{n-1}r_0^{-x_0}$, and
$\binom{x_0}{n-1}\le e^{x_0h((n-1)/x_0)}$.  Put $K=n-1$.  The function
$N\mapsto Nh(K/N)-N\log r_0$ has derivative $\log(N/(N-K))-\log r_0$, which
is negative for $N\ge2.2n$.  Also $h$ increases on $[0,\frac12]$, and
$(n-1)/(2.2n)<1/2.2$.  So
\[
  2(r_0+1)^{n-1}T_{r_0}(n,x)\le\frac{2(n-1)}{0.59(r_0+1)}\,
  e^{n(2.2h(1/2.2)+\log(r_0+1)-2.2\log r_0)}<0.62\,(n-1)e^{-0.088n}<1,
\]
which contradicts (ii).  So $x<2.2n$, and
$\Pi(Z_0)<1+(1.2n+1)/(0.52n)<3.32$.

\emph{Small $n$.}  For $2\le n\le300$ we checked, in exact integer
arithmetic, every $\ell\in[0,n-2]$ satisfying (i) together with the largest
$x\ge n$ satisfying (iii).  In every case $(x-n+\ell+2)/(\ell+1)\le\frac72$.
The bound increases with $x$, so all smaller $x$ are covered too.  The sums
$T_{r_0}(n,x)$ were computed exactly for $s\le8n+40$, and the rest was bounded
by a geometric series; this can only enlarge the set of admissible $x$.  The
largest value found is below $2.54$.

So $\Pi(Z_0)\le\frac72\le r-1$ in all cases.  Lemma~\ref{lem:topzero}(a)
provides a top-row optimum.  Theorem~\ref{thm:rowspi} with this $Z_0$ gives
$V_r(L,L-n+1)\ge1/\Phi_r(Z_0)=\beta_r(n)$, and
Proposition~\ref{prop:rowvalue}(c) gives the reverse inequality.
\end{proof}

By Corollary~\ref{cor:prefixmono}(a), the theorem is the conjecture stated
after Corollary~\ref{cor:belowtwo}, for every $r\ge\frac92$.  It contains the
cases $r\ge\frac92$ of Corollaries~\ref{cor:rowstopranges}
and~\ref{cor:rowspiranges}.  At these roots the conclusions of
Corollaries~\ref{cor:secondrow} and~\ref{cor:crossreduce} hold for every $n$,
without their hypotheses.  The proof never looks at the exempted positions.
It only shows that at these roots every top-row optimum has a long initial
run and few holes.

The same argument works at $r=4$ for bounded $n$.  In the proof, the root
enters only through (i), (ii), (iii) and the target $r-1$.  All three
inequalities become weaker as the root decreases.  So a check at one root
$r_0\ge3$ covers every $r\ge r_0$.  (Lemma~\ref{lem:jensentop} needs $r\ge3$.)

\begin{corollary}[Finite ranges below $\frac92$]\label{cor:bigrootfinite}
$V_r(L,L-n+1)=\beta_r(n)$ for every $L\ge n$ in the following cases:
$r\ge4$ and $n\le1000$; $r\in\{\frac{13}4,\frac72\}$ and $n\le50$; $r=3$
and $n\in\{28,29,30\}$.
\end{corollary}

\begin{proof}
For $r\ge4$, run the small-$n$ part of the proof of
Theorem~\ref{thm:bigroot} with $r_0=4$ and the target $3$.  We checked in
exact integer arithmetic that, for $2\le n\le1000$, every pair $(\ell,x)$
allowed by (i) and (iii) at $r_0=4$ has $(x-n+\ell+2)/(\ell+1)\le3$.  For
$r\in\{\frac{13}4,\frac72\}$ and $r=3$ we found top-row optima $Z_0$ by the
continuation search of
Corollary~\ref{cor:secondrow}, certified them by
Proposition~\ref{prop:toprowcert}, and computed $\Pi(Z_0)$ exactly.  In all
these cases $\Pi(Z_0)\le r-1$; the maxima are $2$ at $r=\frac{13}4$ and $1.82$
at $r=\frac72$.  Theorem~\ref{thm:rowspi} concludes.
\end{proof}

The method has a limit near $r=4$.  In floating point, the largest value of
$1+(x-n+1)/(\ell+1)$ allowed by (i) and (iii) at the level of exponential
rates tends, as $n\to\infty$, to $2.57$ at $r=\frac92$, to exactly $3=r-1$ at
$r=4$ (attained at $\ell\approx n/2$, $x\approx2n$), and to $3.71$ at
$r=\frac72$.  So a proof for all $n$ at $r=4$ needs a sharper bound than the count;
Section~\ref{sec:rootfour} gives one.  At $r_0\in\{3,\frac{13}4,\frac72\}$ the finite check already
fails for every $n$ with $3\le n\le200$.  Near $r=3$ Theorem~\ref{thm:rowspi}
itself stops working.
At $r=3$ the certified optima have $\Pi(Z_0)>2$ for
$n\in[25,27]\cup[31,42]$, and at $r=\frac52$ they have $\Pi(Z_0)>\frac32$ for
every $n\le36$.  The true values are far from the bounds: at $r=4$ and $n=50$
the certified optimum has $\ell=30$, $\max Z_0=76$ and $\Pi(Z_0)=1.65$.

For $2\le r<\frac92$ the second row therefore remains open for large $n$
(Sections~\ref{sec:rootfour} and~\ref{sec:rootbelowfour} lower $\frac92$ to $4$
and then to $\frac{10}3$).
The criterion of Corollary~\ref{cor:secondrow} has a simple form in the
root.  For $W=[1,\ell]\cup Y$ as there, termwise differentiation gives
$\partial_r\Phi_r(W)=-r^{-1}\sum_ss|r_W(s)|r^{-s}$, so
\[
  \frac{\Phi_r(\mathrm{ins}(W,\ell+1))}{\Phi_r(W)}
  =-\frac{\partial_r\log\Phi_r(W)}{(\ell+1)\Pi_Y},\qquad
  \Pi_Y=\prod_{y\in Y}\Bigl(1+\frac1y\Bigr).
\]
The sum $\Phi_r(W)=\sum_s|r_W(s)|r^{-s}$ is log-convex in $\log r$, so for a
fixed set $W$ the ratio, and likewise $\rho$, does not increase with $r$.
At a top-row optimum $Z_0$ for $(r,n)$, $\beta_{r'}(n)\ge1/\Phi_{r'}(Z_0)$ at
every root $r'$, with equality at $r'=r$.  Hence the ratio is at most $\bigl((\ell+1)\Pi_Y\bigr)^{-1}$ times the
lower right Dini derivative of $\log\beta_r(n)$ in $r$.  So the second row
holds at $r$ whenever $\beta_r(n)$ grows slowly enough in the root there.
We could not bound this growth for all $n$.

\subsection{The hole product at roots at least \texorpdfstring{$4$}{4}}\label{sec:rootfour}

At $r=4$ the proof of Theorem~\ref{thm:bigroot} stops because of the count
bound $\Pi(Z_0)\le(x-n+\ell+2)/(\ell+1)$ of Lemma~\ref{lem:firsthole}.  That
bound puts all holes right above the first one.  Such holes make
$|r_{Z_0}(\ell+1)|$ large, and the term $s=\ell+1$ of $\Phi_r(Z_0)$ forbids
this.  This subsection turns the observation into two bounds for the hole
product (Lemma~\ref{lem:holeblock}).  The first one shows that rows equal the
top row for every $n$ at every root $r\ge4$ (Theorem~\ref{thm:rootfour}).  A
bound for the mean of the weight $|r_Z(s)|r^{-s}$ (Lemma~\ref{lem:tailmean})
gives a sufficient condition for the second row that is weaker than
$\Pi(Z_0)\le r-1$ (Corollary~\ref{cor:secondtail}).  With the second bound
this gives finite ranges below $4$ that hold at all roots above a threshold at
once (Corollary~\ref{cor:rootfourfinite}).  Section~\ref{sec:rootbelowfour} extends
them to all $n$ for all rows at $r\ge\frac{25}7$ and for the second row at
$r\ge\frac{10}3$.

Throughout, $Z$ is a set of $n-1$ positive integers with a hole, $x=\max Z$,
$p=\ell+1$ is the first hole of $Z$, and $K=x-n$.  So $Z$ has $K+1$ holes.  By
Lemma~\ref{lem:firsthole},
\[
  Q(Z):=\binom xp\,|r_Z(p)|=\prod_{c\text{ a hole of }Z,\ c>p}\frac c{c-p}\ \ge1 .
\]

\begin{lemma}[Holes above the first hole]\label{lem:holeblock}
Let $Z$, $p$ and $K$ be as above, and put $Q=Q(Z)$.
\begin{enumerate}
\item[(a)] $\Pi(Z)^p\le(1+1/p)^p\,Q$.
\item[(b)] Let $y\ge p+1$ be an integer with $\binom{y+K-1}p\ge Q\binom{y-1}p$.
Then
\[
  \Pi(Z)\le\frac{(p+1)(y+K)}{p\,y} .
\]
\end{enumerate}
\end{lemma}

\begin{proof}
Let $C'$ be the set of holes of $Z$ other than $p$.  It has $K$ elements, all
greater than $p$, and $\Pi(Z)=(1+1/p)\prod_{c\in C'}(1+1/c)$.  For real $c>p$
put $\pi(c)=\log(1+1/c)$ and $\alpha(c)=\log(c/(c-p))$.  Both decrease, and
$\log Q=\sum_{c\in C'}\alpha(c)$.

(a) Since $-\log(1-u)\ge u$ for $0\le u<1$, we have $p\,\pi(c)\le p/c\le\alpha(c)$.
Sum over $C'$.

(b) Let $B=[y,y+K-1]$.  The products telescope:
$\sum_{c\in B}\pi(c)=\log((y+K)/y)$ and
$\sum_{c\in B}\alpha(c)=\log\bigl(\binom{y+K-1}p/\binom{y-1}p\bigr)$.  So the
hypothesis says $\sum_{C'}\alpha\le\sum_B\alpha$, and we must show
$\sum_{C'}\pi\le\sum_B\pi$.  If $K=0$ there is nothing to show.  If $y=p+1$,
list $C'$ as $c_1<\cdots<c_K$.  Then $c_i\ge p+i$, so
$\sum_{C'}\pi\le\sum_i\pi(p+i)=\sum_B\pi$.

Now let $y\ge p+2$ and $K\ge1$.  Put
\[
  \lambda=\frac{\pi(y-1)-\pi(y+K-1)}{\alpha(y-1)-\alpha(y+K-1)}>0,\qquad
  g=\pi-\lambda\alpha\quad\text{on }(p,\infty).
\]
Then $g(y-1)=g(y+K-1)$, and
\[
  g'(c)=\frac{(\lambda p-1)c+p(1+\lambda)}{c(c+1)(c-p)} .
\]
The numerator is linear in $c$ and positive at $c=p$.  It cannot stay
positive, since $g$ takes the same value twice.  So $g$ strictly increases and
then strictly decreases.  With $v=g(y-1)$, this gives $g>v$ on $(y-1,y+K-1)$
and $g<v$ outside $[y-1,y+K-1]$.  The elements of $C'\setminus B$ lie outside
$[y,y+K-1]$, so $g\le v$ there.  The elements of $B\setminus C'$ lie in
$[y,y+K-1]$, so $g\ge v$ there.  The two sets have the same size.  Hence
$\sum_{C'}g\le\sum_Bg$, and
\[
  \sum_{C'}\pi=\sum_{C'}g+\lambda\sum_{C'}\alpha\le\sum_Bg+\lambda\sum_B\alpha
  =\sum_B\pi .\qedhere
\]
\end{proof}

For $y=p+1$, part (b) is the count bound of Lemma~\ref{lem:firsthole}.  The
best $y$ is the largest one allowed.  Part (b) says that, among the hole sets
that cost at most as much in $Q$, a block of consecutive holes has the
largest hole product.

\begin{theorem}[All rows at roots at least $4$]\label{thm:rootfour}
Let $r\ge4$ and $n\ge1$.  Every top-row optimum $Z_0$ for $(r,n)$ satisfies
$\Pi(Z_0)\le r-1$.  Consequently $V_r(L,L-n+1)=\beta_r(n)$ for every $L\ge n$.
\end{theorem}

\begin{proof}
If $Z_0$ has no hole, $\Pi(Z_0)=1$.  Every top-row optimum with a hole
satisfies (i) and (iii) of the proof of Theorem~\ref{thm:bigroot} at $r_0=4$.
So for $n\le1000$ the claim is the check in the proof of
Corollary~\ref{cor:bigrootfinite}, which gives $\Pi(Z_0)\le3$.  Let
$n\ge301$, let $Z_0$ have a hole, and let $x$, $p=\ell+1$ and $K$ be as above.
Put $\theta=p/n$, and let $h$ be as in the proof of Theorem~\ref{thm:bigroot}.

\emph{Step 1: $\theta>0.48$.}  By (i) and $\binom np\le e^{nh(\theta)}$,
$n\bigl((1-\theta)\log4-h(\theta)\bigr)\le\log(8n)$.  The left side is convex
in $\theta$, vanishes at $\theta=\frac12$, and has derivative $-n\log4$ there.
So $n(\frac12-\theta)\log4\le\log(8n)$, and
$\theta\ge\frac12-\log(8n)/(n\log4)>0.48$.

\emph{Step 2.}  If $x\le n+2p-1$, Lemma~\ref{lem:firsthole} gives
$\Pi(Z_0)\le(x-n+1+p)/p\le3$.  So let $x\ge x_0:=n+2p$.

\emph{Step 3.}  By Lemma~\ref{lem:firsthole},
$\beta_r(n)=1/\Phi_r(Z_0)\le r^p/|r_{Z_0}(p)|$.  With
Corollary~\ref{cor:prefixmono}(b) this gives
$\beta_r(n-1)\le r^p/\bigl((r-1)|r_{Z_0}(p)|\bigr)$, and
Lemma~\ref{lem:topzero}(b) gives
\[
  p\,(r-1)\,|r_{Z_0}(p)|\le2r^pT_r(n,x).
\]
This is (iii) before $|r_{Z_0}(p)|$ is bounded below.  As in (iii),
$r^pT_r(n,x)/(r-1)$ decreases in $r$.  With Lemma~\ref{lem:holeblock}(a) and
$(1+1/p)^p<e$,
\[
  \Pi(Z_0)^p<\frac{2e}{3p}\,4^p\,H(x),\qquad H(x)=\binom xp\,T_4(n,x).
\]

\emph{Step 4: $H(x)\le H(x_0)$.}  Put $t_s=s\binom{s-1}{n-2}4^{-s}$ and
$q_s=t_{s+1}/t_s=(s+1)/(4(s-n+2))$.  Since $q_s$ decreases in $s$,
$T_4(n,x+1)=\sum_{s\ge x}q_st_s\le q_xT_4(n,x)$, and
$T_4(n,x)\le t_x/(1-q_x)$ when $q_x<1$.  Hence
\[
  \frac{H(x+1)}{H(x)}\le\frac{(x+1)^2}{4(x+1-p)(x-n+2)} .
\]
The right side decreases in $x$.  At $x_0$ it is at most $1$, because
$4(n+p+1)(2p+2)-(n+2p+1)^2=4p^2+4np+12p+6n+7-n^2>0$ for $p\ge n/4$.

\emph{Step 5: the value at $x_0$.}  Here
$q:=q_{x_0}=(n+2p+1)/(8(p+1))<(1/0.48+2)/8<0.52$ and
$t_{x_0}=(n-1)\binom{x_0}{n-1}4^{-x_0}$.  Also
$\binom{x_0}{n-1}=\binom{x_0}n\,n/(2p+1)<\binom{x_0}n/0.96$ and
$(n-1)/p<1/0.48$.  Use $\binom Nk\le e^{Nh(k/N)}$ and $x_0=(1+2\theta)n$.
Then
\[
  \log\frac{\Pi(Z_0)^p}{3^p}<2.12+n\varphi(\theta),\qquad
  \varphi(\theta)=(1+2\theta)\Bigl(h\Bigl(\frac\theta{1+2\theta}\Bigr)
  +h\Bigl(\frac1{1+2\theta}\Bigr)\Bigr)-(1+\theta)\log4-\theta\log3,
\]
where $2.12>1+\log\frac23+2\log\frac1{0.48}+\log\frac1{0.96}$.  The function
$\varphi$ is concave on $(0,1)$, because $(u,v)\mapsto vh(u/v)$ is concave on
$0<u<v$.  Its derivative is
$\varphi'(\theta)=\log\bigl((1+2\theta)^4/(48\theta^3(1+\theta))\bigr)$.  At
$\theta=\frac56$ we have $\varphi=-0.03665\ldots$ and
$\varphi'=-0.00705\ldots$, so the tangent there gives $\varphi<-0.0307$ on
$(0,1)$.  Hence $\log(\Pi(Z_0)^p/3^p)<2.12-0.03n<0$.

So $\Pi(Z_0)\le3\le r-1$ in all cases.  Lemma~\ref{lem:topzero}(a) provides a
top-row optimum, Theorem~\ref{thm:rowspi} gives
$V_r(L,L-n+1)\ge\beta_r(n)$, and Proposition~\ref{prop:rowvalue}(c) gives the
reverse inequality.
\end{proof}

The theorem contains Theorem~\ref{thm:bigroot} and the case $r\ge4$ of
Corollary~\ref{cor:bigrootfinite}.  By Corollary~\ref{cor:prefixmono}(a) it
proves the conjecture stated after Corollary~\ref{cor:belowtwo} for every
$r\ge4$.  Only Step~2 uses the count bound, and there it is not the limit.

\paragraph{The second row below $4$.}
Corollary~\ref{cor:secondrow} needs $1+\rho\le r\Pi_Y$.  Both sides can be
written with $\Pi(Z_0)$ and the mean of the weight $|r_Z(s)|r^{-s}$.  That mean
is controlled by the zeros near $x$.

\begin{lemma}[The mean of the weight]\label{lem:tailmean}
Let $r>1$, let $Z$ be a nonempty set of $n-1$ positive integers, $x=\max Z$ and
$w(s)=|r_Z(s)|r^{-s}$.  Then
\[
  \sum_{s\ge1}s\,w(s)\le(x+1+M)\,\Phi_r(Z),\qquad M=\frac n{r-1} .
\]
If $Z$ has a hole, $p$ is its first hole and $q=(x+1)/(r(x-p+2))<1$, one may
take $M=\min\{n/(r-1),\ (n-p+1)q/(1-q)\}$.
\end{lemma}

\begin{proof}
Every $s\le x$ is less than $x+1$.  So the claim follows once the weight
$u(d)=w(x+d)$, $d\ge1$, has mean $\sum_ddu(d)/\sum_du(d)$ at most $1+M$.  For
$s>x$,
\[
  \frac{w(s+1)}{w(s)}=\frac1r\prod_{z\in Z}\Bigl(1+\frac1{s-z}\Bigr).
\]
The numbers $s-z$ are $n-1$ distinct integers at least $d=s-x$.  So the product
is at most $\prod_{j=d}^{d+n-2}(1+1/j)=(d+n-1)/d$.  Let
$v(d)=\binom{d+n-2}{n-1}r^{1-d}$.  Then $v(d+1)/v(d)=(d+n-1)/(rd)$, which is at
least $u(d+1)/u(d)$.  So $u/v$ does not increase in $d$.  Under $v$ the
covariance of $d$ and $u/v$ is therefore not positive, and the mean of $u$ is
at most the mean of $v$.  With $j=d-1$, the weight $v$ is proportional to
$\binom{j+n-1}{n-1}r^{-j}$, a negative binomial law with mean $n/(r-1)$.  So
the mean of $v$ is $1+n/(r-1)$.

For the second value of $M$, divide the product over $[1,x]$, which is
$s/(s-x)$, by the product over the set $C$ of holes.  The $x-n+1$ holes are
distinct integers at least $p$, and $1+1/(s-c)$ increases with $c$.  So
$\prod_{c\in C}(1+1/(s-c))\ge\prod_{j=0}^{x-n}(1+1/(s-p-j))=(s-p+1)/(s-p-x+n)$.
With $d=s-x$ this gives
\[
  \frac{u(d+1)}{u(d)}\le\frac{(x+d)(d+n-p)}{r\,d\,(x+d-p+1)}\le q\,\frac{d+n-p}d,
\]
because $(x+d)/(x+d-p+1)$ does not increase in $d$.  Compare with
$v(d)=\binom{d+n-p-1}{n-p}q^{d-1}$.  Its ratio is $q(d+n-p)/d$, and its mean is
$1+(n-p+1)q/(1-q)$.
\end{proof}

In the notation of the remark at the end of Section~\ref{sec:bigroot}, the
lemma says $-\partial_r\log\Phi_r(Z)\le(x+1+M)/r$.

\begin{corollary}[The second row through the tail]\label{cor:secondtail}
Let $r>1$ and $n\ge2$, and let $Z_0$ be a top-row optimum for $(r,n)$ with a
hole.  Let $x=\max Z_0$ and let $M$ be as in Lemma~\ref{lem:tailmean}.  If
\[
  \Pi(Z_0)\,(x+1+M)\le r\,(x+1),
\]
then $V_r(n+1,2)=\beta_r(n)$.  This holds in particular when
$\Pi(Z_0)\le r-1$.  If $r\ge2$ and the top-row optimum is $[1,n-1]$, again
$V_r(n+1,2)=\beta_r(n)$.
\end{corollary}

\begin{proof}
Write $Z_0=[1,\ell]\cup Y$ as in Corollary~\ref{cor:secondrow}, and $p=\ell+1$.
The integers in $[p+1,x]$ are the elements of $Y$ and the holes other than
$p$.  So $\Pi_Y\,\Pi(Z_0)/(1+1/p)=\prod_{z=p+1}^x(1+1/z)=(x+1)/(p+1)$, that is,
$r\Pi_Y=r(x+1)/(p\,\Pi(Z_0))$.  By the remark after
Corollary~\ref{cor:secondrow}, $1+\rho=\sum_ss\,w(s)/(p\,\Phi_r(Z_0))$ with $w$
as in Lemma~\ref{lem:tailmean}.  So $1+\rho\le r\Pi_Y$ is the inequality
$\Pi(Z_0)\sum_ssw(s)\le r(x+1)\Phi_r(Z_0)$, and the lemma gives it.
Corollary~\ref{cor:secondrow} concludes.  If $\Pi(Z_0)\le r-1$, then
$(r-1)(x+1+M)\le r(x+1)$, since $(r-1)M\le n\le x+1$.  For $Z_0=[1,n-1]$ we
have $Y=\varnothing$, $\Pi_Y=1$ and $p=n$, and the lemma gives
$1+\rho\le(n+n/(r-1))/n=r/(r-1)\le r$.
\end{proof}

At the certified optima of Corollary~\ref{cor:secondrow} the condition holds
for every $n\le42$ at $r=3$, where $\Pi(Z_0)\le r-1$ fails for
$n\in[25,27]\cup[31,42]$ (the remark after Corollary~\ref{cor:bigrootfinite}).
It fails at every certified optimum at $r=2$, $\frac94$ and $\frac52$.

\begin{corollary}[Ranges for all roots above a threshold]\label{cor:rootfourfinite}
\begin{enumerate}
\item[(a)] $V_r(L,L-n+1)=\beta_r(n)$ for every $L\ge n$ if $r\ge\frac{18}5$
and $n\le300$, or if $r\ge\frac72$ and $n\le27$.
\item[(b)] $V_r(n+1,2)=\beta_r(n)$ if $r\ge\frac{10}3$ and $n\le300$, or if
$r\ge\frac{13}4$, $n\le300$ and $n\notin\{11,12\}$.
\end{enumerate}
\end{corollary}

\begin{proof}
For $n\le2$ this is Corollary~\ref{cor:tworowsr}, and for $n=3$ it is
Corollary~\ref{cor:rowstopranges}.  In (b) with $r\ge\frac{13}4$ and
$4\le n\le10$ it is Corollary~\ref{cor:rowspiranges}.  In the remaining cases
let $r_0\in\{\frac{18}5,\frac72,\frac{10}3,\frac{13}4\}$ be the threshold, and
let $Z_0$ be a top-row optimum with a hole at a root $r\ge r_0$.
Let $p$ be its first hole and $x=\max Z_0$.  Let $J$ be the largest of
$r_0^n/(8n)$ and the values $((\rho r_0)^n-1)(1-\rho)/\rho$ of
Lemma~\ref{lem:jensentop} for $\rho\in\{1-k/400\}$ with $\rho r_0>1$.  For a
fixed $\rho$, $r^p/(((\rho r)^n-1)(1-\rho)/\rho)$ decreases in $r$.  So
Lemma~\ref{lem:firsthole} and Step~3 of the proof of
Theorem~\ref{thm:rootfour} give
\[
  \binom np^{-1}\le|r_{Z_0}(p)|\le\min\Bigl\{\frac{r_0^p}J,\
  \frac{2r_0^pT_{r_0}(n,x)}{p(r_0-1)}\Bigr\}.
\]
Put $Q$ equal to $\binom xp$ times this minimum, so $Q\ge Q(Z_0)$.  For every
pair $(p,x)$ allowed by these inequalities we checked, in exact rational
arithmetic, that the count bound, Lemma~\ref{lem:holeblock}(a) or
Lemma~\ref{lem:holeblock}(b) bounds $\Pi(Z_0)$ by a target.  The target is
$r_0-1$ in (a).  In (b) it is $r_0(x+1)/(x+1+M)$, with $M$ taken at $r_0$ in
Lemma~\ref{lem:tailmean}.  In (b) we used the least $y$ allowed by the target.
The sums $T_{r_0}(n,x)$ were bounded above as in the proof of
Theorem~\ref{thm:bigroot}, and $x$ ran until the right bound fell below
$\binom np^{-1}$; it decreases in $x$.  All bounds only get better as the root
grows, and $M$ decreases in the root.  So the check at $r_0$ covers every
$r\ge r_0$.  Theorem~\ref{thm:rowspi}, resp.\ Corollary~\ref{cor:secondtail},
concludes.
\end{proof}

The check covers between $6.7$ and $8.5$ million pairs at each threshold, and
part (b) of Lemma~\ref{lem:holeblock} is needed at $11$ to $17\%$ of them.  Without it, (a)
fails from $n=15$ on already at $r_0=\frac{15}4$.  At $r_0=\frac72$ the check
for (a) fails for every $n$ with $28\le n\le40$.  At $r_0=\frac{13}4$ the
check for (b) fails at $n=11,12$, and at $r_0=3$ it fails for every $n$ with
$4\le n\le80$.

Theorem~\ref{thm:rootbelowfour} extends the ranges at $\frac{10}3$ and $\frac{13}4$
to all $n$, and gives all rows for every $n$ at $r\ge\frac{25}7$.  In floating
point, at the level of exponential rates as $n\to\infty$, the largest hole
product allowed by (i), (iii) and Lemma~\ref{lem:holeblock}(b) stays below
$r-1$ exactly for $r$ above about $3.55$.  With part (a) instead of (b) the
limit is about $3.94$, so part (a) alone cannot go much below $4$.
The condition of Corollary~\ref{cor:secondtail}, with part (b), holds for large
$n$ exactly for $r$ above about $3.2$.  Theorem~\ref{thm:rootbelowfour}
uses a two-parameter version of Steps~4 and~5.  The worst largest zero is
no longer $x_0$: near $r=3.55$ it is about $2.5n$, with $p$ about $0.75n$.  Two
natural sharpenings do not help.  Counting the holes in $T_r(n,x)$ exactly,
instead of putting them at the top, lowers the limiting hole product at
$r=\frac72$ only from $2.575$ to $2.559$.  In tests at finite $n$, the terms of
$\Phi_r(Z_0)$ at $s>x$ and at the largest hole added nothing either.  The worst
configurations have a long block of holes and a largest zero near $2.5n$, far
from the certified optima (at $r=4$, $n=50$ the optimum has
$\max Z_0=76$).  Excluding them needs the terms of $\Phi_r(Z_0)$ at the inner
holes, and those are not linear in the hole set.

\subsection{The hole product at roots at least \texorpdfstring{$\frac{25}7$}{25/7}}\label{sec:rootbelowfour}

Corollary~\ref{cor:rootfourfinite} gives the rows below $4$ only for
$n\le300$.  This subsection removes the bound on $n$
(Theorem~\ref{thm:rootbelowfour}).  For large $n$ we replace Steps~4 and~5 of
the proof of Theorem~\ref{thm:rootfour} by an inequality between exponential
rates in two variables, $\theta=p/n$ and $\xi=x/n$.  It is checked on a box by
a finite computation in interval arithmetic.  The count bound, part (b) of
Lemma~\ref{lem:holeblock} and the two upper bounds for $|r_{Z_0}(p)|$ from the
proof of Corollary~\ref{cor:rootfourfinite} are all we use.

As before, $h(t)=-t\log t-(1-t)\log(1-t)$ on $[0,1]$.  For a root $r_0>1$
put $L=\log r_0$ and, for $\xi>1$, $u>t>0$ and $\kappa\ge0$,
\[
  \tau(\xi)=\xi h(1/\xi)-\xi L,\qquad
  F(u,t)=u\log u-(u-t)\log(u-t),\qquad
  B(\eta,\kappa,t)=F(\eta+\kappa,t)-F(\eta,t).
\]
The derivatives are
$\tau'(\xi)=\log(\xi/(\xi-1))-L$, $\partial_\kappa B=\log((\eta+\kappa)/(\eta+\kappa-t))$,
$\partial_\eta B=\log((\eta+\kappa)/(\eta+\kappa-t))-\log(\eta/(\eta-t))$ and
$\partial_tB=\log((\eta+\kappa-t)/(\eta-t))$ for $\eta>t$.  So $\tau$
decreases for $\xi>r_0/(r_0-1)$, and $B$ decreases in $\eta$ and increases in
$\kappa$ and in $t$.

\begin{theorem}[All rows at roots at least $\frac{25}7$]\label{thm:rootbelowfour}
\begin{enumerate}
\item[(a)] Let $r\ge\frac{25}7$ and $n\ge1$.  Then $V_r(L,L-n+1)=\beta_r(n)$
for every $L\ge n$.  If $n\ge4$, every top-row optimum $Z_0$ for $(r,n)$
satisfies $\Pi(Z_0)\le\frac{18}7$.
\item[(b)] Let $n\ge1$, and let $r\ge\frac{10}3$, or $r\ge\frac{13}4$ and
$n\notin\{11,12\}$.  Then $V_r(n+1,2)=\beta_r(n)$.
\end{enumerate}
\end{theorem}

\begin{proof}
\emph{Small $n$.}  In (a), for $n\le2$ this is Corollary~\ref{cor:tworowsr},
and for $n=3$ it is Corollary~\ref{cor:rowstopranges}.  For $4\le n\le369$ we
ran the exact check of the proof of Corollary~\ref{cor:rootfourfinite} at
$r_0=\frac{25}7$ with the target $\frac{18}7$.  It passed for all of the
$12.8$ million pairs $(p,x)$.  So every top-row optimum with a hole has
$\Pi(Z_0)\le\frac{18}7\le r-1$, and Theorem~\ref{thm:rowspi} concludes as in
the proof of Theorem~\ref{thm:rootfour}.  In (b), for $n\le300$ this is
Corollary~\ref{cor:rootfourfinite}.

\emph{Large $n$.}  Let $N_0=370$ in (a) and $N_0=301$ in (b), and let
$n\ge N_0$.  Let $r_0$ be $\frac{25}7$ in (a) and $\frac{10}3$ or $\frac{13}4$
in (b), and let $r\ge r_0$.  Let $Z_0$ be a top-row optimum for $(r,n)$.  If
$Z_0$ has no hole, then $\Pi(Z_0)=1$, and Corollary~\ref{cor:secondtail}
gives (b).  So let $Z_0$ have a hole, and let $p$, $x$ and $K=x-n$ be as in
Section~\ref{sec:rootfour}.  Put $\theta=p/n$, $\xi=x/n$, $d=1/n$ and
$e=\log(8n)/n$.  Then $d\le1/N_0$ and $e\le e_1:=\log(8N_0)/N_0$.

\emph{Step 1: the target.}  In (a) put $g=r_0-1$.  In (b) let $M_0$ be the
number $M$ of Lemma~\ref{lem:tailmean} for $Z_0$ at the root $r_0$, and put
$g=r_0(x+1)/(x+1+M_0)$.  We show $\Pi(Z_0)\le g$.  This is enough.  In (a),
$g\le r-1$, and Theorem~\ref{thm:rowspi} concludes.  In (b), the number $M$
at $r$ is at most $M_0$, because $n/(r-1)$ and $q=(x+1)/(r(x-p+2))$ decrease
in $r$, and $q/(1-q)$ increases in $q$.  So $\Pi(Z_0)(x+1+M)\le r(x+1)$, and
Corollary~\ref{cor:secondtail} concludes.  In both cases $g\ge r_0-1$, since
$M_0\le n/(r_0-1)\le(x+1)/(r_0-1)$.

\emph{Step 2: three bounds at $r_0$.}  Put $u=|r_{Z_0}(p)|$.  By
Lemma~\ref{lem:firsthole}, $u\ge\binom np^{-1}\ge e^{-nh(\theta)}$.  As in (i)
of the proof of Theorem~\ref{thm:bigroot}, Lemmas~\ref{lem:firsthole}
and~\ref{lem:jensentop} give $u\le r^p/\beta_r(n)\le8nr^{p-n}\le8nr_0^{p-n}$.
By Step~3 of the proof of Theorem~\ref{thm:rootfour},
$u\le2r_0^pT_{r_0}(n,x)/(p(r_0-1))$.

\emph{Step 3: the count bound.}  If $x-n+1\le(g-1)p$, then
Lemma~\ref{lem:firsthole} gives $\Pi(Z_0)\le1+(x-n+1)/p\le g$.  So let
$x-n+1>(g-1)p\ge(r_0-2)p$.  In the variables, $\xi-1+d>(g-1)\theta$.

\emph{Step 4: a box.}  The first two bounds of Step~2 give
$(1-\theta)L-h(\theta)\le e$.  The left side has derivative
$-L-\log((1-\theta)/\theta)$, which is negative for
$\theta<\theta^*:=r_0/(1+r_0)$.  Let $\theta_0$ be the number in the table
below.  It is less than $\theta^*$, and $(1-\theta_0)L-h(\theta_0)>e_1$.  So
$\theta>\theta_0$.  By Step~3, $\xi>\xi_0:=1+(r_0-2)\theta_0-1/N_0$.

Next we bound $T_{r_0}(n,x)$.  With $t_s=s\binom{s-1}{n-2}r_0^{-s}$ and
$q_s=t_{s+1}/t_s=(s+1)/(r_0(s-n+2))$, as in Step~4 of the proof of
Theorem~\ref{thm:rootfour}, $T_{r_0}(n,x)\le t_x/(1-q_x)$ when $q_x<1$.  Put
$a=x-n+2$.  Then $a>(r_0-2)\theta_0n$, so
\[
  q_x=\frac{a+n-1}{r_0a}<\frac1{r_0}\Bigl(1+\frac1{(r_0-2)\theta_0}\Bigr)
  =:\bar q<1 .
\]
Also $t_x=(n-1)\binom x{n-1}r_0^{-x}$ and
$\binom x{n-1}=\binom xn\,n/(x-n+1)\le e^{xh(n/x)}n/(x-n+1)$, and
$xh(n/x)-xL=n\tau(\xi)$.  Since $p>\theta_0n$ and
$x-n+1>(r_0-2)\theta_0n$,
\[
  \frac{2T_{r_0}(n,x)}{p(r_0-1)}
  \le\frac{2n(n-1)\,e^{n\tau(\xi)}}{p(r_0-1)(x-n+1)(1-q_x)}
  <\frac{2e^{n\tau(\xi)}}{c},\qquad
  c=\theta_0^2(r_0-1)(r_0-2)(1-\bar q).
\]
The table gives $4cN_0\ge1$, so $2/c\le8n$.  So the third bound of Step~2
becomes $u\le e^{n(\theta L+\tau(\xi)+e)}$.  With $u\ge e^{-nh(\theta)}$,
\[
  -h(\theta)-\theta L\le\tau(\xi)+e_1 .
\]
The left side is convex in $\theta$ with least value $-\log(1+r_0)$ at
$\theta^*$.  Let $\xi_1$ be the number in the table.  It exceeds $\xi_0$, and
$\xi_0>r_0/(r_0-1)$, and $\tau(\xi_1)<-\log(1+r_0)-e_1$.  As $\tau$ decreases
for $\xi>r_0/(r_0-1)$, we get $\xi<\xi_1$.  So
$(\theta,\xi)$ lies in the box $[\theta_0,1]\times[\xi_0,\xi_1]$.
\[
\begin{array}{ccccccc}
  r_0&N_0&\theta_0&\xi_0&\xi_1&\bar q&4cN_0\\\hline
  \frac{25}7&370&\frac{11}{25}&\frac{21869}{12950}&\frac{64}{25}&\frac{2072}{3025}&364.7\ldots\\[2pt]
  \frac{10}3&301&\frac{41}{100}&\frac{34841}{22575}&\frac{273}{100}&\frac{174}{205}&95.2\ldots\\[2pt]
  \frac{13}4&301&\frac25&\frac{901}{602}&\frac{14}5&\frac{12}{13}&41.6\ldots
\end{array}
\]

\emph{Step 5: the block bound.}  By Step~2 and $\binom xp\le e^{xh(p/x)}$,
$Q(Z_0)=\binom xpu$ satisfies
\[
  \log Q(Z_0)\le n\bigl(\xi h(\theta/\xi)+\theta L+\min\{-L,\tau(\xi)\}+e\bigr).
\]
Let $g'\le g$ be positive with $p(g'-1)>1$, and put
$y_0=(p+1)K/(p(g'-1)-1)$.  The inequality $(p+1)(y+K)\le g'py$ holds exactly
for $y\ge y_0$.  If $y_0\le p+1$, then
$\Pi(Z_0)\le(p+1+K)/p\le g'\le g$ by the count bound.  Otherwise let
$y=\lceil y_0\rceil\ge p+2$.  If
$\binom{y+K-1}p\ge Q(Z_0)\binom{y-1}p$, then Lemma~\ref{lem:holeblock}(b) gives
$\Pi(Z_0)\le(p+1)(y+K)/(py)\le g$.  Now $\log(c/(c-p))$ decreases in $c>p$, so
\[
  \log\frac{\binom{y+K-1}p}{\binom{y-1}p}=\sum_{c=y}^{y+K-1}\log\frac c{c-p}
  \ge\int_y^{y+K}\log\frac c{c-p}\,dc=F(y+K,p)-F(y,p)
  =nB(\eta,\xi-1,\theta),
\]
with $\eta=y/n$.  The last step uses $F(nu,nt)=nF(u,t)+nt\log n$.  Also
$\eta<y_0/n+d=(\theta+d)(\xi-1)/((g'-1)\theta-d)+d$.  Hence $\Pi(Z_0)\le g$
as soon as
\[
  B(\eta,\xi-1,\theta)\ge\xi h(\theta/\xi)+\theta L+\min\{-L,\tau(\xi)\}+e_1 .
\]

\emph{Step 6: the certificate.}  We cover the box by rectangles
$[\theta_a,\theta_b]\times[\xi_a,\xi_b]$ with rational corners.  On each one
we use a lower bound $g'$ for $g$.  In (a), $g'=r_0-1$.  In (b), $M_0\le nm$
with $m=\min\{1/(r_0-1),(1-\theta_a+1/N_0)\bar q_c/(1-\bar q_c)\}$, where
$\bar q_c=\xi_a/(r_0(\xi_a-\theta_b))$; the second term is used only when
$\bar q_c<1$.  Indeed $(x+1)/(x-p+2)\le x/(x-p)=\xi/(\xi-\theta)$, and
$(n-p+1)/n\le1-\theta+d$.  Since $(x+1)/(x+1+M_0)\ge\xi/(\xi+m)$, we may take
$g'=r_0\xi_a/(\xi_a+m)$.  A rectangle is settled by one of three tests.
\begin{itemize}
\item[(E1)] The least value of $-h(\theta)-\theta L$ on
$[\theta_a,\theta_b]$ exceeds $\tau(\xi_a)+e_1$.  Then no point of the
rectangle satisfies the inequality of Step~4, since
$\tau(\xi)\le\tau(\xi_a)$.
\item[(E2)] $\xi_b-1+1/N_0\le(g'-1)\theta_a$.  Then Step~3 applies.
\item[(E3)] $(g'-1)\theta_a>1/N_0$, and with
$\bar\eta=\max\{\theta_a+\frac1{N_0},\,(\theta_a+\frac1{N_0})(\xi_b-1)/((g'-1)\theta_a-\frac1{N_0})+\frac1{N_0}\}$,
\[
  B(\bar\eta,\xi_a-1,\theta_a)\ge\xi_bH+\theta_bL+\min\{-L,\tau(\xi_a)\}+e_1,
\]
where $H$ is the largest value of $h$ on $[\theta_a/\xi_b,\theta_b/\xi_b]$.
The bound for $\eta$ in Step~5 increases in $\xi$ and in $d$, and decreases in
$\theta$, so $\eta\le\bar\eta$.  By the monotonicity of $B$, the left side is
at most $B(\eta,\xi-1,\theta)$.  The function $\xi h(\theta/\xi)$ increases in
$\xi$, so the right side is at least the right side of Step~5.
\end{itemize}
We started from a $48\times48$ grid and quartered every rectangle that no
test settled.  All quantities were evaluated in interval arithmetic with
outward rounding at $100$ bits.  Every rectangle was settled.
\[
\begin{array}{cccccc}
  r_0&\text{(E1)}&\text{(E2)}&\text{(E3)}&\text{rounds}&\text{least margin in (E3)}\\\hline
  \frac{25}7&421&1091&132945&7&1.6\cdot10^{-8}\\
  \frac{10}3&205&1503&596&0&0.035\\
  \frac{13}4&198&1376&1438&3&2.6\cdot10^{-5}
\end{array}
\]
Here ``rounds'' is the largest number of quarterings of one grid rectangle.
So $\Pi(Z_0)\le g$ for every $n\ge N_0$, and Step~1 concludes.  In (a),
$g=\frac{18}7$.
\end{proof}

By Corollary~\ref{cor:prefixmono}(a), part (a) proves the conjecture stated
after Corollary~\ref{cor:belowtwo} for every $r\ge\frac{25}7$.  At
$r_0=\frac{25}7$ the certificate did not close for $N_0=340$ within eight
rounds.  At $r_0=\frac{18}5$ it closes for $N_0=301$, so there
Corollary~\ref{cor:rootfourfinite} suffices for small $n$.  The limit of the method is near $3.55$ for
all rows and near $3.2$ for the second row (the remarks after
Corollary~\ref{cor:rootfourfinite}).  At $r_0=\frac{16}5$ the certificate for
the second row did not close even for $N_0=2000$ within six rounds.  Closer to these limits the box
certificate needs a larger $N_0$, and the exact check must then reach $N_0-1$.
Section~\ref{sec:allrowstwo} proves all rows at every $r\ge2$ by another
argument.


\begin{thebibliography}{9}

\bibitem{coefficient-mass}
B.~Pham,
\newblock Displaced-zero certificates and the coefficient mass of polynomial
multiples,
\newblock manuscript in preparation.

\end{thebibliography}

\end{document}
